% NT-Disk: Astrophysics of Black Holes
% Original: I. D. Novikov and K. S. Thorne (1973)
% LaTeX Translation in Physical Review D format
%
% From: "Black Holes" (Les Houches 1972), eds. C. DeWitt and B. S. DeWitt
% Gordon and Breach, New York, 1973, pp. 345--450
%
\documentclass[aps,prd,twocolumn,superscriptaddress,nofootinbib,%
  longbibliography]{revtex4-2}

% === Packages ===
\usepackage{amsmath,amssymb,amsfonts}
\usepackage{mathrsfs}       % Script letters (\mathscr)
\usepackage{bm}             % Bold math (\bm)
\usepackage{graphicx}       % Figures
\usepackage{hyperref}       % Hyperlinks
\usepackage{xcolor}         % Colors (for placeholder figures)
\usepackage{array}          % Enhanced tables
\usepackage{booktabs}       % Professional tables

% === Custom Commands ===
% Fundamental constants
\newcommand{\me}{m_e}          % electron mass
% \mp is already defined in LaTeX (minus-plus); use m_p directly in equations
\newcommand{\Ry}{R_y}          % Rydberg energy
\newcommand{\kB}{k}            % Boltzmann constant (paper uses just k)

% Vectors and tensors
\newcommand{\fvec}[1]{\mathbf{#1}}     % 4-vectors (bold upright)
\newcommand{\tvec}[1]{\mathbf{#1}}     % 3-vectors (bold)

% Kerr metric script functions (Section 5.4)
\newcommand{\scrA}{\mathscr{A}}
\newcommand{\scrB}{\mathscr{B}}
\newcommand{\scrC}{\mathscr{C}}
\newcommand{\scrD}{\mathscr{D}}
\newcommand{\scrE}{\mathscr{E}}
\newcommand{\scrF}{\mathscr{F}}
\newcommand{\scrG}{\mathscr{G}}
\newcommand{\scrP}{\mathscr{P}}        % Power/luminosity script P

% Accretion notation
\newcommand{\Mdot}{\dot{M}}
\newcommand{\Mdotzero}{\dot{M}_0}
\newcommand{\Msun}{M_\odot}
\newcommand{\Lcrit}{L_{\text{crit}}}
\newcommand{\rg}{r_g}                  % gravitational radius
\newcommand{\rs}{r_s}                  % sonic radius
\newcommand{\rms}{r_{\text{ms}}}       % marginally stable orbit
\newcommand{\Lms}{\tilde{L}_{\text{ms}}}

% Differential operators
\newcommand{\grad}{\boldsymbol{\nabla}}
\newcommand{\divg}{\boldsymbol{\nabla} \cdot}

% Averages
\newcommand{\avg}[1]{\langle #1 \rangle}

% Dimensionless parameters for Section 5.9
\newcommand{\mstar}{M_*}
\newcommand{\mdotstar}{\dot{M}_*}
\newcommand{\rstar}{r_*}

% Text shortcuts
\newcommand{\ie}{i.e.}
\newcommand{\eg}{e.g.}
\newcommand{\cf}{cf.}

% === Equation numbering ===
% The paper uses (section.subsection.number) format
% We use \tag{} for explicit equation numbers to match original
\numberwithin{equation}{section}

% === Document ===
\begin{document}

\title{Astrophysics of Black Holes}

\author{Igor D. Novikov}
\thanks{Supported in part by the Academy of Sciences of the USSR}
\affiliation{Institute of Applied Mathematics, Academy of Sciences, Moscow}

\author{Kip S. Thorne}
\thanks{Supported in part by the U.S. National Science Foundation [GP-27304, GP-28027]}
\affiliation{California Institute of Technology, Pasadena}

\date{1973}

\begin{abstract}
Lectures on the astrophysics of black holes, covering radiation processes in plasmas,
hydrodynamics, accretion onto isolated and binary black holes, relativistic disk structure,
and cosmological aspects of black holes. Originally published in \textit{Black Holes}
(Les Houches 1972), eds.\ C.~DeWitt and B.~S.~DeWitt (Gordon and Breach, New York, 1973),
pp.~345--450.
\end{abstract}

\maketitle
\tableofcontents

% === SECTION 1: Introductory Remarks ===
% Section 1: Introductory Remarks
% Original pages: 347 (book page numbering)
% Translated from NT-Disk.pdf

\section{Introductory Remarks}
\label{sec:1}%
To analyze astrophysical aspects of black holes one must bring input from two
directions: from the elegant realm of general relativity (lectures of Hawking,
Carter, Bardeen, and Ruffini), and from the more mundane realm of astro-
physics and plasma physics (\S\,2 of these lectures). We shall here combine these
inputs to discuss the origin of stellar black holes (\S\,3); and to analyze observable
aspects of black holes in the interstellar medium (\S\,4), in binary star systems and
the nuclei of galaxies (\S\,5), and in cosmological contexts (\S\,6).

Readers with strong astrophysical backgrounds will wish to skip over our
brief treatment of the fundamentals of astrophysics and plasma physics (\S\,2),
and dig immediately into our discussion of black-hole astrophysics (\S\S\,3--8).
However, for the reader whose previous training has focussed primarily on
gravitation theory, the material of \S\,2 is an important prerequisite for under-
standing the rest of the notes.

The authors are deeply indebted to Mr.\ Alan Lightman for removing a large
number of errors from the manuscript.


% === SECTION 2: Radiation and Plasma Physics ===
% Section 2: Thermal Bremsstrahlung (``Free-Free Radiation'' From a Plasma)
% Merged from agent_01.tex through agent_11.tex
% Subsections 2.1--2.9

\section{Thermal Bremsstrahlung (``Free-Free Radiation'' From a Plasma)}
\label{sec:2}

\subsection{Thermal Bremsstrahlung (``Free-Free Radiation'' From a Plasma)}
\label{sec:2.1}

\textit{Basic references}\quad chapter~6 of Shkarofsky, Johnston, and
Bachynski~(1966)\cite{Shkarofsky1966};
chapter~15 of Jackson~(1962)\cite{Jackson1962}; \S4 of
Ginzburg~(1967)\cite{Ginzburg1967}.

\bigskip

\textit{Notation for fundamental constants}

\medskip

\noindent
\begin{tabular}{ll}
$h$            & Planck's constant \\
$m_e$          & rest mass of electron \\
$m_p$          & rest mass of proton \\
$e$            & charge of electron \\
$c$            & speed of light \\
$\alpha$       & fine-structure constant \\
$r_0$          & classical electron radius \\
$R_y$          & Rydberg energy \\
$Z$            & charge of ions in plasma \\
$\mathcal{A}$  & atomic weight of ions \\
$k$            & Boltzmann's constant \\
$\xi$          & Euler constant ($= 1.78\ldots$)
\end{tabular}

\bigskip

\textit{Radiation from a single classical collision}

\medskip

Consider a single ion of charge~$Z$, at rest in the laboratory frame. Let a single
electron of speed $v \ll c$ (nonrelativistic) scatter off the nucleus (Coulomb
scattering), with impact parameter~$b$. Let $I(\omega)$ be the total energy per unit
circular frequency, $\omega$, radiated by the electron as it scatters. A simple classical
calculation [e.g.\ Jackson~(1962)\cite{Jackson1962}, p.~507] gives
%
\begin{equation}
I(\omega) = \frac{2}{3} \frac{e^2}{m_e c^3} \left|\Delta \mathbf{v}_\omega\right|^2,
\tag{2.1.1}
\end{equation}
%
where $\Delta \mathbf{v}_\omega$ is the change in electron velocity during a time $\tau = 1/\omega$ centered
about the electron's point of closest approach to the nucleus.

It is useful to examine two limiting cases: small-angle scattering ($\theta \ll 1$)
corresponding to large impact parameter ($b \gg b_{s\text{-}l}$); and large-angle scattering
($\theta \sim 2\pi$), corresponding to small impact parameter ($b \ll b_{s\text{-}l}$). The impact
parameter, $b_{s\text{-}l}$, which separates small-angle from large-angle scattering, is given
by
%
\begin{equation}
Ze^2 / b_{s\text{-}l} = \tfrac{1}{2}\, m_e v^2.
\tag{2.1.2}
\end{equation}

\textit{For small-angle scattering} the total scattering angle $\theta$ is
%
\begin{equation}
\theta = \frac{1}{v}\, |\Delta \mathbf{v}|
      = \frac{Ze^2/b}{\frac{1}{2}\, m_e v^2}
      = \frac{b_{s\text{-}l}}{b}\,,
\tag{2.1.3}
\end{equation}
%
and the time during which the scattering occurs is $\Delta t \approx b/v$.  Hence, for
$\tau = 1/\omega \gg \Delta t$, the change in velocity during time $\tau$, $|\Delta \mathbf{v}_\tau|$,
is the full change $|\Delta \mathbf{v}|$ of~(2.1.3); while for $\tau = 1/\omega \ll \Delta t$ the change is
essentially zero.  Inserting these changes into equation~(2.1.1), we obtain
%
\begin{align}
I(\omega) &= 0  &\text{for}&\quad \omega \gg v/b, \notag \\[4pt]
I(\omega) &= \frac{8}{3}\,\frac{Z^2 e^2}{m_e c}\left(\frac{r_0}{v\,b}\right)^{\!2}
           &\text{for}&\quad \omega \ll v/b.
\tag{2.1.4}
\end{align}

\textit{For large-angle scattering} the electron orbit is very nearly a parabola, with
``perihelion'' at radius
%
\begin{equation}
r_p = b^2 / b_{s\text{-}l}
\tag{2.1.5a}
\end{equation}
%
and with speed at perihelion
%
\begin{equation}
v_p = c\bigl(2 Z r_0 / r_p\bigr)^{1/2}.
\tag{2.1.5b}
\end{equation}
%
During time intervals $\tau = 1/\omega \ll r_p/v_p$, a negligible amount of velocity change
occurs; hence
%
\begin{equation}
I(\omega) = 0 \qquad \text{for}\quad \omega \gg v_p/r_p
           = c(2Zr_0\, b_{s\text{-}l}^{\,3})^{1/2}/b^3.
\tag{2.1.6a}
\end{equation}

During time intervals $\tau = 1/\omega \gg r_p/v_p$ but $\tau < b/v$, the electron is initially
($\tau = -\tau/2$) headed directly toward the nucleus with ``parabolic'' speed

\noindent it swings into a sharp turn around the nucleus; and it emerges
at $t = +\tau/2$ with its initial velocity precisely reversed.
Consequently, $|\Delta v_\omega|$ is equal to~$2v_f$, and $I(\omega)$~is
\begin{equation}
  I(\omega)
  = \frac{32}{3}\,\frac{e^{2}}{c^{3}}
    \left(\frac{Zr_0\,\omega\,c^{2}}{3}\right)^{\!2/3}
  \quad\text{for}\quad
  \frac{v}{b_{s-1}} < \omega < \frac{v_p}{r_p}
  = \frac{c(2Zr_0\,b_{s-1}^{3})^{1/2}}{b^{3}}\,.
  \tag{2.1.6b}
\end{equation}

%=============================================================================
% Radiation from a nonrelativistic beam of electrons
%=============================================================================

\bigskip
\noindent\textit{Radiation from a nonrelativistic beam of electrons}

\medskip

Consider a single ion of charge~$Z$, at rest in the laboratory frame.
Bombard the ion with a monochromatic beam of electrons with speed
$v \ll 1$ and flux (number per unit time per unit area)~$S$.  Examine the
radiation of frequency~$\nu$ (circular frequency $\omega = 2\pi\nu$)
emitted by the electrons as they scatter off the nucleus.
Let $\mathscr{P}_\nu$ be the total power emitted per unit frequency
(ergs~sec$^{-1}$~Hz$^{-1}$); and define the emission cross section per
unit frequency, $d\sigma/d\nu$, by
\begin{equation}
  \mathscr{P}_\nu = (d\sigma/d\nu)\,S\,h\nu.
  \tag{2.1.7}
\end{equation}
The dependence of the emission cross section on photon frequency~$\nu$
and on electron speed~$v$ is most conveniently expressed in terms of
the two dimensionless parameters
\begin{equation}
  \frac{h\nu}{\tfrac{1}{2}m_e v^{2}}\,,
  \qquad
  \frac{\tfrac{1}{2}m_e v^{2}}{2R_y}
  = \left(\frac{v/c}{\alpha Z}\right)^{\!2}.
  \tag{2.1.8}
\end{equation}

On a sheet of paper (Figure~\ref{fig:2.1.1}) plot
$(h\nu)/(\tfrac{1}{2}m_e v^{2})$ vertically, and plot
$(\tfrac{1}{2}m_e v^{2})/(2R_y)$ horizontally.  This 2-dimensional plot
can be split into several different regions, in each of which the details
of the emission process are different.

%--- Forbidden region ---------------------------------------------------
\bigskip
\noindent\textit{Forbidden region}\quad
The individual photons emitted while scattering cannot have energies
greater than the electron kinetic energy---unless the electron gets
captured into a bound state around the ion.  But capture produces
``free-bound'' rather than bremsstrahlung radiation (\S\,2.2).
Consequently, the region
\begin{equation}
  (h\nu)/(\tfrac{1}{2}m_e v^{2}) > 1
  \tag{2.1.9a}
\end{equation}
of Figure~\ref{fig:2.1.1} is a ``forbidden region''; no bremsstrahlung
is emitted in this region; $d\sigma/d\nu$ vanishes in this region.

%--- Photon-discreteness region -----------------------------------------
\bigskip
\noindent\textit{Photon-discreteness region}\quad In the region
\begin{equation}
  \tfrac{1}{3} < (h\nu)/(\tfrac{1}{2}m_e v^{2}) \le 1
  \tag{2.1.9b}
\end{equation}
no more than 3~photons can be emitted by each scattering.  This ``photon
discreteness'' has a significant effect on the emission cross section,
cutting it down toward zero as one approaches the edge of the forbidden
region, $(h\nu)/(\tfrac{1}{2}m_e v^{2}) = 1$.  No classical calculation
can reveal such an effect; the photon-discreteness region must be
analyzed in a quantum mechanical manner; see, e.g., Heitler
\cite{Heitler1954}, and see below.

%=============================================================================
% Large-angle region (page 350)
%=============================================================================

\bigskip
\noindent\textit{Large-angle region}\quad
All small-angle scatterings last too long
($\Delta\tau \sim b/v > b_{s-1}/v$; cf.\ eq.~(2.1.4)) to produce any
radiation of $\tau = 1/\omega < b_{s-1}/v$; cf.\ eq.~(2.1.4).
Consequently, in the region
\begin{equation}
  \frac{h\nu}{\tfrac{1}{2}m_e v^{2}}
  > \frac{h\omega}{\tfrac{1}{2}m_e v^{2}}
  = \left(\frac{\tfrac{1}{2}m_e v^{2}}{Z^{2}R_y}\right)^{\!1/2}
  \tag{2.1.10}
\end{equation}
only large-angle scatterings produce radiation.  This region is shown in
Figure~\ref{fig:2.1.1}.  The power per unit frequency emitted at a fixed
frequency~$\nu$ in this region is
\begin{equation}
  \mathscr{P}_\nu
  = 2\pi \int_0^{b_{\max}} I(\omega)\,S\,2\pi\,b\,db\,,
  \tag{2.1.11a}
\end{equation}
where $I(\omega)$ is given by the large-angle formula~(2.1.6), and
$b_{\max}$ is the ``cutoff''
\begin{equation}
  b_{\max} = (2c^{2}\,Zr_0\,b_{s-1}^{3})^{1/6}
  \tag{2.1.11b}
\end{equation}
in that formula.

Performing the integration, dividing by $S\,h\nu$ to obtain
$d\sigma/d\nu$, and reexpressing the result in terms of fundamental
atomic constants, one obtains
\begin{equation}
  \frac{d\sigma}{d\nu}
  = \frac{2^{1/3}\,16\,\alpha\,c^{2}\,Z^{2}\,r_0^{2}}
         {3^{5/3}\;v^{2}}\,.
\end{equation}

A more exact classical calculation gives a slightly different numerical
coefficient:
\begin{equation}
  \left(\frac{d\sigma}{d\nu}\right)_{\!\text{LA}}
  = \frac{16\pi}{3\sqrt{3}}\;
    \frac{\alpha\,c^{2}\,Z^{2}\,r_0^{2}}{v^{2}}\,.
  \tag{2.1.12}
\end{equation}

\noindent(Here ``LA'' stands for ``large-angle region''.)

In other regions of Figure~\ref{fig:2.1.1}, one gets other formulas for
$d\sigma/d\nu$ (see below).  It is conventional in all regions to lump
the deviations from this large-angle cross section into a correction
$G(\nu,v)$ which is called the ``Gaunt factor'':
\begin{equation}
  G(\nu,v)
  \equiv \frac{(d\sigma/d\nu)}{(d\sigma/d\nu)_{\text{LA}}}\,.
  \tag{2.1.13}
\end{equation}

Thus, in the large-angle region $G = 1$; and in the forbidden region
$G = 0$.  Throughout the large-angle and photon-discreteness portion of
the large-angle region, $G$~remains approximately~1; see
Figure~\ref{fig:2.1.1}.

%=============================================================================
% Small-angle, classical region (pages 350--351)
%=============================================================================

\bigskip
\noindent\textit{Small-angle, classical region}\quad
In the region $\tau = 1/\omega \gg b_{s-1}/v$---i.e.
\begin{equation}
  \frac{h\nu}{\tfrac{1}{2}m_e v^{2}}
  \ll \left(\frac{\tfrac{1}{2}m_e v^{2}}{2R_y}\right)^{\!1/2},
  \tag{2.1.14}
\end{equation}
radiation is produced by small-angle scatterings as well as by
large-angle scatterings.  As a function of impact parameter~$b$, the
energy radiated by a single electron is constant in the large-angle
region, $b < b_{s-1}$ [eq.~(2.1.6b)], and decreases as~$1/b^{2}$ in
the small-angle region, $b > b_{s-1}$ [eq.~(2.1.4)].  Hence, when one
integrates over $2\pi\,b\,db$ to get the power radiated,
$\mathscr{P}_\nu$, one obtains

\smallskip
\noindent(contribution from small angles)
$\displaystyle{}= \ln\!\left(\frac{b_{\max}}{b_{s-1}}\right)$,

\smallskip
\noindent(contribution from large angles) ${}= \text{const}$,

\smallskip
\noindent where $b_{\max}$ is the cutoff in the small-angle region
(eq.~2.1.4):
\begin{equation}
  b_{\max} = v/\omega.
  \tag{2.1.15}
\end{equation}

Since the small-angle contribution is logarithmically dominant, one can
ignore the contribution from large-angle scatterings and calculate
\begin{equation}
  \frac{d\sigma}{d\nu}
  = \frac{\mathscr{P}_\nu}{S\,h\nu}
  = \frac{16\,\alpha\,c^{2}\,Z^{2}\,r_0^{2}}{3\,v^{2}}\,
    \ln\!\left(\frac{b_{\max}}{b_{s-1}}\right).
  \tag{2.1.16}
\end{equation}

\noindent A more exact classical calculation gives a slightly different
argument in the logarithm:
\begin{equation}
  \frac{d\sigma}{d\nu}
  = \frac{16\,\alpha\,c^{2}\,Z^{2}\,r_0^{2}}{3\,v^{2}}\,
    \ln\!\left(\frac{2\,b_{\max}}{\xi\;b_{s-1}}\right),
  \tag{2.1.18}
\end{equation}
where $\xi = 1.781\ldots$ is the ``Euler constant''.  Consequently, the
Gaunt factor for the small-angle, classical region is
\begin{equation}
  G(\nu,v)
  = \frac{\sqrt{3}}{\pi}\,
    \ln\!\left[
      \frac{2}{\xi}
      \left(\frac{\tfrac{1}{2}m_e v^{2}}{h\nu}\right)
      \left(\frac{\tfrac{1}{2}m_e v^{2}}{2R_y}\right)^{\!1/2}
    \right].
  \tag{2.1.19}
\end{equation}

Actually, this small-angle classical result is not valid throughout the
region~(2.1.14).  The uncertainty principle requires a modification at
electron energies $\tfrac{1}{2}m_e v^{2}$ larger than~$Z^{2}R_y$.

%=============================================================================
% Small-angle, uncertainty-principle region (pages 351--352)
%=============================================================================

\bigskip
\noindent\textit{Small-angle, uncertainty-principle region}\quad
Consider an electron with impact parameter~$b$ and speed~$v$.  The
electron is actually not a classical object; rather, it is a
quantum-mechanical wave packet.  To scatter with impact parameter~$b$,
the wave packet (i)~must have transverse dimensions, $\Delta x$, less
than~$b$
\begin{equation}
  \Delta x < b,
\end{equation}
and (ii)~must not spread transversely by more than~$b$ during the
classical scattering time $\Delta t = b/v$
\begin{equation}
  b > \left(\frac{\text{spreading}}{\text{during }\Delta t}\right)
  = \left(\frac{\Delta p_x}{m_e}\right)\Delta t
  > \left(\frac{\hbar}{m_e\,\Delta x}\right)\Delta t
  > \frac{\hbar}{m_e\,v}
  = \frac{b}{m_e\,v}\,\frac{\hbar}{b}\,.
\end{equation}

\noindent Thus, the uncertainty principle (``$\ge$'' in the above
equation) prevents the existence of scatterings with~$b$ less than the
electron de~Broglie wavelength
\begin{equation}
  b_{dB} = \hbar / m_e v.
  \tag{2.1.20}
\end{equation}

If $b_{dB} < b_{s-1}$, this limitation has no effect on small-angle
scatterings; and the small-angle, classical Gaunt factor~(2.1.19) is
valid.  If $b_{dB} > b_{s-1}$, the uncertainty principle comes into
play, and one must use~$b_{dB}$ rather than~$b_{s-1}$ as the lower
limit on the small-angle integral~(2.1.17) for $d\sigma/d\nu$.  The
result is
\begin{equation}
  \frac{d\sigma}{d\nu}
  = \frac{16\,\alpha\,c^{2}\,Z^{2}\,r_0^{2}}{3\,v^{2}}\,
    \ln\!\left(\frac{2\,b_{\max}}{\xi\;b_{dB}}\right),
  \tag{2.1.21}
\end{equation}
corresponding to the Gaunt factor
\begin{equation}
  G(\nu,v)
  = \frac{\sqrt{3}}{\pi}\,
    \ln\!\left(\frac{2\,b_{\max}}{\xi\;b_{dB}}\right)
  = \frac{\sqrt{3}}{\pi}\,
    \ln\!\left[
      \frac{2}{\xi}
      \left(\frac{\tfrac{1}{2}m_e v^{2}}{h\nu}\right)
    \right].
  \tag{2.1.22}
\end{equation}

The ``small-angle, uncertainty-principle region,'' in which this
expression holds, is separated from the ``small-angle, classical
region'' by the line
\begin{equation}
  b_{dB}/b_{s-1}
  = (\tfrac{1}{2}m_e v^{2})/(Z^{2}R_y) = 1.
\end{equation}

\noindent See Figure~\ref{fig:2.1.1}.

From Figure~\ref{fig:2.1.1} it is clear that the uncertainty principle
cannot have any effect on the large-angle region.

%=============================================================================
% Figure 2.1.1
%=============================================================================

\begin{figure}[htbp]
\centering
% \includegraphics[width=\columnwidth]{fig_2_1_1}
\fbox{\parbox{0.9\columnwidth}{\centering
[Figure 2.1.1 placeholder]\\[6pt]
Horizontal axis:
$\bigl(\tfrac{1}{2}m_e v^{2}\bigr)\big/\bigl(Z^{2}R_y\bigr)$.\quad
Vertical axis:
$h\nu\big/\bigl(\tfrac{1}{2}m_e v^{2}\bigr)$.\\[4pt]
Five regions shown:\\[2pt]
(1) Forbidden region ($G = 0$, $d\sigma/d\nu = 0$)
    --- above the line $h\nu/(\tfrac{1}{2}m_ev^2) = 1$.\\
(2) Photon-discreteness region ($G \lesssim 1$)
    --- between $\tfrac{1}{3}$ and~$1$.\\
(3) Large-angle region ($G = 1$)
    --- below discreteness, above
    $h\nu/(\tfrac{1}{2}m_ev^2) =
     [(\tfrac{1}{2}m_ev^2)/(2R_y)]^{1/2}$.\\
(4) Small-angle, classical region:
    $G = \dfrac{\sqrt{3}}{\pi}
    \ln\!\Bigl[\dfrac{4}{\xi}
    \bigl(\dfrac{\frac{1}{2}m_ev^{2}}{h\nu}\bigr)
    \bigl(\dfrac{\frac{1}{2}m_ev^{2}}{2R_y}\bigr)^{\!1/2}\Bigr]$
    --- lower-left.\\[4pt]
(5) Small-angle, U.P.\ (uncertainty-principle) region:
    $G = \dfrac{\sqrt{3}}{\pi}
    \ln\!\Bigl[\dfrac{4}{\xi}
    \bigl(\dfrac{\frac{1}{2}m_ev^{2}}{h\nu}\bigr)\Bigr]$
    --- lower-right.
}}
\caption{Regions and Gaunt factors for bremsstrahlung of frequency~$\nu$
emitted by an electron of kinetic energy~$\tfrac{1}{2}m_e v^{2}$
when it impinges on an ion of charge~$Ze$.}
\label{fig:2.1.1}
\end{figure}

out that throughout the photon-discreteness portion of the large-angle region,
$G$ remains approximately 1; see Figure~2.1.1.

\textit{Small-angle, classical region}\quad In the region $\tau = 1/\omega \gg b_{s-l}/v$---i.e.
\begin{equation}
\frac{h\nu}{\frac{1}{2}m_e v^2}
\ll \left(\frac{\frac{1}{2}m_e v^2}{Z^2 R_y}\right)^{1/2}
\tag{2.1.14}
\end{equation}

In the ``photon-discreteness'' portion of the small-angle, uncertainty-principle
region, discreteness effects modify the Gaunt factor (2.1.22) into the form
\begin{subequations}
\begin{align}
G(\nu,v) &= \frac{\sqrt{3}}{\pi}\,
\ln\!\left[\frac{\frac{1}{2}m_e(v_i + v_f)^2}{h\nu}\right],
\tag{2.1.24a}\\[6pt]
G(\nu,v) &= \frac{\sqrt{3}}{\pi}\,
\ln\!\left[\frac{v_i + v_f}{v_i - v_f}\right]^2,
\tag{2.1.24b}
\end{align}
\end{subequations}

Here $v_i$ and $v_f$ are the electron speeds before and after the collision:
\begin{equation}
\tfrac{1}{2}m_e v_f^2 = \tfrac{1}{2}m_e v_i^2 - h\nu.
\end{equation}

[Expression (2.1.24) is the result of a quantum-mechanical Born-approximation
calculation, valid throughout the small-angle, uncertainty-principle region.]

%%---------------------------------------------------------------------------
%% Radiation from a plasma
%%---------------------------------------------------------------------------

\bigskip
\textit{Radiation from a plasma}

Consider a plasma which may contain a variety of ionic species.  Focus attention
on radiation from all ions of a particular type, with charge $Ze$ and atomic weight
$A$.  Let $C_A$ be the concentration by mass of the species of atom that gives rise to
this ion, and let $f_i$ be the fraction of all such atoms ionized to the state of interest.
Let $f_e$ be the number of unbound electrons per baryon in the gas.  Then
\begin{equation}
n_I = \text{(number of ions per unit volume)} = f_i C_A \rho_0 / A m_p.
\tag{2.1.25}
\end{equation}

Here $\rho_0$ is the density of rest mass.  (Of course, $f_i \leq 1$; $C_A \leq 1$.)  The ions
and electrons both have Maxwell velocity distributions; but because $m_e \ll A m_p$,
the velocities of the electrons and their accelerations during Coulomb scattering
are far greater than those of the ions.  Therefore, in calculating bremsstrahlung
one can regard the ions as at rest.  The number of photons per unit mass of plasma
is $f_e/m_p$, so the total power per unit frequency emitted from one gram of plasma
by electron-ion scatterings is
\begin{equation}
\varepsilon_\nu
= \frac{C_A f_i}{A m_p}
\underbrace{\vphantom{\bigg|}\quad}_{\substack{
\text{fractions of all} \\
\text{ions per} \\
\text{unit mass}
}}
\cdot \int
\underbrace{\vphantom{\bigg|}\quad}_{\substack{
\text{number} \\
\text{of electrons}
}}
\underbrace{\vphantom{\bigg|}\quad}_{\substack{
\text{electrons that have} \\
\text{speeds $v$ in range $dv$}
}}
\frac{f_e \rho_0}{m_p}\;v\,h\nu\;
\frac{d\sigma}{dv}\;dv.
\tag{2.1.26}
\end{equation}

Here $T$ is the temperature in the plasma.  By inserting into the integrand
expressions (2.1.12) and (2.1.13) for $d\sigma/dv$ and performing the integration, one
obtains
\begin{align}
\varepsilon_\nu
&= \frac{C_A f_i f_e^2}{A m_p}\,
\frac{32\pi e^6}{3\,m_e c^3}\,
\left(\frac{2\pi m_e}{3kT}\right)^{\!1/2}\!
\frac{\rho_0^2}{m_p^2}\;
e^{-h\nu/kT}\;
\bar{G}(\nu,T)
\notag\\[4pt]
&= 2.5 \times 10^{10}\;
\frac{\text{erg}}{\text{g\;sec\;Hz}}\;
\left(\frac{C_A f_i f_e^2}{A}\right)
\left(\frac{\rho_0}{\text{g\;cm}^{-3}}\right)
\left(\frac{\rho_0}{A}\right)
T_K^{\,1/2}\;e^{-h\nu/kT}\;\bar{G}(\nu,T).
\tag{2.1.27}
\end{align}

%%---------------------------------------------------------------------------
%% Mean Gaunt factor \bar{G}
%%---------------------------------------------------------------------------

Here $T_K$ is temperature measured in ${}^\circ$K, and $\bar{G}(\nu,T)$ is the ``Maxwell--Boltzmann-averaged
Gaunt factor'', also called the ``mean Gaunt factor'':
\begin{equation}
\bar{G}(\nu,T) = \int_0^\infty
G\!\left(\nu,\;v = \left[\frac{2(xkT + h\nu)}{m_e}\right]^{1/2}\right)
e^{-x}\,dx.
\tag{2.1.28}
\end{equation}

Corresponding to the 2-dimensional plot of Gaunt factors (Figure~2.1.1) one
can make a 2-dimensional plot of mean Gaunt factors (Figure~2.1.2).  It is
straightforward to calculate the mean Gaunt factors shown there (up to
quantities of order unity) from the Gaunt factors of Figure~2.1.1.  Notice that,

%%---------------------------------------------------------------------------
%% Figure 2.1.2 — Mean Gaunt factors
%%---------------------------------------------------------------------------

\begin{figure}[htbp]
\centering
% \includegraphics[width=\columnwidth]{fig_2_1_2}
\fbox{\parbox{0.9\columnwidth}{\centering [Figure 2.1.2 placeholder]\\[6pt]
Mean Gaunt factors for bremsstrahlung of frequency $\nu$ emitted by
electron-ion collisions in a plasma of temperature $T$.\\[4pt]
Horizontal axis: $h\nu/kT$.\\
Vertical axis: $(kT)/(Z^2 R_y)$.\\[4pt]
Regions shown:\\
--- ``Small-angle, u.p., tail region'',
$\bar{G} = \bigl(\frac{3}{\pi}\bigr)\,\bigl(\frac{kT}{h\nu}\bigr)^{1/2}$\\
--- ``Small-angle, u.p.\ region'',
$\bar{G} = \frac{\sqrt{3}}{\pi}\ln\!\left[\frac{4}{\xi}\,\frac{kT}{h\nu}\right]$\\
--- ``Small-angle, classical region'',
$\bar{G} = \frac{\sqrt{3}}{\pi}\,\ln\!\left[\left(\frac{4}{\xi^{3/2}}\right)
\left(\frac{kT}{h\nu}\right)\left(\frac{kT}{Z^2 R_y}\right)^{1/2}\right]$\\
--- ``Large-angle region'', $\bar{G} = 1$\\
--- ``Large-angle, tail region'', $\bar{G} = 1$\\
--- ``Forbidden region''}}
\caption{\textit{Figure~2.1.2.}\ Mean Gaunt factors for bremsstrahlung of frequency $\nu$ emitted by
electron-ion collisions in a plasma of temperature~$T$.}
\label{fig:2.1.2}
\end{figure}

aside from quantities of order unity, the mean Gaunt factors in the large-angle
and small-angle regions are obtained by simply taking the Gaunt
factors and making the replacement $\frac{1}{2}m_e v^2 \approx kT$:
\begin{equation}
\bar{G}(\nu,kT) \simeq G\bigl(\nu,\;\tfrac{1}{2}m_e v^2 = kT\bigr).
\end{equation}

Notice also that the ``forbidden region'' of the Gaunt-factor plot is replaced, in
the mean-Gaunt-factor plot, by ``tail regions''.  The radiation in these tail regions
is produced by the high-energy tail, $\frac{1}{2}m_e v^2 \approx h\nu \gg kT$, of the Maxwell--Boltzmann
distribution.

Tables of mean Gaunt factors, calculated with much higher accuracy than
Figure~2.1.2, are given by \citet{Green1959} and by \citet{Karzas1961}.

From Figure~2.1.1 it is clear that the uncertainty principle cannot have any
effect on the large-angle region.

%%---------------------------------------------------------------------------
%% Total emissivity and frequency-averaged mean Gaunt factor
%%---------------------------------------------------------------------------

The total energy emitted by electron-ion collisions is obtained by integrating
expression (2.1.27) over frequency:
\begin{align}
e &= \int \varepsilon_\nu\,d\nu
= \frac{16\,C_A f_i f_e^2}{3}\;
\frac{\alpha\,m_e c^2_0}{A}\;
\left(\frac{2\pi kT}{m_e c^2}\right)^{\!1/2}\!
\frac{\rho_0}{m_p^2}\;
\left(\frac{C_A f_i f_e^2}{A}\right)
\left(\frac{\rho_0}{\text{g\;cm}^{-3}}\right)
T_K^{\,1/2}\;\bar{G}(T)
\notag\\[4pt]
&= 5.2 \times 10^{20}\;
\frac{\text{ergs}}{\text{g\;sec}}\;
\left(\frac{C_A f_i f_e^2}{A}\right)
\left(\frac{\rho_0}{\text{g\;cm}^{-3}}\right)
T_K^{\,1/2}\;\bar{G}(T).
\tag{2.1.29}
\end{align}

Here $\bar{G}(T)$ is the frequency-averaged mean Gaunt factor
\begin{equation}
\bar{G}(T) = \int_0^\infty \bar{G}(\nu = xkT/h,\;T)\;e^{-x}\,dx.
\tag{2.1.30}
\end{equation}

From a value of $\bar{G} \simeq 1.2$ at $kT/Z^2 R_y \sim 0.01$, it increases gradually to a maximum
of $1.42$ at $kT/Z^2 R_y \sim 1$, and then decreases gradually toward an asymptotic,
high-temperature value of $1.103$ at $kT/Z^2 R_y > 100$; see Figure~2 of \citet{Green1959}.

In a plasma at nonrelativistic temperatures, the radiation from electron-electron
collisions and from ion-ion collisions is negligible compared to electron-ion
radiation.  This is because (i) when identical particles scatter, the first time
derivative of the electric dipole moment is conserved, so only radiation of
quadrupole order and higher [with intensity $\sim(v/c)^2$ less than dipole radiation]
can be emitted; (ii) the speeds and accelerations of colliding ions are far less than
those of electrons because of their much larger masses.

Equations (2.1.27) and (2.1.29), together with Figures~2.1.2 and~2.1.3, are
the chief results of astrophysical interest from this section.  When dealing with
atoms that are only partially ionized, one must sometimes correct these results
for screening of the nuclear charge by the bound electrons; see, e.g.,
\citet{Ginzburg1967}.  For temperatures approaching and exceeding $kT = m_e c^2$,
relativistic effects and electron-electron collisions modify the emissivity
(2.1.29) to the form
\begin{align}
e &= 5.2 \times 10^{20}\;
\frac{\text{ergs}}{\text{g\;sec}}\;
\left(\frac{C_A f_i f_e^2}{A}\right)
\left(\frac{\rho_0}{\text{g\;cm}^{-3}}\right)
T_K^{\,1/2}\bigl(1 + 4.4 \times 10^{-10}\,T_K\bigr)\,\bar{G}(T).
\tag{2.1.31}
\end{align}

\noindent (\citealt{Ginzburg1967}).


%%---------------------------------------------------------------------------
%%  SECTION 2.2 — Free-Bound Radiation
%%---------------------------------------------------------------------------

\subsection{Free-Bound Radiation}
\label{sec:2.2}

\textit{Basic references}\quad \citet{Jackson1962}; \citet{Ginzburg1967}.

\bigskip
\textit{Basic physics and formulas}

When an electron of kinetic energy
\begin{equation}
\tfrac{1}{2}m_e v^2 \leq Z^2 R_y
\end{equation}
impinges on an ion, there is a significant probability that, instead of scattering
with the emission of bremsstrahlung, it will get captured into a bound state with
an accompanying emission of ``free-bound'' radiation.

We shall not compute here the details of the radiation.  Instead we cite the
results of such a computation from the above references.

Consider a plasma at temperature $T$ and density of rest mass $\rho_0$.  Let $C_A$ be
the concentration by mass of a particular atom in the plasma; let $f_i$ be the fraction
of all such atoms ionized into a given state, with charge $Ze$; and let $n_q$ be the
principal quantum number of that state, and $E_i$ its
ionization energy.  Then the emissivity due to such transitions is
\begin{align}
\varepsilon_\nu &= 8c^2
\left(\frac{C_A f_i f_e^2}{An_q^5}\right)
\left(\frac{\alpha^2 m_e e^2}{m_p}\right)^{\!3/2}
\rho_0
\left(\frac{\rho_0}{3kT}\right)^{\!3/2}
e^{-(h\nu - E_i)/kT}\;G_{bf}
\notag\\[4pt]
&= 3.9 \times 10^{15}\;
\frac{\text{ergs}}{\text{g\;sec\;Hz}}\;
\left(\frac{C_A f_i f_e^2}{An_q^5}\right)
\left(\frac{\rho_0}{\text{g\;cm}^{-3}}\right)
T_K^{-3/2}\;e^{-(h\nu - E_i)/kT}\;G_{bf}.
\tag{2.2.1}
\end{align}

Here $G_{bf}$ is a ``bound-free Gaunt factor'', analogous to the ``free-free'' Gaunt
factors of \S\ref{sec:2.1}.  It depends on the kinetic energy
\begin{equation}
\tfrac{1}{2}m_e v^2 = h\nu - E_i
\tag{2.2.2}
\end{equation}
of the electron that is captured, and on the structure of the bound state (which
one characterizes by its quantum numbers $n_q, l_q, \ldots$):
\begin{equation}
G_{bf} = G_{bf}(h\nu - E_i;\;n_q, l_q, \ldots).
\tag{2.2.3}
\end{equation}

Gaunt factors for various bound-free transitions of astrophysical interest are
tabulated and graphed by \citet{Karzas1961}.  Energy conservation
requires that all photons emitted have energies greater than $E_i$; consequently,
\begin{equation}
G_{bf} = 0 \quad \text{for} \quad h\nu - E_i < 0.
\tag{2.2.4}
\end{equation}

In general, $G_{bf}$ ``turns on'' discontinuously at $h\nu = E_i$, with an initial value
between ${\sim}\,0.5$ and ${\sim}\,1.5$:
\begin{equation}
G_{bf} \simeq 1 \quad \text{for} \quad 0 < h\nu - E_i \lesssim E_i.
\tag{2.2.5}
\end{equation}

For $h\nu - E_i \gtrsim Z^2 R_y$, $G_{bf}$ typically remains within an order of magnitude of
unity; as $h\nu - E_i \geq 10\,Z^2 R_y$, $G_{bf}$ falls rapidly.
See \citet{Karzas1961}.

The total emissivity, $e = \int \varepsilon_\nu\,d\nu$, due to a given free-bound transition is
\begin{align}
e &= 8\sqrt{3}\,c^2
\left(\frac{C_A f_i f_e^2}{An_q^5}\right)
\left(\frac{\alpha^2 m_e e^2}{m_p^2}\right)^{\!3/2}
\rho_0
\left(\frac{2\pi m_e c^2}{3kT}\right)^{\!1/2}
T_K^{-1/2}\;\bar{G}_{bf}
\notag\\[4pt]
&= 4.2 \times 10^{26}\;
\frac{\text{ergs}}{\text{g\;sec}}\;
\left(\frac{C_A f_i f_e^2}{An_q^5}\right)
\left(\frac{\rho_0}{\text{g\;cm}^{-3}}\right)
T_K^{-1/2}\;\bar{G}_{bf}
\tag{2.2.6}
\end{align}

where $\bar{G}_{bf}$, the mean Gaunt factor, depends only on temperature and on the
bound state, and is approximately equal to one.
\begin{equation}
\bar{G}_{bf} = \int_0^\infty G_{bf}(h\nu - E_i = xkT;\;n_q, l_q, \ldots)\;e^{-x}\,dx \simeq 1.
\tag{2.2.7}
\end{equation}

The total emissivity (2.2.6) due to a given free-bound transition, compared
to the total emissivity (2.1.29) of free-free radiation from the same type of ion,
is
\begin{equation}
\frac{\varepsilon_{fb}}{\varepsilon_{ff}}
\simeq \frac{Z^2}{n_q^5}\,
\left(\frac{8 \times 10^5}{T}\right).
\tag{2.2.8}
\end{equation}

\citet{Ginzburg1967} states that for astrophysical plasmas the ratio of total free-bound
radiation to total free-free radiation (summed over all states and all ions)
generally is less than
\begin{equation}
\frac{\varepsilon_{fb\,\text{total}}}{\varepsilon_{ff\,\text{total}}}
< \left(\frac{8 \times 10^5\,\text{K}}{T}\right).
\tag{2.2.9}
\end{equation}

\subsection{Thermal Cyclotron and Synchrotron Radiation}
\label{sec:2.3}

\textit{Basic references}\quad Jackson~\cite{Jackson1962};\
Ginzburg~\cite{Ginzburg1967}.

\bigskip
\noindent\textit{Power radiated by a single electron}

\medskip
Consider an electron of speed~$v$ and total mass-energy~$\gamma m_e c^2$,
\begin{equation}
\gamma \equiv (1 - v^2)^{-1/2},
\tag{2.3.1}
\end{equation}
spiralling in a magnetic field of strength~$B$. The Lorentz 4-acceleration of the
electron is
\begin{equation}
a^0 = 0, \qquad
\mathbf{a} = (e / m_e c)\,\gamma\,\mathbf{v} \times \mathbf{B}.
\tag{2.3.2}
\end{equation}
Consequently, the total power radiated---as obtained from the standard
Lorentz-invariant equation,
\begin{equation}
\frac{dE}{dt} = \frac{2e^2}{3c^3}\,a^2
\tag{2.3.3}
\end{equation}

\noindent---is
\begin{equation}
\frac{dE}{dt} = \frac{2\,r_0^2}{3\,c}\,(\gamma v_\perp)^2\,B^2.
\tag{2.3.4}
\end{equation}
Here $v_\perp$ is the component of the electron velocity perpendicular to the magnetic
field. If the electron is nonrelativistic, $v \ll c$, this radiation is called ``cyclotron
radiation''. If the electron is ultrarelativistic, $\gamma \gg 1$, it is called ``synchrotron
radiation''.

\bigskip
\noindent\textit{Cyclotron radiation from a nonrelativistic plasma}

\medskip
Consider a plasma with temperature
\begin{equation}
kT \ll m_e c^2, \quad \text{i.e.}\ T \lesssim 6 \times 10^9\;\text{K},
\tag{2.3.5}
\end{equation}
and with a magnetic field of strength~$B$. As in the nonrelativistic case, radiation
from protons spiralling
in the magnetic field is smaller by a factor $(m_e/m_p)^3 \simeq 10^{-10}$ than that from
electrons, since
\begin{equation}
\frac{dE}{dt} \propto \frac{v^2}{m^2} \propto \frac{1}{m^3}\,.
\tag{2.3.6}
\end{equation}
(Recall: in thermal equilibrium the mean kinetic energies of protons and electrons
are the same.) Hence, we can ignore protons and other ions when calculating
cyclotron radiation. Let $f_e$ be the number of unbound electrons per baryon in
the plasma. Then the total emissivity (ergs per second per gram of plasma) is
\begin{equation}
\varepsilon = f_e \left\langle \frac{dE}{dt} \right\rangle
\bigg/\!\left(\frac{1}{m_p}\right),
\tag{2.3.7}
\end{equation}
where $\langle\;\rangle$ denotes an average over the Maxwell-Boltzmann velocity distribution
of the electrons. Since $\gamma = 1$, since 2 of the 3 spatial directions are orthogonal
to~$B$, and since the mean kinetic energy per electron is $\tfrac{3}{2}kT$, we have
\begin{equation}
\langle \gamma v_\perp^2 \rangle
= \frac{2}{3}\,\frac{2\,\bar{E}}{m_e}
= \frac{2}{3}\,\frac{3kT}{m_e}
= \frac{2kT}{m_e}\,.
\tag{2.3.8}
\end{equation}
Combining this with equations~(2.3.4) and~(2.3.7), we obtain for the total
emissivity of the plasma
\begin{align}
\varepsilon &= \frac{4}{3}\left(\frac{f_e\,r_0^2}{m_p}\right)
\left(\frac{kT}{m_e c^2}\right) B^2
\tag{2.3.9}\\
&= (0.32\;\text{ergs/g\,sec})\;f_e\;T_{\rm K}\;B_G^2\,,
\notag
\end{align}
where $T_{\rm K}$ is temperature in~${}^\circ$K and $B_G$ is magnetic field in Gauss.

Each electron emits its radiation monochromatically, with the cyclotron
frequency (orbital frequency of electron's spiral motion)
\begin{equation}
\nu_{\text{cyc}} = \frac{eB}{2\pi m_e c}
= (2.79\;\text{MHz})\;B_G\,.
\tag{2.3.10}
\end{equation}

Thus, if the magnetic field is uniform, the radiation will be monochromatic
with this frequency. But in Nature the magnetic field will always be inhomo\-geneous,
and the spectrum will show a spread proportional to the spread in~$B^2$.
(Of course, there are always other sources of spread in the spectrum---e.g., doppler
shifts and relativistic generation of harmonics.)

\textit{Synchrotron radiation from a relativistic plasma.}

\medskip
\noindent
Consider a plasma with temperature
\begin{equation}
kT \gg m_e c^2, \quad \text{i.e.}\ T \gg 6 \times 10^9\,\text{K},
\tag{2.3.11}
\end{equation}
and with a magnetic field of strength $B$.
As in the nonrelativistic case, radiation
from protons and ions is totally negligible.
The total emissivity due to electrons
and positrons (recall: at $kT \gg m_e c^2$, there will be many electron-positron pairs!)
is given, as before, by equation~(2.3.7); but now the factor $f_e$ must be the number
of electrons and positrons per baryon,\footnote{The case of an optically thin plasma
is of particular interest in astrophysics. The kinetic theory of the creation and
annihilation of positrons in this case is somewhat complex; see, e.g.,
\citet{BisnovatyiKogan1971}.}
and the relevant, ultrarelativistic Maxwell--Boltzmann average is

\begin{equation}
\langle (\gamma v_\perp)^2 \rangle = \tfrac{2}{3}\langle\gamma^2\rangle
= 8\left(\frac{kT}{m_e c^2}\right)^{\!2}.
\tag{2.3.12}
\end{equation}
[One evaluates $\langle\gamma^2\rangle$ using the distribution function in phase space
\begin{equation}
\mathscr{N} = \frac{dN}{d^3 x\,d^3 p} \propto e^{-E/kT} = e^{-\gamma m_e c^2/kT},
\tag*{}
\end{equation}
with
\begin{equation}
d^3 p = 4\pi c^{-3}E^2\,dE
\tag*{}
\end{equation}
in ultrarelativistic limit;
in particular
\begin{equation}
\langle\gamma^2\rangle = \frac{1}{(m_e c^2)^2}\,
\frac{\displaystyle\int_0^\infty E^2\,e^{-E/kT}\,E^2\,dE}
{\displaystyle\int_0^\infty e^{-E/kT}\,E^2\,dE}
= 12\left(\frac{kT}{m_e c^2}\right)^{\!2}.\,]
\tag*{}
\end{equation}

Combining equations~(2.3.12), (2.3.7), and~(2.3.4), we obtain for the total
emissivity of the plasma
\begin{align}
\varepsilon &= \frac{16}{3}\left(\frac{r_0^2 c}{m_p}\right)
\left(\frac{kT}{m_e c^2}\right)^{\!2} B^2 \notag \\
&= (2.2 \times 10^{-10}~\text{ergs/g\,sec})\,f_e\,T_K^2\,B_G^2.
\tag{2.3.13}
\end{align}

In the ultrarelativistic case the electrons do not emit monochromatically.
An electron of energy $\gamma m_e c^2$ beams most of its radiation into a forward cone of
half-angle $\alpha = 1/\gamma$ (special-relativistic ``headlight effect'').
Since the electron is
headed toward the observer with speed $v$ when its cone is directed toward him,
the cone sweeps past the observer in time
\begin{equation}
\Delta t = (1-v^2)\left(\frac{2\alpha}{\omega_{\text{cyc}}}\right) = (1-v^2)\,\frac{1}{\gamma^3\,\omega_{\text{cyc}}}
= \frac{1}{\gamma^3\,\omega_{\text{cyc}}}.
\tag{2.3.14}
\end{equation}
Here $\omega_{\text{cyc}}$ is the angular velocity of the electron in its spiralling orbit
\begin{equation}
\omega_{\text{cyc}} = \frac{eB}{\gamma m_e c}.
\tag{2.3.15}
\end{equation}

Thus, the radiation comes in short bursts, of duration $\sim\!\Delta t = 1/(\gamma^3\omega_{\text{cyc}})$,
separated by long intervals of time $2\pi/\omega_{\text{cyc}}$.
When Fourier analyzed, such
radiation must be concentrated near the ``critical frequency''
\begin{equation}
\nu_{\text{crit}} = \tfrac{3}{2}\,\nu_{\text{cyc}} = \gamma^2\,eB/m_e c.
\tag{2.3.16}
\end{equation}

A detailed calculation [chapter~14 of \citet{Jackson1962}] reveals a fairly broad
spectrum which rises as $\nu^{1/3}$ at low frequencies, which peaks at
\begin{equation}
\nu_{\text{peak}} = 0.29\,\nu_{\text{crit}},
\tag{2.3.17}
\end{equation}
and which decays exponentially $\sim\nu^{1/2}\,e^{-2\nu/\nu_{\text{crit}}}$, at $\nu \gg \nu_{\text{crit}}$.
When averaged over the relativistic plasma, such a spectrum will have a broad
peak, with maximum located at
\begin{equation}
\nu_{\text{peak, plasma}} = \frac{\gamma_{\text{peak}}^2\,eB}{4\,m_e c}
\simeq 12(kT/m_e c^2)^2
\quad\left[\text{value of $\gamma_{\text{peak}}^2$ for electrons of}\right]
\tag{2.3.18}
\end{equation}
\begin{equation}
\nu_{\text{peak, plasma}} \simeq (100~\text{MHz})\,T_{10}^2 B_G.
\tag*{}
\end{equation}
Here $T_{10}$ is temperature in units of $10^{10}$\,K.
At $\nu \ll \nu_{\text{peak, plasma}}$
the spectrum will rise as a power law.

%%% ====================================================================
%%% SECTION 2.4 — Electron Scattering of Radiation
%%% ====================================================================

\subsection{Electron Scattering of Radiation}
\label{sec:2.4}

\textit{Basic references}

\citet{Jackson1962}, \S\,14.7 of \citet{Jackson1962b};
\S\,12.3 of \citet{Leighton1959};
\citet{Kompaneets1957}; \citet{Weyman1965};
\citet{Sunyaev1973}.

\medskip
\noindent
\textit{Basic physics and formulas}

\medskip
\noindent
In discussing electron scattering, we shall confine attention to the nonrelativistic
case---i.e.\ we shall demand that the temperature of the gas and the photon
frequencies satisfy
\begin{equation}
kT \ll m_e c^2 \quad (T \ll 6 \times 10^9\,\text{K}); \quad
h\nu \ll m_e c^2 = 500~\text{keV}.
\tag{2.4.1}
\end{equation}

The differential cross section for a free electron to scatter a photon has the
form
\begin{equation}
\frac{d\sigma_e}{d\Omega} = \tfrac{1}{2}\,r_0^2\,(1 + \cos^2\theta),
\tag{2.4.2}
\end{equation}
where $\theta$ is the angle between incoming and outgoing photon.  The total
cross section (``Thomson cross section''), obtained by integrating over all
directions, is
\begin{equation}
\sigma_T = (8\pi/3)\,r_0^2 = 0.665 \times 10^{-24}~\text{cm}^2.
\tag{2.4.3}
\end{equation}
It is important for astrophysical applications that the scattering cross section
is ``color blind'' (no dependence on frequency).  See, e.g., \S\,4.

It is also important that the differential cross section is symmetric between
forward and backward directions.  This guarantees, for example, that on the
average a scattered photon transmits \textit{all} of its momentum to the electron
(Here $\mathbf{n}_s$ is a unit vector in the direction of the initial photon motion).  By
contrast, only a tiny fraction of the photon energy is transmitted to the electron.
In a frame where the electron is initially at rest, it receives a kinetic energy
\begin{subequations}
\begin{equation}
\Delta E_\gamma = (h\nu/m_e c^2)\,h\nu\,(1-\cos\theta).
\tag{2.4.4a}
\end{equation}
\begin{equation}
(\Delta\mathbf{p}_e) = \mathbf{p}_\gamma = (h\nu/c)\,\mathbf{n}_s.
\tag{2.4.4b}
\end{equation}
\end{subequations}
(Recall: we have assumed $h\nu/m_e c^2 \ll 1$.)

Consider a monochromatic beam of photons, with frequency $\nu$, moving
through a thermalized plasma of temperature $T$.  Scattering by protons and ions
will be negligible compared to scattering by electrons, since
\begin{equation}
\sigma_e \propto r_0^2 = (e^2/m_e c^2)^2 \propto 1/(\text{mass of scatterer})^2.
\tag{2.4.5}
\end{equation}

On the average, how much energy $\langle\Delta E_\gamma\rangle$ is transmitted to the electron in a
single scattering?  The answer for electrons nearly at rest ($kT \ll h\nu$) is obtained
by averaging expression~(2.4.4b) over all angles.  Because of the forward-backward
symmetry of the differential cross section, $\cos\theta$ averages to zero and one obtains
\begin{equation}
\langle\Delta E_\gamma\rangle = (h\nu/m_e c^2)\,h\nu, \quad \text{if } kT \ll h\nu.
\tag{2.4.6}
\end{equation}

One can calculate $\langle\Delta E_\gamma\rangle$ for higher temperatures by invoking conservation
of 4-momentum, and averaging over all angles and electron speeds.  However,
such a calculation is rather long and messy.  To obtain the answer more easily,
one can use the following trick: An examination of the law of 4-momentum
conservation convinces one that $\langle\Delta E_\gamma\rangle$ must have the temperature dependence
\begin{equation}
\langle\Delta E_\gamma\rangle \propto (h\nu - \alpha kT),
\tag{2.4.7}
\end{equation}
where $\alpha$ is a constant to be calculated.  (Thus, if $h\nu > \alpha kT$, the electrons get
heated by the photons; if $h\nu < \alpha kT$, they get cooled.)
This law will reduce to
expression~(2.4.6) for $kT \ll h\nu$ if and only if the proportionality constant is
$h\nu/m_e c^2$:
\begin{equation}
\langle\Delta E_\gamma\rangle = (h\nu/m_e c^2)(h\nu - \alpha kT).
\tag*{}
\end{equation}

To calculate the constant $\alpha$, imagine the following experiment.  Place a large
number of photons and electrons into a box with perfectly reflective walls.
Require that the photons and electrons interact only by electron scattering.  Then
photons cannot be created or destroyed; only scattered.  Hence, when thermal
equilibrium is reached, the photons acquire a Boltzmann energy distribution,
rather than a Planck distribution.  (Recall that stimulated emissions are responsible
for Planckian deviations from the Boltzmann law; \S\,2.5.)  Since photons have
zero rest mass, their Boltzmann distribution
\begin{equation}
\mathscr{N} = \frac{dN}{d^3 x\,d^3 p} \propto e^{-h\nu/kT},
\tag*{}
\end{equation}
corresponds to a number of photons per unit frequency given by
\begin{equation}
\frac{dN}{d\nu} \propto \nu^2\,e^{-h\nu/kT},
\tag{2.4.8}
\end{equation}
and corresponds to
\begin{equation}
\langle h\nu\rangle = 3kT, \qquad
\langle (h\nu)^2\rangle = 12(kT)^2.
\tag*{}
\end{equation}
Thus, for our thought experiment the energy transfer in each collision, expression
(2.4.7) averaged over the equilibrium photon distribution, is
\begin{equation}
\langle\langle\Delta E_\gamma\rangle\rangle = (3kT/m_e c^2)(4-\alpha)\,kT.
\tag*{}
\end{equation}

But in equilibrium there must be no average energy transfer between photons
and electrons; $\langle\langle\Delta E_\gamma\rangle\rangle$ must be zero.
This condition tells us that $\alpha = 4$.
Thus, turning back to the general situation, we conclude that monochromatic
photons passing through a plasma of temperature $T$ transfer an average energy
per collision
\begin{equation}
\langle\Delta E_\gamma\rangle = (h\nu/m_e c^2)(h\nu - 4kT)
\tag{2.4.9}
\end{equation}
to the electrons.  When $4kT \gg h\nu$, the photon energies get boosted by the
collisions, and one says that the radiation is being ``Comptonized''; see \S\,5.10;
see, e.g., \citet{Illarionov1972}.

This concludes, for the moment, our discussion of the interaction between
radiation and plasmas.  We must point out that we have ignored a number of
plasma effects and instabilities which can be important in astrophysical
situations.  See, e.g., \citet{Bekefi1966} and \citet{KaplanTsytovich1972}.

%%% ====================================================================
%%% SECTION 2.5 — Hydrodynamics and Thermodynamics
%%% ====================================================================

\subsection{Hydrodynamics and Thermodynamics}
\label{sec:2.5}

\textit{Basic references}\quad
\S\S\,22.2 and 22.3 of \citet{MTW1973};
\citet{Lichnerowicz1967} and references cited therein.
We adopt the notational conventions of \citet{MTW1973},
including signature ``$-+++$''; geometrized units ($c = G = 1$); Greek
indices ranging from 0 to 3 and Latin from 1 to 3; ``hats'' on indices that refer
to local orthonormal frames, e.g.\ $\mu^{\hat{\alpha}}$ and $T^{\hat{\alpha}\hat{\beta}}$; extra-bold, sans-serif type for
4-vectors and 4-tensors, e.g.\ $\fvec{u}$ and $\mathbf{T}$; normal bold-face type for 3-vectors and
3-tensors (local Euclidean geometry).

\medskip
\noindent
\textit{Parameters describing the ``fluid.''}

\medskip
\noindent
``LRF'' $=$ local rest frame of the baryons; i.e.\ local orthonormal frame in which
there is no net baryon flux in any direction.

\begin{itemize}
\item[$\fvec{u}$] 4-velocity of LRF; i.e.\ 4-velocity of the ``fluid''.
\item[$n$] number density of baryons (number of baryons per unit volume) as
measured in LRF.
\item[$m_B$] mean rest mass of a baryon in the fluid; i.e.
\end{itemize}
\begin{equation}
m_B = \left(\!\!
\begin{array}{c}
\text{rest mass of}\\
\text{hydrogen atom}
\end{array}\!\!\right)
\times
\left(\!\!
\begin{array}{c}
\text{fraction of all baryons that are}\\
\text{in the form of hydrogen nuclei}
\end{array}\!\!\right)
+
\frac{1}{4}
\left(\!\!
\begin{array}{c}
\text{rest mass of}\\
\text{helium atom}
\end{array}\!\!\right)
\times
\left(\!\!
\begin{array}{c}
\text{fraction of all baryons that}\\
\text{are in helium nuclei}
\end{array}\!\!\right)
+ \cdots
\tag{2.5.1}
\end{equation}

\textit{Note}: in this section we restrict ourselves to fluids that are chemically
homogeneous and in which no nuclear reactions occur (``standard
fluid'').  Thus, $m_B$ is constant.

\begin{align}
\rho_0 &\equiv m_B\,n, && \text{rest-mass density, defined by}
\tag{2.5.2} \\[4pt]
\mathcal{V} &\equiv 1/n, && \text{specific volume, i.e.\ volume per baryon, defined by}
\tag{2.5.3} \\[4pt]
\mathcal{V}_0 &= 1/\rho_0, && \text{specific volume, i.e.\ volume per unit rest mass, defined by}
\tag{2.5.4}
\end{align}

\begin{align}
\rho &= \rho_0(1 + \Pi), && \text{total density of mass-energy, as measured in LRF.}
\tag{2.5.5}\\[4pt]
\Pi &\quad && \text{specific internal energy, defined by}
\tag*{}
\end{align}

\textit{Note}: here and throughout this section we set the speed of light equal to
one.
\begin{align}
p  &\quad && \text{isotropic pressure, as measured in LRF.}
\tag*{} \\
T  &\quad && \text{temperature, as measured in LRF.}
\tag*{} \\
s  &\quad && \text{entropy per baryon, as measured in LRF; of course,}
\tag*{} \\
s_0 &= m_B^{-1}\,s. && \text{entropy per unit mass, as measured in LRF.}
\tag{2.5.6}
\end{align}

$\mu$ \quad chemical potential, as measured in LRF; defined by
%
\begin{equation}
\tag{2.5.7}
\mu \equiv \left(\frac{\partial\rho}{\partial n}\right)_{\!s}
  = \frac{\rho + p}{n}\,.
\end{equation}

(The second equality follows from the first law of thermodynamics,
below.)

\medskip
\noindent
$\mathbf{q}$ \quad flux of energy (due to heat conduction, radiation,
convection, etc.)\ as measured in LRF.  This flux
(ergs\,cm$^{-2}$\,sec$^{-1}$) is a purely spatial vector as measured
in LRF; i.e.
%
\begin{equation}
\tag{2.5.8}
\mathbf{q} \cdot \mathbf{u} = 0.
\end{equation}

\noindent
$\mathbf{S}$ \quad entropy density-flux vector.  In LRF this vector has
time component equal to the entropy density, $S^0 = ns$, and space
components equal to the entropy flux, $S^j = q^j / T$.  Hence, in
frame-independent notation
%
\begin{equation}
\tag{2.5.9}
\mathbf{S} = ns\,\mathbf{u} + \mathbf{q}/T.
\end{equation}

\noindent
Note that
%
\begin{equation}
\tag{2.5.10}
\boldsymbol{\nabla}\cdot\mathbf{S}
= \left\{
\begin{array}{l}
\text{rate at which entropy is being generated}\\
\text{per unit volume as measured in LRF}
\end{array}
\right\}.
\end{equation}

\bigskip
\noindent
\emph{First law of thermodynamics}

\smallskip
\noindent
Follow a fluid element, containing $A$ baryons, along its world tube.
The total mass-energy in the fluid element, $\rho A/n$, changes as a
result of compression (change in volume, $A/n$) and as a result of
influx of heat:
%
\begin{equation}
\tag{2.5.11}
d(\rho A/n) = -p\,d(A/n) + T\,d(As).
\end{equation}

\noindent
[Here and throughout this section we assume for simplicity no internal
generation of entropy in the fluid element---e.g., no irreversible
chemical reactions; this enables us to write the influx of heat in
terms of the change in entropy, $T\,d(As)$.]

One often uses $A = \text{const.}$ and $\mu \equiv (\rho + p)/n$ to
rewrite this first law of thermodynamics in the equivalent forms
%
\begin{align}
\tag{2.5.12}
d\rho &= \frac{\rho + p}{n}\, dn + nT\, ds, \\[4pt]
\tag{2.5.12$'$}
d\mu  &= V\, dp + T\, ds.
\end{align}

\noindent
From the first law one can read off partial derivatives, e.g.
%
\begin{equation}
\tag{2.5.12$''$}
(\partial\rho/\partial n)_s = (\rho + p)/n, \qquad
(\partial\mu/\partial p)_s = T.
\end{equation}

\bigskip
\noindent
\emph{Fundamental relation and equations of state}

\smallskip
\noindent
The peculiar thermodynamic properties of the particular fluid being
studied are determined by a ``fundamental thermodynamic relation''
%
\begin{equation}
\tag{2.5.13}
\rho = \rho(n,\,s) \qquad \text{or} \qquad
\mu = \mu(p,\,s).
\end{equation}

\noindent
Once this relation has been specified---and once the constant $m_B$ has
been specified---one can use the first law of thermodynamics
(2.5.12) or (2.5.12$'$) to derive explicit expressions
for all other thermodynamic variables as functions of $n$, $s$ or $p$,
$s$.  For example, by combining with the first law one can derive the
``equations of state''
%
\begin{equation}
\tag{2.5.8$'$}
T(n,s) = \frac{1}{n}\left(\frac{\partial\rho}{\partial s}\right)_{\!n},
\qquad
p(n,s) = n\!\left(\frac{\partial\rho}{\partial n}\right)_{\!s} - \rho.
\end{equation}

\bigskip
\noindent
\emph{Adiabatic index and speed of sound}

\smallskip
\noindent
One defines the adiabatic index $\Gamma_1$ by
%
\begin{equation}
\tag{2.5.14}
\Gamma_1
\equiv \left(\frac{\partial\ln p}{\partial\ln n}\right)_{\!s}
= \frac{\rho + p}{p}\,
  \left(\frac{\partial p}{\partial\rho}\right)_{\!s}\,,
\end{equation}
%
where the third equality follows from the first law of thermodynamics.
It turns out that weak adiabatic perturbations (weak ``sound waves'')
propagate, in the LRF, with ordinary velocity
%
\begin{equation}
\tag{2.5.15}
c_s
= \left[\left(\frac{\partial p}{\partial\rho}\right)_{\!s}\right]^{1/2}
= \left(\frac{\Gamma_1\, p}{\rho + p}\right)^{\!1/2}.
\end{equation}

\bigskip
\noindent
\emph{Second law of thermodynamics}

\smallskip
\noindent
Equation~(2.5.10) allows one to write the second law of thermodynamics
in the form
%
\begin{equation}
\tag{2.5.16}
\boldsymbol{\nabla}\cdot\mathbf{S} \geq 0.
\end{equation}

\bigskip
\noindent
\emph{Decomposition of 4-velocity}

\smallskip
\noindent
One decomposes the gradient of the 4-velocity, $\nabla\mathbf{u}$,
into its ``irreducible'' tensorial parts:
%
\begin{equation}
\tag{2.5.17}
u_{\alpha;\beta}
= \sigma_{\alpha\beta} + \omega_{\alpha\beta}
  + \tfrac{1}{3}\,\theta\, h_{\alpha\beta}
  - a_\alpha\, u_\beta\,.
\end{equation}

\noindent
Here $\mathbf{a}$ is the 4-acceleration of the fluid
%
\begin{equation}
\tag{2.5.18a}
\mathbf{a} \equiv \nabla_{\mathbf{u}}\,\mathbf{u},
\qquad\text{i.e.}\quad
a^\alpha = u^\alpha{}_{;\beta}\, u^\beta\,,
\end{equation}
%
and $-a_\alpha u_\beta$ is that portion of $u_{\alpha;\beta}$ which is
not orthogonal to $\mathbf{u}$.  The remainder of $u_{\alpha;\beta}$
(i.e.\ $\sigma_{\alpha\beta} + \omega_{\alpha\beta} +
\tfrac{1}{3}\theta\, h_{\alpha\beta}$) is orthogonal to $\mathbf{u}$;
i.e.\ in the LRF it has only spatial components.  This orthogonal part
is decomposed into an isotropic part,
%
\begin{equation}
\tag{2.5.18b}
\theta \equiv \nabla \cdot \mathbf{u}
= u^\alpha{}_{;\alpha}\,,
\end{equation}
%
expansion, $\tfrac{1}{3}\theta h_{\alpha\beta}$, where $\theta$ is the
``expansion''; and $h_{\alpha\beta}$ is the ``projection tensor''
%
\begin{equation}
\tag{2.5.18c}
h_{\alpha\beta} \equiv g_{\alpha\beta} + u_\alpha\, u_\beta\,;
\end{equation}
%
plus a symmetric, trace-free ``shear''
%
\begin{equation}
\tag{2.5.18d}
\sigma_{\alpha\beta}
= \tfrac{1}{2}(u_{\alpha;\mu}\, h^\mu{}_\beta
              + u_{\beta;\mu}\, h^\mu{}_\alpha)
  - \tfrac{1}{3}\theta\, h_{\alpha\beta}\,;
\end{equation}
%
plus an antisymmetric ``rotation'' or ``vorticity''
%
\begin{equation}
\tag{2.5.18e}
\omega_{\alpha\beta}
= \tfrac{1}{2}(u_{\alpha;\mu}\, h^\mu{}_\beta
              - u_{\beta;\mu}\, h^\mu{}_\alpha).
\end{equation}

An observer in the LRF sees all fluid elements in his neighborhood to
move with low (nonrelativistic) velocities.  Let him use standard
Newtonian methods to calculate or measure the expansion $\theta$; shear
$\sigma^{\hat{j}}_{\;\hat{k}}$; and rotation $\omega^{\hat{j}}_{\;\hat{k}}$
of the fluid, and let a relativist calculate $\theta$,
$\sigma^{\hat{j}}_{\;\hat{k}}$, and
$\omega^{\hat{j}}_{\;\hat{k}}$ in the LRF from the above equations.
The two calculations will give the same answers.
[See, e.g., Ellis (1971).]

\bigskip
\noindent
\emph{Frozen-in magnetic field}

\smallskip
\noindent
Consider an ionized plasma which contains a ``frozen-in'' magnetic field.
The field is pure magnetic (no electric field) in the LRF.  Therefore
it can be described by a magnetic-field 4-vector $\mathbf{B}$
orthogonal to $\mathbf{u}$,
%
\begin{equation}
\tag{2.5.19}
\mathbf{B}\cdot\mathbf{u} = 0.
\end{equation}

\noindent
[For simplicity we assume that the magnetic permeability of the plasma
is the same as that of vacuum; so we do not distinguish between
``$\mathbf{B}$'' and ``$\mathbf{M}$''.  For a more general treatment, see
Lichnerowicz (1967).]  As the fluid moves, carrying with it the
frozen-in B-field, $\mathbf{B}$ must change as
%
\begin{equation}
\tag{2.5.20}
\frac{DB^\alpha}{d\tau}
= u_{\alpha;\beta}\, B^\beta
  + (\sigma_{\alpha\beta} - \tfrac{2}{3}\theta\, h_{\alpha\beta})\, B^\beta\,.
\end{equation}

\noindent
The term $u_{\alpha;\beta}\, B^\beta$ is required to keep $\mathbf{B}$
orthogonal to $\mathbf{u}$; by the term $\sigma_{\alpha\beta}\, B^\beta$,
the rotation of the fluid rotates the field lines; by the term
$-\tfrac{2}{3}\theta\, h_{\alpha\beta}\, B^\beta$, the compression of the
fluid orthogonal to $\mathbf{B}$ magnifies $\mathbf{B}$, conserving
flux.

\bigskip
\noindent
\emph{Stress-energy tensor}

\smallskip
\noindent
Consider a fluid with isotropic pressure, with shear and bulk viscosity,
with energy flowing between fluid elements, and with a frozen-in
magnetic field.  The stress-energy tensor for such a fluid is
%
\begin{align}
\tag{2.5.21}
T^{\alpha\beta}
&= \rho u^\alpha u^\beta + (p - \zeta\theta)\, h^{\alpha\beta}
   - 2\eta\, \sigma^{\alpha\beta}
   + q^\alpha u^\beta + u^\alpha q^\beta + \pi^{\alpha\beta}
\notag \\
&\quad
   + \frac{1}{8\pi}(B^2 u^\alpha u^\beta + B^2 h^{\alpha\beta}
   - 2B^\alpha B^\beta).
\end{align}

\noindent
The term $\rho u^\alpha u^\beta$ is the total density of mass-energy
(excluding only that of the frozen-in B-field) as measured in the LRF.
The term $p\, h^{\alpha\beta}$ is the isotropic pressure that would be
measured in the LRF if the gas were not changing volume (if $\theta$
were zero).  The quantities $\zeta$ and $\eta$ are the
\emph{coefficients of bulk viscosity and of dynamic viscosity},
respectively.  The term $-\zeta\theta\, h^{\alpha\beta}$ is the
isotropic viscous stress which resists isotropic expansion
($\theta > 0$) or compression ($\theta < 0$) of the fluid.
The term $-2\eta\, \sigma^{\alpha\beta}$ is the viscous shear stress
which resists shearing motions.  The term
$q^\alpha u^\beta + u^\alpha q^\beta$ is the energy density and
stresses associated with the flowing energy $\mathbf{q}$, as measured
in the LRF, are here neglected by comparison with $\rho u^\alpha u^\beta$
and $p\, h^{\alpha\beta}$.

(\emph{Note:} the energy density and stresses associated with the
flowing energy $\mathbf{q}$, as measured in the LRF, are here
neglected by comparison with $\rho u^\alpha u^\beta$ and
$p\, h^{\alpha\beta}$.  This neglect is valid with enormous accuracy in
most contexts of interest in these lectures.  An exception is the
radiation pressure which produces ``self-regulation'' of accretion onto
black holes when the inflowing mass is sufficiently large; see
\S\S\,4.5 and 5.13.)  The term
%
\begin{equation}
\tag{2.5.22}
T^{\alpha\beta}_{\text{MAG}}
= \frac{1}{8\pi}(B^2 u^\alpha u^\beta + B^2 h^{\alpha\beta}
  - 2B^\alpha B^\beta)
\end{equation}
%
is the Maxwell stress-energy associated with the frozen-in B-field (in
LRF: energy density $B^2/8\pi$, pressure $B^2/8\pi$ orthogonal to
field lines; tension $-B^2/8\pi$ along field lines).

\bigskip
\noindent
\emph{Equations of hydrodynamics}

\smallskip
\noindent
The fundamental equations governing the motion of a fluid in a given
gravitational field (spacetime geometry) are (i) the law of
\emph{baryon conservation}
%
\begin{equation}
\tag{2.5.23}
\nabla \cdot (n\,\mathbf{u}) = 0,
\quad\text{i.e.}\quad dn/d\tau = -\theta\, n,
\end{equation}
%
or equivalently rest-mass conservation
%
\begin{equation}
\tag{2.5.23$'$}
\nabla \cdot (\rho_0\, \mathbf{u}) = 0,
\quad\text{i.e.}\quad d\rho_0/d\tau = -\theta\, \rho_0\,;
\end{equation}

\noindent
(ii) the law of \emph{local energy conservation}
%
\begin{equation}
\tag{2.5.24}
\mathbf{u}\cdot(\nabla\cdot\mathbf{T}) = 0;
\end{equation}

\noindent
(iii) the \emph{Euler equations} (i.e.\ law of local momentum
conservation)
%
\begin{equation}
\tag{2.5.25}
\mathbf{h}\cdot(\nabla\cdot\mathbf{T}) = 0;
\end{equation}

\noindent
(iv) the laws of thermodynamics, (2.5.11)--(2.5.16); (v) the laws of
energy transport for the frozen-in B-field, (2.5.20); the laws of
energy transport, which govern $\mathbf{q}$ (see \S\,2.6, 2.6.43).

\bigskip
\noindent
\emph{Law of local energy conservation}

\smallskip
\noindent
When one evaluates the law of local energy conservation,
$\mathbf{u}\cdot(\nabla\cdot\mathbf{T}) = 0$,
for the stress-energy tensor (2.5.21), one finds that the Maxwell
stress-energy
%
\begin{equation}
\tag{2.5.26}
\mathbf{u}\cdot(\nabla\cdot\mathbf{T}_{\text{MAG}}) = 0
\end{equation}
%
gives zero contribution (work done to compress magnetic field is
precisely equal to increase in magnetic field energy); and that the
remainder of the stress-energy tensor gives
%
\begin{equation}
\tag{2.5.27}
d\rho/d\tau = -(\rho + p)\,\theta
+ \zeta\,\theta^2 + 2\eta\,\sigma_{\alpha\beta}\,\sigma^{\alpha\beta}
- \nabla\cdot\mathbf{q} - \mathbf{a}\cdot\mathbf{q}\,.
\end{equation}

\noindent
Here $d/d\tau$ is derivative with respect to proper time along the
world lines of the fluid; i.e., in the language of the differential
geometer, $d/d\tau = \mathbf{u}$.  The terms $-(\rho + p)\,\theta$ is
the increase in mass-energy density due to compression.  The term
$\zeta\,\theta^2$ and $2\eta\,\sigma_{\alpha\beta}\,\sigma^{\alpha\beta}$
are the increases in mass-energy density due to viscous heating
(conversion of relative kinetic energy of adjacent fluid elements into
heat).  The term $-\nabla\cdot\mathbf{q}$ is the influx of
mass-energy from neighboring fluid elements.  The term
$-\mathbf{a}\cdot\mathbf{q}$ is a special relativistic correction to
the ``redshift'' of $\mathbf{q}$.  To understand this term, consider a
flux of energy $\mathbf{q}$---or, equivalently, with the inertia of
the energy flux $\mathbf{q}$ associated with it, a flux of mass-energy
from the viewpoint of an accelerated fluid (photons do not interact with
them).  The acceleration gives rise to a redshift (a gravitational
redshift) and hence to a nonzero $\nabla\cdot\mathbf{q}$ (no
interaction between fluid and photons), one must include the correction
factor $-\mathbf{a}\cdot\mathbf{q}$.  For a similar reason, a similar
relativistic correction appears in the law of heat conduction.

\bigskip
\noindent
\emph{Euler equation for a perfect fluid}

\smallskip
\noindent
Consider a perfect fluid, flowing adiabatically through spacetime
($\zeta = \eta = \mathbf{B} = \mathbf{q} = 0$).  In this case the
Euler equations $\mathbf{h}\cdot(\nabla\cdot\mathbf{T}) = 0$ reduce to
%
\begin{equation}
\tag{2.5.30}
(\rho + p)\,\mathbf{a} = -\nabla p - \mathbf{u}\,(\mathbf{u}\cdot\nabla p).
\end{equation}
%
In words:
\[
(\rho + p)\,\mathbf{a}
= -\left(\!\text{inertial mass per}\atop\text{unit volume}\!\right)
  \times\text{(4-acceleration)}
= -\left(\!\text{pressure gradient}\atop\text{projected orthogonal to }\mathbf{u}\!\right).
\]

\textit{Bernoulli equation}

Consider a perfect fluid undergoing stationary, adiabatic flow in a stationary
spacetime.  More particularly, set $\zeta = \eta = \fvec{B} = \fvec{q} = 0$; assume that spacetime is
endowed with a Killing vector field $\boldsymbol{\xi}$,

\begin{equation}
\xi_{\alpha;\beta} + \xi_{\beta;\alpha} = 0\,,
\tag{2.5.31}
\end{equation}

which need not be timelike; and assume that the flow is adiabatic, $ds/d\tau = 0$,
and stationary in the sense that

\begin{equation}
\mathscr{L}_{\boldsymbol{\xi}}\,p = 0, \quad
\nabla_{\boldsymbol{\xi}}\,\rho = \nabla_{\boldsymbol{\xi}}\,p = \cdots = 0\,.
\tag{2.5.32}
\end{equation}

Here $\mathscr{L}_{\boldsymbol{\xi}}$ is the Lie derivative along $\boldsymbol{\xi}$.  In this case the Euler equations~(2.5.30),
together with the first law of thermodynamics~(2.5.12), imply the relativistic
Bernoulli equation

\begin{equation}
d(\mu\,\fvec{u}\cdot\boldsymbol{\xi})/d\tau = 0\,;
\tag{2.5.33}
\end{equation}

i.e., $\mu\,\fvec{u}\cdot\boldsymbol{\xi}$ is constant along flow lines.

\bigskip
\textit{Newtonian limit}

The Newtonian limit of relativistic hydrodynamics is obtained when, in a nearly
global Lorentz frame, the following approximations hold:

\begin{gather}
g_{00} = -(1 + 2\Phi), \quad |\Phi| \ll 1, \quad
B^2/\rho_0 \ll 1,
\notag\\
p/\rho_0 \ll 1, \quad \Pi \ll 1, \quad v^2 \ll 1\,.
\tag{2.5.34}
\end{gather}

Here $\Phi$ is the Newtonian gravitational potential with sign $\Phi < 0$, and $v$ is the
ordinary velocity of the fluid

\begin{equation}
\fvec{v} \equiv v_j = (u^j/u^0)\,\fvec{e}_j\,.
\tag{2.5.35}
\end{equation}

In the Newtonian limit the chemical potential $\mu$ reduces to

\begin{equation}
\mu = m_B(1 + w)\,,
\tag{2.5.36}
\end{equation}

%%% ============================================================
%%% Book page 370
%%% (PDF p.15, left column)
%%% ============================================================

where $w$ is the \textit{enthalpy}

\begin{equation}
w = \Pi + p/\rho_0\,.
\tag{2.5.37}
\end{equation}

The Bernoulli equation~(2.5.33) reduces to the familiar form

\begin{equation}
\Phi + \tfrac{1}{2}v^2 + w = \text{constant along flow lines}\,.
\tag{2.5.38}
\end{equation}

The Euler equation for a perfect fluid in adiabatic flow, (2.5.30), reduces to

\begin{equation}
\frac{d\fvec{v}}{dt} = -\nabla\Phi - \frac{1}{\rho_0}\nabla p\,,
\tag{2.5.39}
\end{equation}

where $d/dt$, the derivative with respect to proper time along the flow lines, has
the Newtonian form

\begin{equation}
d/dt = \partial/\partial t + \fvec{v}\cdot\boldsymbol{\nabla}\,.
\tag{2.5.40}
\end{equation}

The first law of thermodynamics, (2.5.12) and (2.5.12$'$), reduces to

\begin{equation}
d\Pi = -p\,dV_0 + T\,ds_0\,,
\tag{2.5.41}
\end{equation}

\begin{equation}
dw = V_0\,dp + T\,ds_0\,.
\tag{2.5.42}
\end{equation}

(In Newtonian theory one usually adopts a per-unit-mass viewpoint rather than a
per-baryon viewpoint; and thus one uses $\rho_0$, $V_0$ rather than $n$, $\mathcal{V}$, $s_0$.)

%%% ============================================================
%%% Section 2.6: Radiative Transfer
%%% Book page 370–371 (PDF p.15, left–right columns)
%%% ============================================================

\subsection{Radiative Transfer}
\label{sec:2.6}

\textit{Basic references}\quad
Chandrasekhar (1960); Appendix~1 of Pacholczyk (1970); Mihalas (1970).

\bigskip
\textit{Notation and terminology}

At an arbitrary event in spacetime pick an arbitrary
local Lorentz frame.  In that frame pick an arbitrary spatial direction $\fvec{n}$ (a
unit vector; see Figure~2.6.1).  Examine the amount of energy $dE$ that is (i) carried
by photons across a unit surface area $dA$ orthogonal to $\fvec{n}$ (the surface $\mathscr{S}_m$ of
Figure~2.6.1) during unit time $dt$, with (ii) the photons having frequencies $\nu$ in
the range $d\nu$, and (iii) the photons being directed into a solid angle $d\Omega$ about
$\fvec{n}$.  The ratio

\begin{equation}
I_\nu = \frac{dE}{dt\;dA\;d\nu\;d\Omega}
\tag{2.6.1}
\end{equation}

is called the \textit{specific intensity} or simply the \textit{intensity}.  It depends on (a) the
event in spacetime; (b) the choice of Lorentz frame; (c) the direction $\fvec{n}$; (d)
the frequency $\nu$.  One also defines the \textit{total intensity} $I$ by

\begin{equation}
I = \int I_\nu\;d\nu = \frac{dE}{dt\;dA\;d\Omega}\,;
\tag{2.6.2}
\end{equation}

the \textit{average specific intensity} or \textit{average intensity} $J_\nu$ (averaged over all directions)
by

\begin{equation}
J_\nu = \frac{1}{4\pi}\int I_\nu\;d\Omega\,;
\tag{2.6.3}
\end{equation}

and the \textit{average total intensity} $J$ by

\begin{equation}
J = \frac{1}{4\pi}\int I\;d\Omega = \frac{1}{4\pi}\int I_\nu\;d\nu\;d\Omega\,.
\tag{2.6.4}
\end{equation}

It is easy to see that the \textit{energy density per unit frequency} in the radiation is

\begin{equation}
\rho^{(\text{rad})} = \frac{dE}{d^3x\;d\nu} = \frac{4\pi J_\nu}{c}\,,
\tag{2.6.5}
\end{equation}

and that the \textit{total energy density} in the radiation is

\begin{equation}
\rho^{(\text{rad})} = \frac{dE}{d^3x} = \frac{4\pi J}{c}\,.
\tag{2.6.6}
\end{equation}

%%% ============================================================
%%% Book page 371 (right column) — Figure 2.6.1, flux, invariance
%%% ============================================================

Pick a 2-surface $\mathscr{S}_m$ in the chosen Lorentz frame, with unit normal $\fvec{m}$ (Fig.\
2.6.1).  Examine the total energy $dE$ per unit area $dA$ that is carried across the
chosen surface during unit time $dt$, by photons having frequency $\nu$ in the range
$d\nu$.  Make no restrictions on the photon direction or solid angle; but count
negatively those photons that cross $dA$ from the front side toward the back.
The ratio

\begin{equation}
F_\nu = \frac{dE}{dt\;dA\;d\nu}
\tag{2.6.7}
\end{equation}

is called the \textit{specific flux}, or simply the \textit{flux}.  It is easy to see from Figure~2.6.1
that

\begin{equation}
F_\nu = \int I_\nu\,\cos\theta\;d\Omega\,.
\tag{2.6.8}
\end{equation}

\begin{figure}[htbp]
\centering
\fbox{\parbox{0.9\columnwidth}{\centering [Figure 2.6.1 placeholder]\\
The surfaces, directions, and solid angle used in the definition of intensity
and of flux.  $\mathscr{S}_m$ is a surface with unit normal $\fvec{m}$;
$\fvec{n}$ is a unit vector in the photon propagation direction;
$\theta$ is the angle between $\fvec{m}$ and $\fvec{n}$; $d\Omega$ is the solid angle element.}}
\caption{The surfaces, directions, and solid angle used in the definition of intensity
and of flux.}
\label{fig:2.6.1}
\end{figure}

The specific flux depends on (a) the chosen event in spacetime; (b) the chosen
Lorentz frame at that event; (c) the chosen 2-surface $\mathscr{S}_m$ in that Lorentz frame.

%%% ============================================================
%%% Book page 372 — Invariance of I_nu/nu^3, equation of radiative transfer
%%% (PDF p.16, left column)
%%% ============================================================

The integral of the flux over all frequencies is called the \textit{total flux}

\begin{equation}
F = \int F_\nu\;d\nu = \int I\,\cos\theta\;d\Omega\,.
\tag{2.6.9}
\end{equation}

Notice that the total flux is the ``first moment'' of the total intensity (the
``moment'' being taken with respect to the unit normal $\fvec{m}$).  One can readily
verify that the ``second moment'' is equal to the radiation pressure that acts
across the surface $\mathscr{S}_m$:

\begin{equation}
T^{(\text{rad})} \equiv \fvec{m}\cdot\fvec{T}^{(\text{rad})}\cdot\fvec{m}
= \int I\,\cos^2\theta\;d\Omega\,.
\tag{2.6.10}
\end{equation}

\bigskip
\textit{Invariance of $I_\nu/\nu^3$}

Choose a particular photon that passes through a given event in spacetime.
Observe that particular photon, and all others in the vicinity, from several different
Lorentz frames at the given event.  The frequency $\nu$ of the photon will depend
on the choice of Lorentz frame (Doppler shift from one frame to another).
The specific intensity $I_\nu$ in the frame neighborhood of the photon will also depend
on the choice of Lorentz frame.  However, the ratio $I_\nu/\nu^3$ will be Lorentz-
invariant. (Aside from a factor $h^{-4}$, $I_\nu/\nu^3$ is the invariant number density in
phase space

\begin{equation}
\mathscr{N} = \frac{dN}{d^3x\,d^3p} = \frac{1}{h^3}\,\frac{I_\nu}{\nu^3}\,;
\tag{2.6.11}
\end{equation}

see, e.g., \S22.6 of MTW.)

When photons propagate freely through curved spacetime (no interaction with
matter), the ratio $I_\nu/\nu^3$ is conserved along the world line of each photon
(Liouville's theorem).

\paragraph*{Equation of Radiative Transfer}

Consider the propagation of photons through a medium (e.g., through gas that
is falling into a black hole).\ Analyze the photon propagation from the viewpoint
of observers at rest in the medium (local rest frame; ``LRF'').\ At each event
denote by $\mathbf{u}$ the 4-velocity of the medium---and hence also of the LRF.\ Focus
attention on all photons in the (phase-space) neighborhood of a given null
geodesic ray.\ Denote by $\mathbf{p}$ the 4-momentum of that ray.\ Then at each event
along the given ray $\mathbf{p}$ is given by
%
\begin{equation}
\mathbf{p} = h\nu(\mathbf{u} + \mathbf{n}).
\tag{2.6.12}
\end{equation}
%
Here $\mathbf{n}$ is a unit vector;
%
\[
\mathbf{n} \cdot \mathbf{u} = 0,
\]
%
and (i) is interpreted in the LRF as the direction of propagation of the ray
(same as 3-vector $\mathbf{n}$ of Figure~2.6.1).\ Also, in eq.~(2.6.12), $\nu$ is the frequency of
any photon which propagates along the ray, and $h\nu$ is its energy, as measured in
the LRF.

Because of the acceleration and shear of the medium, the frequency
$\nu$ of the chosen ray changes from point to point along the ray.\ The change $d\nu$,
when the ray propagates a proper spatial distance $dl$ as seen in the LRF, is
%
\begin{equation}
d\nu = \nu(-\,\mathbf{p} \cdot \mathbf{u}/h)\,dl.
\tag{2.6.13}
\end{equation}
%
(This frame-independent equation is derived easily by geometrical arguments
in the LRF.)\ A straightforward calculation, using the geodesic equation for the
ray
%
\[
\boldsymbol{\nabla}_{\mathbf{p}}\,\mathbf{p} = 0,
\]
%
and using expansion (2.5.17) for $\boldsymbol{\nabla}\mathbf{u}$, reveals
%
\begin{equation}
d\nu/dl = -\nu(\mathbf{n} \cdot \mathbf{a} + \tfrac{1}{3}\theta
  + n^\alpha n^\beta \sigma_{\alpha\beta}).
\tag{2.6.14}
\end{equation}
%
The first term, $-\nu\,\mathbf{n} \cdot \mathbf{a}$, is the ``gravitational redshift'' produced by the acceleration
of the LRF; the second and third terms, $-\nu(\tfrac{1}{3}\theta + n^\alpha n^\beta \sigma_{\alpha\beta})$, are the
``cosmological redshift'' due to the expansion of the medium along the direction
of the ray.

If there were no interaction with the medium, then $I_\nu/\nu^3$ would be conserved
along the ray.\ Hence, the equation of radiative transfer along the given ray must
have the form
%
\begin{equation}
\frac{dI_\nu}{dl} - \left(\frac{3}{\nu}\,\frac{d\nu}{dl}\right) I_\nu
  = \text{(effects of interaction with medium)};
\tag{2.6.15}
\end{equation}
%
i.e.,
%
\[
dI_\nu/dl + (3\mathbf{n} \cdot \mathbf{a} + \theta
  + 3n^\alpha n^\beta \sigma_{\alpha\beta})\,I_\nu
  = \text{(interaction effects)}.
\]
%
Four types of interaction can occur: spontaneous emission of radiation by the
matter; stimulated emission; absorption; and scattering.

We shall assume that on all length scales of interest the medium is isotropic;
and we shall denote its emissivity by
%
\begin{equation}
\varepsilon_\nu \equiv \frac{dE}{d\nu\,dt\,dm_0}
  = \begin{pmatrix} \text{energy emitted spontaneously per unit frequency}\\
    \text{during unit time by a unit rest mass, integrated}\\
    \text{over all angles, and measured in LRF} \end{pmatrix}.
\tag{2.6.16}
\end{equation}
%
(The emissivities for free-free and free-bound transitions were discussed in \S\S\,2.1
and 2.2.)\ Then spontaneous emission contributes
%
\begin{equation}
\left(\frac{dI_\nu}{dl}\right)_{\!\text{spontaneous emission}}
  = \frac{1}{4\pi}\,\rho_0\,\varepsilon_\nu
\tag{2.6.17}
\end{equation}
%
to the specific intensity along the given ray.\ Here $\rho_0$ is the density of rest mass
in the medium.

The rate of absorption and the rate of stimulated emission are both proportional
to the intensity of the passing beam.\ Hence, the effects of absorption and of
stimulated emission can be lumped together into a single \textit{absorption coefficient}
$\kappa_\nu$, defined by
%
\begin{equation}
\left(\frac{dI_\nu}{dl}\right)_{\!\substack{\text{absorption plus}\\[2pt]
  \text{stimulated emission}}}
  = -\rho_0\,\kappa_\nu\,I_\nu.
\tag{2.6.18}
\end{equation}
%
The dimensions of the absorption coefficient are $\mathrm{cm}^2/\mathrm{g}$ (``absorption cross section
per unit mass'' at the given frequency).\ When absorption dominates over
stimulated emission, $\kappa_\nu$ is positive; when stimulated emission dominates,
$\kappa_\nu$ is negative (``negative absorption'').

Scattering is more complicated to treat than emission and absorption.\ Let
%
\begin{equation}
\frac{d\kappa_s}{d\Omega'\,d\nu'}\,(\mathbf{p},\,\mathbf{p}')
\tag{2.6.19}
\end{equation}
%
be the differential cross section (as measured in the LRF) for a unit rest mass to
convert a photon of momentum $\mathbf{p}$ into a photon of momentum $\mathbf{p}'$; and let
%
\[
\kappa_s(\nu) = \int \frac{d\kappa_s}{d\Omega'\,d\nu'}\,d\Omega'\,d\nu'
\]
%
be the total scattering cross-section (``scattering opacity'') per unit rest mass.
For example, in the case of electron scattering by a nonrelativistic plasma
($h\nu \ll m_e c^2$; $kT \ll m_e c^2$; see \S\,4.4), when one neglects the tiny change in photon
frequency the cross sections per unit rest mass are
%
\begin{equation}
\frac{d\kappa_s}{d\Omega'\,d\nu'}\,(\mathbf{p},\,\mathbf{p}')
  = \frac{f_e}{2m_p}\,\sigma_T\,\tfrac{3}{8\pi}\,
    [1 + (\mathbf{n} \cdot \mathbf{n}')^2]\,\delta(\nu' - \nu)
\tag{2.6.20}
\end{equation}
%
\[
\kappa_s = (f_e/m_p)\,\sigma_T = (0.40\;\mathrm{cm}^2/\mathrm{g})\,f_e
\]
%
(Recall: $n = \rho/h\nu$, $n' = \rho'/h\nu'$; $f_e$ is the number of free electrons per baryon in
the plasma.)\ Scattering, as described by the appropriate differential cross section,
can increase $I_\nu$ (scattering into the beam from other directions), or can decrease
$I_\nu$ (scattering out of the beam into other directions):
%
\begin{equation}
\left(\frac{dI_\nu}{dl}\right)_{\!\text{scattering}}
  = +\rho_0 \!\int \frac{d\kappa_s}{d\Omega\,d\nu}\,
    (\mathbf{p}',\,\mathbf{p})\,I_{\nu'}\,d\Omega'\,d\nu'
    \;-\; \rho_0\,\kappa_s\,I_\nu.
\tag{2.6.21}
\end{equation}
%
By combining equations (2.6.15)--(2.6.21), we obtain our final form of the
equation of radiative transfer
%
\begin{align}
dI_\nu/dl &+ (3\mathbf{n} \cdot \mathbf{a} + \theta
  + 3n^\alpha n^\beta \sigma_{\alpha\beta})\,I_\nu \notag\\
  &= +\rho_0 \!\int \frac{d\kappa_s}{d\Omega\,d\nu}\,
     (\mathbf{p}',\,\mathbf{p})\,I_{\nu'}\,d\Omega'\,d\nu'
     - \rho_0\,\kappa_s\,I_\nu \notag\\
  &\quad + \tfrac{1}{4\pi}\,\rho_0\,\varepsilon_\nu
     - \rho_0\,\kappa_\nu\,I_\nu.
\tag{2.6.22}
\end{align}

%%% -----------------------------------------------------------------
%%% Relationship between emissivity and absorption
%%% -----------------------------------------------------------------

\paragraph*{Relationship between emissivity and absorption}

Into an insulated box place a material medium and a radiation field; and then wait
until complete thermodynamic equilibrium is achieved.\ If $T$ is the equilibrium
temperature, then the radiation field as seen in the LRF will be isotropic with the
standard black-body intensity, $I_\nu = B_\nu$, where
%
\begin{equation}
B_\nu = \frac{2h\nu^3/c^2}{e^{h\nu/kT} - 1}.
\tag{2.6.23}
\end{equation}
%
In equilibrium there must be no net change of $I_\nu$ along any ray: scattering into
the beam must be completely balanced by scattering out of the beam; and
emission must be completely balanced by absorption.\ The requirement of
``detailed balance'' for emission and absorption can be met only if the emissivity
and the absorption coefficient are related by
%
\begin{equation}
\frac{dI_\nu}{dl} = \rho_0\,\kappa_\nu
  \!\left(\frac{\varepsilon_\nu}{4\pi\kappa_\nu}
  - I_\nu\right) = 0;
\tag{2.6.24}
\end{equation}
%
i.e.,
%
\[
\varepsilon_\nu / 4\pi\kappa_\nu = B_\nu.
\]
%
Since $\varepsilon_\nu$ and $\kappa_\nu$ depend only on the thermodynamic state of the matter, and
have nothing to do with the state of the radiation, relation~(2.6.24) must be
satisfied not only inside our insulated box, but also in all other cases where
the matter by itself is in thermodynamic equilibrium but the radiation might
not be.

For matter not in thermodynamic equilibrium one can use a more sophisticated
version of the principle of detailed balance (``Einstein A and B coefficients'') to
derive a more complicated relationship between the emissivity and the absorption
coefficient.\ See, e.g., Chandrasekhar~\cite{Chandrasekhar1960}.

As an application of the equilibrium relationship $\varepsilon_\nu/4\pi\kappa_\nu = B_\nu$, consider free-free
transitions in an ionized gas.\ Whenever the free electrons (which do the
emitting and absorbing) are in thermodynamic equilibrium with each other
(Maxwell--Boltzmann velocity distribution), the absorption coefficient---as
derived from the emissivity (2.1.27)---must be
%
\begin{equation}
\kappa_\nu^{\mathrm{ff}} = (1.50 \times 10^{25}\;\mathrm{cm}^5\;\mathrm{g}^{-1})
  \left(\frac{\rho_0}{\mathcal{A}}\right)
  \left(\frac{(eZ)^2}{\mathcal{A}}\right)
  \bar{G}\;T^{-7/2}\;\nu^{-3}
  \left(\frac{1 - e^{-x}}{x^{-1}}\right),
\tag{2.6.25}
\end{equation}
%
where
\[
x = h\nu/kT.
\]
%
See \S\,2.1 for notation.

Notice that for matter in thermodynamic equilibrium the relationship
$\varepsilon_\nu / 4\pi\kappa_\nu = B_\nu$ permits one to rewrite the equation
of radiative transfer~(2.6.22) in the form
%
\begin{equation}
\tag{2.6.26}
\frac{dI_\nu}{dl} + (3n\cdot a + \dot\theta + 3\sigma^{\alpha\beta}
n_\alpha n_\beta)\, I_\nu
= \rho_0\,\kappa_\nu\bigl(B_\nu - I_\nu\bigr)
  + \left.\frac{dI_\nu}{dl}\right|_{\text{scattering}}
  + \left.\frac{d(I_\nu/\nu^3)}{d\tau_\nu}\right|_{\text{scattering}} ;
\end{equation}
%
or, equivalently,
%
\begin{equation}
\tag{2.6.27}
\frac{d(I_\nu/\nu^3)}{d\tau_\nu}
= \rho_0\,\kappa_\nu \left(\frac{B_\nu}{\nu^3}
  - \frac{I_\nu}{\nu^3}\right)
  + \left.\frac{d(I_\nu/\nu^3)}{d\tau_\nu}\right|_{\text{scattering}}.
\end{equation}

%%% ---------------------------------------------------------------
%%% Optical depth
%%% ---------------------------------------------------------------

\bigskip
\noindent
\emph{Optical depth}

\medskip
\noindent
Consider radiation propagating out of a medium into surrounding empty space.
Follow a given ray ``backward'', from ``infinity'' into the medium.  Let the ray
have frequency $\nu_\infty$ at infinity; then its frequency at location~$l$
($l = {}$proper distance measured in LRF of medium) is
%
\begin{equation}
\tag{2.6.28}
\nu(l) = \nu_\infty \exp\!\left[\,-\int_i^l
\bigl(\dot\theta + n^\alpha n^\beta \sigma_{\alpha\beta}\bigr)\,dl\right].
\end{equation}
%
[cf.\ eq.~(2.6.14)].  In calculating the change of intensity along the ray, it is
often useful to replace the proper-length parameter~$l$ by the
``optical-depth'' parameter
%
\begin{equation}
\tag{2.6.29}
\tau_\nu = \int_i^l \rho_0\,\kappa_\nu\,dl.
\end{equation}

In the integration $\kappa_\nu$ must be evaluated at the frequency
$\nu(l)$.  The optical depth $\tau_\nu$ depends on (i) the world line of
the chosen ray, (ii) the frequency of the ray in spacetime, (iii) location
along that ray.  Then the equation of transfer
world line; and (iii) location along the ray at ``infinity'' $\nu_\infty$\ldots

One can also introduce an optical depth for scattering radiation out of the
beam, $\tau_s$:
%
\begin{equation}
\tag{2.6.30}
\tau_s = \int_i^l \rho_0\,\kappa_s\,dl.
\end{equation}

In terms of optical depths, the law of radiative transfer~(2.6.27) reads
%
\begin{equation}
\tag{2.6.31}
\frac{d(I_\nu/\nu^3)}{d\tau_\nu}
= \frac{B_\nu}{\nu^3} - \frac{I_\nu}{\nu^3}
  - \frac{\kappa_s}{\kappa_\nu}
  \left.\frac{d(I_\nu/\nu^3)}{d\tau_\nu}\right|_{\text{scattering}},
\end{equation}
%
where
%
\begin{equation}
\tag{2.6.32}
\left.\frac{d(I_\nu/\nu^3)}{d\tau_\nu}\right|_{\text{scattering}}
= \frac{\kappa_s}{\kappa_\nu}\left[\frac{I_\nu}{\nu^3} - \frac{1}{\kappa_s}
  \int \frac{dk_s}{d\Omega\,d\nu'}\,
  (\mathbf{p}',\mathbf{p})\,I_\nu'\,d\Omega'\,d\nu'
  - \rho_0\,\kappa_s\,I_\nu \right].
\end{equation}

%%% ---------------------------------------------------------------
%%% Optically thin / optically thick
%%% ---------------------------------------------------------------

One says that the medium is \emph{optically thin} if everywhere along
the ray $\tau_\nu \ll 1$ (or $\tau_s \ll 1$).

Consider a medium that (i) has its matter in thermodynamic equilibrium with
a spatially uniform temperature, (ii) is optically thick along a chosen ray, and
(iii) has absorption dominant over scattering, i.e.,
%
\begin{equation}
\tag{2.6.33}
\kappa_\nu \gg \kappa_s.
\end{equation}

One says that a medium is \emph{optically thick} to emission and absorption (or to
scattering) along a given ray if optical depths $\tau_\nu \gg 1$
(or $\tau_s \gg 1$) are achieved.

\medskip
\noindent
Along that ray, then the radiation emerging to infinity along the chosen ray must
have the blackbody form
%
\begin{equation}
\tag{2.6.34}
\frac{I_\nu}{\nu^3} = \frac{B_\nu}{\nu^3}
= \frac{2h\nu^3/c^3}{e^{h\nu/kT} - 1}.
\end{equation}
%
[One can prove this easily from the equation of transfer~(2.6.31).]\
Moreover, at a point in the medium where all rays are optically thick, the
radiation will have a blackbody intensity at all frequencies and in all
directions; so the energy density and pressure in the radiation will be
%
\begin{equation}
\tag{2.6.35}
\rho^{\text{rad}} = 3p^{\text{rad}} = bT^4,
\end{equation}
%
where $b$ is the universal constant
%
\begin{equation}
\tag{2.6.36}
b = \frac{8\pi^5 k^4}{15c^3 h^3}
  = 7.56 \times 10^{-15}\;\frac{\text{ergs}}{\text{cm}^3\,\text{K}^4}.
\end{equation}

Consider, alternatively, a medium that is optically thin along a chosen ray,
and has negligible scattering ($\tau_s \ll 1$) along that ray.  Then the
equation of transfer~(2.6.31) predicts for the radiation emerging to infinity
%
\begin{equation}
\tag{2.6.37}
I_\nu\,\nu^3 = \int_{\text{entire ray}} (\varepsilon_\nu/\nu^3)\,\rho_0\,dl
= (1/4\pi)\int (\varepsilon_\nu/\nu^3)_{\text{in region of strongest emission}}
  \rho_0\,dl.
\end{equation}

\begin{equation}
\tag{2.6.38}
\simeq \left(\int \rho_0\,dl\right)\!(1/4\pi)\,
(\varepsilon_\nu/\nu^3)_{\text{in region of strongest emission}}.
\end{equation}

In summary, the radiation from an optically thin source with negligible
scattering $I_\nu \propto \varepsilon_\nu$, and has an intensity proportional
to the amount of matter along the line of sight.  But radiation from an
optically thick source in thermodynamic equilibrium with negligible scattering
has the blackbody form independent of the nature of its emissivity.  Media with
non-negligible scattering will be studied in~\S5.10.

%%% ---------------------------------------------------------------
%%% Radiative transfer in the diffusion approximation
%%% ---------------------------------------------------------------

\bigskip
\noindent
\emph{Radiative transfer in the diffusion approximation}

\medskip
\noindent
Consider the interior of an optically thick medium $\tau_\nu \gg 1$ which is
in local thermodynamic equilibrium.  Suppose that the medium has a temperature
gradient and/or an acceleration, but that the characteristic length scale
%
\begin{equation}
\tag{2.6.39}
l_T \equiv \frac{T}{|\boldsymbol{\nabla}T + \dot{\mathbf{a}}\,T|}
\end{equation}
%
over which the temperature changes are long compared to the mean-free path of a
photon,
%
\begin{equation}
\tag{2.6.40}
l_T \gg 1/\kappa_\nu\,\rho_0 \equiv l_\nu.
\end{equation}
%
Then the radiation distribution will consist of a large, isotropic, blackbody
component, plus a tiny ``correction'' due to the temperature gradient
%
\begin{equation}
\tag{2.6.41}
I_\nu \simeq B_\nu + I_\nu^{(1)};\qquad
I_\nu^{(1)} \ll B_\nu.
\end{equation}

There is no net flux $F$ across any surface associated with the blackbody
component~$B_\nu$; but the ``correction'' term $I_\nu^{(1)}$ will lead to a
flux in the direction of
$\mathbf{h}\cdot(\boldsymbol{\nabla}T+\dot{\mathbf{a}}\,T)$.

(Here $\mathbf{h}$ is the projection operator of~\S2.5.)  One can calculate
the magnitude of that flux by inserting expression~(2.6.41) for~$I_\nu$ into
the equation of transfer~(2.6.22), and by then integrating over
$\cos\theta\,d\Omega\,d\nu$.  (Here $\theta$ is the angle between
$d\Omega$~and the direction of
$\mathbf{h}\cdot(\boldsymbol{\nabla}T+\dot{\mathbf{a}}\,T)$.)
The result is
%
\begin{equation}
\tag{2.6.43}
\mathbf{q} = (1/\kappa_R\rho_0)(\tfrac{1}{3}\partial T^4/\partial \mathbf{x})
  \,\mathbf{h}\cdot(\boldsymbol{\nabla}T + \dot{\mathbf{a}}\,T).
\end{equation}

Here $\mathbf{q}$ is the energy flux vector of~\S2.5, and its magnitude is the
flux of energy in the
$\mathbf{h}\cdot(\boldsymbol{\nabla}T+\dot{\mathbf{a}}\,T)$ direction
%
\begin{equation}
\tag{2.6.44}
|\mathbf{q}| = F.
\end{equation}

Also, in eq.~(2.6.43) $\bar{\kappa}$ is the ``Rosseland mean opacity'',
defined by
%
\begin{equation}
\tag{2.6.45}
\frac{1}{\bar{\kappa}}
\equiv \frac{\displaystyle\int_0^\infty
  (\kappa_\nu + \kappa_s)^{-1}\,
  (dB_\nu/dT)\,d\nu}
  {\displaystyle\int_0^\infty (dB_\nu/dT)\,d\nu}.
\end{equation}

For free-free transitions by themselves (eq.~2.6.25) the Rosseland mean
opacity is
%
\begin{equation}
\tag{2.6.46}
\kappa_{ff} = (0.645\times 10^{23}\;\text{cm}^2/\text{g})
  \left(\frac{C_f\,\bar{f}\,\bar{g}^2}{\mathcal{A}}\right)
  \left(\frac{\rho_0}{\text{g\;cm}^{-3}}\right) T^{-7/2}.
\end{equation}

See~\S2.1 for notation.  For electron scattering by itself in the nonrelativistic
case, the Rosseland mean opacity is
%
\begin{equation}
\tag{2.6.47}
\bar{\kappa}_{es} = \kappa_{es} = 0.40\;\text{cm}^2/\text{g}.
\end{equation}


%%% ===============================================================
%%% SECTION 2.7: Shock Waves
%%% ===============================================================

\subsection{Shock Waves}
\label{sec:2.7}

\emph{Basic references}\quad
Zel'dovich and Raizer~(1966)~\cite{ZeldovichRaizer1966};
Landau and Lifshitz~(1959)~\cite{LandauLifshitz1959};
Taub~(1948)~\cite{Taub1948};
Lichnerowicz~(1967, 1970, 1971)~\cite{Lichnerowicz1967};
Thorne~(1973a)~\cite{Thorne1973a}.

\bigskip
\noindent
\emph{Types of shock waves}

\medskip
\noindent
When gas in supersonic flow encounters an obstacle (e.g., the surface of a star,
or the geometric structure of the ergosphere of a black hole), a shock front
develops.  In the shock front the flow decelerates sharply from supersonic to
subsonic, and some of the kinetic energy of the flow gets converted into heat
(increase in entropy).

The structure of the shock depends on the nature of the forces which decelerate
the gas particles (atoms, ions, electrons).  If those forces are collisions between
the particles themselves (the usual case), one has an ordinary shock.  But if the
particles are decelerated without colliding---e.g., by impact onto the dipole
magnetic field of a neutron star, which swings the particles into Larmour
orbits---one has a \emph{collisionless} shock.  These notes will be confined to ordinary
shocks.  For the theory of collisionless shocks see, e.g., the end of~\S12 of Kaplan
(1966).

Shock waves in a partially ionized gas (e.g., the interstellar medium) can have
a somewhat different form than ordinary shocks.  As gas passes through the shock
front, much of its kinetic energy of supersonic flow can be converted into
excitation energy of atoms and ions, and can be quickly radiated away as ``line''
radiation.  The result is a bright glow from the shock front itself.  See, e.g.,
Pikel'ner~(1961) and~\S10 of Kaplan~(1966).  In this section we shall ignore such
energy losses to radiation---i.e., we shall demand that the gas behind the shock
have the same total energy per unit rest mass as the gas in front of the shock.

%%% ---------------------------------------------------------------
%%% The relativistic Rankine--Hugoniot equations
%%% ---------------------------------------------------------------

\bigskip
\noindent
\emph{The relativistic Rankine--Hugoniot equations}

\medskip
\noindent
Pick a particular event~$\mathscr{P}$ on a shock front.  In the neighborhood
of~$\mathscr{P}$ introduce a local Lorentz frame (``rest frame of the shock'')
in which (i) the shock is the surface $y = z = 0$; and (iii) on both sides of
the shock the fluid is moving in the $x$~direction, i.e., perpendicular to the
shock front (``normal shock''; see Figure~\ref{fig:2.7.1}).  That such a local
Lorentz frame exists in general one can prove quite easily [see, e.g.,
Taub~(1948)].

Denote the ``front'' side of the shock (side \emph{from} which the fluid moves)
by a ``1'', and denote the ``back'' side (side \emph{toward} which the fluid
moves) by a ``2''; see Figure~\ref{fig:2.7.1}.  Denote the velocity of the
fluid, as measured in the rest frame of the shock, by
%
\begin{subequations}
\begin{align}
\tag{2.7.1a}
v_1 &= (dx/dt)_1 = \text{ordinary velocity on front side},\\[4pt]
\tag{2.7.1b}
v_2 &= (dx/dt)_2 = \text{ordinary velocity on back side}.
\end{align}
\end{subequations}
%
(Note that these ``4-velocities'' are scalars, not vectors.)

In the rest frame of the shock the law of baryon conservation is equivalent
to continuity of the \emph{baryon flux}:
%
\begin{equation}
\tag{2.7.2a}
j \equiv n_1\,u_1 = n_2\,u_2,
\end{equation}
%
and the relations $n = 1/V_1$, $\mu_1 = p_1$,
%
\begin{subequations}
\begin{align}
\tag{2.7.1c}
u_1 &= v_1\,\gamma_1 = v_1/(1 - v_1^2)^{1/2}
  = \text{``4-velocity'' on front side},\\[4pt]
\tag{2.7.1d}
u_2 &= v_2\,\gamma_2 = v_2/(1 - v_2^2)^{1/2}
  = \text{``4-velocity'' on back side}.
\end{align}
\end{subequations}

\noindent
Similarly, energy and momentum conservation are equivalent to continuity of
energy flux and momentum flux:
%
\begin{subequations}
\begin{align}
\tag{2.7.2b}
\bigl[(\rho + p)\,v\,u\bigr]_1
  &= \bigl[(\rho + p)\,v\,u\bigr]_2,\\[4pt]
\tag{2.7.2c}
\bigl[(\rho + p)\,u^2 + p\bigr]_1
  &= \bigl[(\rho + p)\,u^2 + p\bigr]_2.
\end{align}
\end{subequations}

%%% ---------------------------------------------------------------
%%% Figure 2.7.1
%%% ---------------------------------------------------------------

\begin{figure}[htbp]
\centering
% \includegraphics[width=\columnwidth]{fig_2_7_1}
\fbox{\parbox{0.9\columnwidth}{\centering [Figure 2.7.1 placeholder]\\[6pt]
Diagram of a shock wave viewed in the local-Lorentz rest frame of the
shock front at an event~$\mathscr{P}$.  The front (``1'' side) is on the left,
and the back (``2'' side) is on the right.  Arrows indicate gas flow with
velocity~$v_1$ entering from the left and velocity~$v_2$ exiting to the right.
The $x$-axis is perpendicular to the shock surface, which lies in the
$y$-$z$ plane.  Hatching indicates the compressed gas behind the shock.}}
\caption{\textit{Figure~2.7.1.}  A shock wave viewed in the local-Lorentz
rest frame of the shock front at an event~$\mathscr{P}$.}
\label{fig:2.7.1}
\end{figure}

These junction conditions are more easily understood by writing them in a
form analogous to the Rankine--Hugoniot equations of Newtonian theory:
First take the law of baryon conservation~(2.7.2a) and turn it into equations
for the fluid 4-velocities
%
\begin{equation}
\tag{2.7.3a}
u_1 = jV_1, \qquad u_2 = jV_2.
\end{equation}
%
(See~\S2.5 for notation.)  Then take the law of momentum conservation~(2.7.2a)
and rewrite it in terms of $\mu$, $V_1$, and $p$ using the law of baryon
conservation~(2.7.2a) to obtain
%
\begin{subequations}
\begin{align}
\tag{2.7.3b}
j^2 &= -\frac{p_2 - p_1}{\mu_2\,V_2 - \mu_1\,V_1},\\[6pt]
\end{align}
\end{subequations}
%
and the relations
$j^2 = (p_2-p_1)/(\mu_2 V_2 - \mu_1 V_1)$.

Finally, use the law of baryon conservation~(2.7.2b) to rewrite the energy
equation~(2.7.3b) by $(\mu_1 V_1 + \mu_2 V_2)$, and combine with
$j = u_1/V_1 = u_2/V_2$; divide equation~(2.7.3b) by the square of the energy
equation $(\mu_2\gamma_2)^2 - (\mu_1\gamma_1)^2 = (p_1 - p_2)(\mu_1 V_1 + \mu_2 V_2)$;
then subtract this from [eq.~(2.7.2a)] to obtain
%
\begin{equation}
\tag{2.7.3c}
\mu_2^2 - \mu_1^2 = (p_2 - p_1)(\mu_1\,V_1 + \mu_2\,V_2).
\end{equation}

\emph{Equations~(2.7.3) are Taub's~(1948) junction conditions for shock waves.}
Their Newtonian limits are the standard Rankine--Hugoniot equations:
%
\begin{subequations}
\begin{align}
\tag{2.7.4a}
v_1 &= j_0\,V_{01},\qquad v_2 = j_0\,V_{02}\\[4pt]
\tag{2.7.4b}
j_0 &\equiv \frac{p_2 - p_1}{V_{01} - V_{02}}\quad
\text{(``mass flux'' of Newtonian theory)};\\[4pt]
\tag{2.7.4c}
w_2 - w_1 &= \tfrac{1}{2}(p_2 - p_1)(V_{01} + V_{02}).
\end{align}
\end{subequations}
%
The form of the Newtonian junction conditions~(2.7.4) motivates one to use $p$
and $V_0$ as one's independent thermodynamic variables when analyzing Newtonian
shocks; similarly, the form of the relativistic junction conditions~(2.7.3)
motivates one to use $p$ and~$\mu V$ as one's independent variables for
relativistic shocks.

\subsubsection*{The Rankine-Hugoniot curve}

Consider a family of shocks, each with the same thermodynamic state on the
front face (same $\mu_1$, $V_1$, $p_1$, etc.), but with different states on the
back face (different $\mu_2$, $V_2$, $p_2$, etc.).  This family of shocks is
a one-parameter family.  Thus, if one plots all back-face states
$(\mu_2 V_2,\,p_2)$ in the $\mu V$--$p$ plane, they lie on a single
curve---the ``Rankine-Hugoniot curve''---passing through the point
$(\mu_1 V_1,\,p_1)$.  (In the relativistic case this curve is also called the
``Taub adiabat'' or ``Poisson adiabat''.)

One can also plot, in the $\mu V$--$p$ plane, the Poisson adiabat (curve of
constant entropy) passing through $(\mu_1 V_1,\,p_1)$.  These two curves
typically have the relative shapes and locations shown in Figure~2.7.2.  In
particular, from the Rankine-Hugoniot equations one can derive the following
general properties of the Rankine-Hugoniot
curve:\footnote{See \citet{Thorne1973a} for derivation.  (i)~The
  Rankine-Hugoniot curve is tangent to the Poisson adiabat, and has the same
  second derivative at point~``1''.  The derivation requires the assumption
  $[\partial^2(\mu V)/\partial p^2]_{s} > 0$; a condition which is satisfied
  by all materials of astrophysical or laboratory interest; see
  \citet{LandauLifshitz1959}.  Some, but not all, of the listed properties
  remain true without this condition; see \citet{ZeldovichRaizer1966} for
  nonrelativistic details, and \citet{Thorne1973a} for relativistic details.}

\begin{enumerate}
\item[(i)] The Rankine-Hugoniot curve is tangent to the Poisson adiabat, and
  has the same second derivative at point~``1''; i.e., for weak shocks the
  increase in entropy is third-order in the pressure jump:
  %
  \begin{equation}
    s_2 - s_1 = \left[\frac{1}{12\mu T}\,
      \frac{\partial^2(\mu V)}{\partial p^2}\bigg|_{s,\,1}\right]
      (p_2 - p_1)^3 + O\!\left[(p_2-p_1)^4\right].
    \tag{2.7.5}
  \end{equation}

\item[(ii)] As the gas passes from the front of the shock to the back, its
  entropy, pressure, and chemical potential increase
  %
  \begin{equation}
    s_2 > s_1\,,\qquad p_2 > p_1\,,\qquad \mu_2 > \mu_1\,;
    \tag{2.7.6a}
  \end{equation}
  %
  while its specific volume and the product $\mu V$ decrease
  %
  \begin{equation}
    V_2 < V_1\,,\qquad \mu_2 V_2 < \mu_1 V_1\,.
    \tag{2.7.6b}
  \end{equation}
  %
  (Hence, only the ``upper branch'' of the Rankine-Hugoniot curve,
  Figure~2.7.2, is physically relevant.)

\item[(iii)] The flow on the front side is always supersonic; on the back side
  is always subsonic
  %
  \begin{equation}
    u_1/u_{S1} > 1\,,\qquad u_2/u_{S2} < 1\,;
    \tag{2.7.7}
  \end{equation}
  %
  \[
    u_S = c_S/(1 - c_S^2)^{1/2}\,.
  \]

\item[(iv)] As one moves up the Rankine-Hugoniot curve, away from point
  ``1''---i.e., as one studies a sequence of ever ``stronger''
  shocks---the following quantities increase monotonically:
  %
  \begin{equation}
  \begin{gathered}
    \text{the baryon flux across the shock, } j\,; \\
    \text{the jump in entropy across the shock, } s_2 - s_1\,; \\
    \text{the ``relativistic'' Mach number on the front side of the shock, }
      M_1 = u_1/u_{S1}\,.
  \end{gathered}
    \tag{2.7.8}
  \end{equation}
\end{enumerate}

Notice that, once the thermodynamic state on the front face of the shock
has been specified, the shock has only one free parameter.  If one fixes the
baryon flux across the shock~$j$, or the speed on the front face~$u_1$, or
the pressure on the back face~$p_2$, or any other single parameter, then all
other properties of the shock are uniquely determined.  To calculate them, one
need merely invoke the Rankine-Hugoniot equations~(2.7.4), the laws of
thermodynamics, and the equation of state of the gas.

%%% Shocks in an ideal gas %%%

\subsubsection*{Shocks in an ideal gas}

Consider, as a special but important case, a nonrelativistic ideal gas with
constant adiabatic index $\Gamma = -(\partial\ln p/\partial\ln V_0)_s$, so
that
%
\begin{equation}
  pV_0 = kT/\bar{m}\,.
  \tag{2.7.9}
\end{equation}

By a straightforward calculation (\S85 of \citealt{LandauLifshitz1959}), one
can derive the following relationships between various quantities along the
Rankine-Hugoniot curve:
%
\begin{equation}
\begin{aligned}
  \frac{V_{02}}{V_{01}} &= \frac{(\Gamma+1)\,p_1 + (\Gamma-1)\,p_2}
                               {(\Gamma-1)\,p_1 + (\Gamma+1)\,p_2}
  \;\to\; \frac{\Gamma-1}{\Gamma+1}
  &&\text{for strong shock,} \\[6pt]
  %
  \frac{T_2}{T_1} &= \frac{p_2}{p_1}\,
    \frac{(\Gamma-1)\,p_2 + (\Gamma+1)\,p_1}
         {(\Gamma+1)\,p_2 + (\Gamma-1)\,p_1}
  \;\to\; \frac{\Gamma-1}{\Gamma+1}\,\frac{p_2}{p_1}
  &&\text{for strong shock,} \\[6pt]
  %
  \hat{\jmath}_0^{\,2} &=
    \frac{(\Gamma-1)\,p_1 + (\Gamma+1)\,p_2}{2V_{01}}
  \;\to\; \frac{\Gamma+1}{2}\,\frac{p_2}{V_{01}}
  &&\text{for strong shock,} \\[6pt]
  %
  f^2 &= \tfrac{1}{2}\,V_{01}
    \bigl[(\Gamma+1)\,p_1 + (\Gamma-1)\,p_2\bigr]
  \;\to\; \tfrac{1}{2}\,(\Gamma-1)\,V_{01}\,p_2
  &&\text{for strong shock,} \\[6pt]
  %
  \hat{\jmath}_0^{\,2} &=
    \frac{V_{01}\bigl[(\Gamma+1)\,p_1 + (\Gamma-1)\,p_2\bigr]^2}
         {2(\Gamma+1)}
  \;\to\; \frac{(\Gamma-1)^2}{2(\Gamma+1)}\,V_{01}\,p_2^{\,2}
  &&\text{for strong shock.}
\end{aligned}
  \tag{2.7.10}
\end{equation}
%
Here ``strong shock'' means ``in the limit $p_2/p_1 \gg 1$''.

\begin{figure}[htbp]
\centering
\fbox{\parbox{0.9\columnwidth}{\centering [Figure 2.7.2 placeholder]\\
The Poisson adiabat $\mathscr{P}\,\mathscr{P}'$ and the Rankine-Hugoniot
curve $\mathscr{R}\,\mathscr{R}'$ plotted in the $\mu V$--$p$ plane.}}
\caption{The Poisson adiabat $\mathscr{P}\,\mathscr{P}'$ and the
  Rankine-Hugoniot curve $\mathscr{R}\,\mathscr{R}'$ plotted in the
  $\mu V$--$p$ plane.}
\label{fig:2.7.2}
\end{figure}


%%%%%%%%%%%%%%%%%%%%%%%%%%%%%%%%%%%%%%%%%%%%%%%%%%%%%%%%%%%%%%%%
%%% SECTION 2.8  TURBULENCE
%%%%%%%%%%%%%%%%%%%%%%%%%%%%%%%%%%%%%%%%%%%%%%%%%%%%%%%%%%%%%%%%

\subsection{Turbulence}
\label{sec:2.8}

In the accretion of gas onto a black hole, turbulence probably plays an
important role.  (See \S\S4 and~5.)  Unfortunately, the theory of turbulence
is in a very uncertain state.  Little is known with confidence about the
astrophysical circumstances under which turbulence should develop, or about
the strength of the turbulence; see, e.g., Chapter~3 of
\citet{LandauLifshitz1959}; also \citet{Pikelner1961}.


%%%%%%%%%%%%%%%%%%%%%%%%%%%%%%%%%%%%%%%%%%%%%%%%%%%%%%%%%%%%%%%%
%%% SECTION 2.9  RECONNECTION OF MAGNETIC FIELD LINES
%%%%%%%%%%%%%%%%%%%%%%%%%%%%%%%%%%%%%%%%%%%%%%%%%%%%%%%%%%%%%%%%

\subsection{Reconnection of Magnetic Field Lines}
\label{sec:2.9}

\emph{Basic references}\quad \S5.3 of \citet{Cowling1965};
\citet{Sonnerup1970}; \citet{Yeh1970};
\citet{VainshteinZeldovich1972}.

\subsubsection*{Basic ideas}

In flat spacetime consider a plasma that is macroscopically neutral, that has a
high electrical conductivity~$\sigma$, that is sufficiently dilute for one to
ignore its dielectric properties: $\varepsilon = \mu = 1$.  Let the plasma be
endowed with a large-scale magnetic field, and examine the evolution of that
field in a Lorentz frame where the plasma has low velocity
$|\mathbf{v}| \ll c$.  Maxwell's equations for the electromagnetic field then
read
%
\begin{align}
  \nabla \times \mathbf{E} + (1/c)(\partial\mathbf{B}/\partial t) &= 0\,,
  \tag{2.9.1} \\[4pt]
  \nabla \times \mathbf{B} - (1/c)(\partial\mathbf{E}/\partial t)
    &= 4\pi\mathbf{J}/c\,.
  \tag{2.9.2}
\end{align}

In the local rest frame of the plasma the current is proportional to the
electric field, $\mathbf{J} = \sigma\mathbf{E}$.  When transformed to the
Lorentz frame where the plasma moves with velocity~$\mathbf{v}$, this
equation says
%
\begin{equation}
  \mathbf{J} = \sigma\bigl[\mathbf{E}
    + (\mathbf{v}/c) \times \mathbf{B}\bigr].
  \tag{2.9.3}
\end{equation}

A straightforward calculation from the above equations leads to the following
law for the rate of change of magnetic field along the world lines of the
plasma:
%
\begin{equation}
  \frac{d\mathbf{B}}{d\tau}
  = (\omega + \sigma - \tfrac{2}{3}\theta\,\mathbf{1}) \cdot \mathbf{B}
    - \frac{1}{4\pi\sigma}\left(
      \frac{\partial^2\mathbf{B}}{\partial t^2}
      - c^2\,\nabla^2\mathbf{B}\right).
  \tag{2.9.4}
\end{equation}
%
Here $\mathbf{1}$ is the unit 3-tensor; $\omega$, $\sigma$, and~$\theta$ are
the rotation, shear, and expansion of the plasma
[cf.\ eqs.~(2.5.17) and~(2.5.18)]; and
%
\[
  d/d\tau = \partial/\partial t + \mathbf{v}\cdot\nabla\,.
\]

Because of the very high electrical conductivity, one can ignore the
``wave-equation'' part of eq.~(2.9.4) almost everywhere in the plasma;
and the magnetic field evolves in a ``frozen-in'' manner:
%
\begin{equation}
  d\mathbf{B}/d\tau
  = (\omega + \sigma - \tfrac{2}{3}\theta\,\mathbf{1}) \cdot \mathbf{B}\,.
  \tag{2.9.5}
\end{equation}

[Cf.\ Eq.~(2.5.20) and associated discussion.]  However, in regions of plasma
where the field has strong gradients, the ``wave-equation'' part must come
into play.

Consider, as the case of greatest importance, a magnetic field that is chaotic
with an ordered structure on some scale~$l_c$ (subscript $c$ for ``cell
size'').  At the interface between adjacent cells, the magnetic field must
reverse sign (``neutral point'' or ``neutral sheet''), and its gradients will
be very large.  Hence, at the interface, ``wave-equation'' behavior can come
into play destroying the frozen-in behavior of the field.  The result is a
reconnection of magnetic field lines, as

\begin{figure}[htbp]
\centering
\fbox{\parbox{0.9\columnwidth}{\centering [Figure 2.9.1 placeholder]\\
(a)~Top view and (b)~perspective view of reconnecting magnetic field lines
in a plasma.  The reconnection region is stippled; the magnetic field lines
are drawn solidly and the flow lines of the plasma are dotted.  The innermost
field line is in the process of reconnecting into the dashed form.
The perspective view~(b) shows the closed curve~$\mathscr{E}$ around which
one integrates to derive eq.~(2.9.6) for the EMF in the connecting region.}}
\caption{The region in a plasma where magnetic field lines are
  reconnecting (schematic picture).  In the top view, (a), the
  reconnection region is stippled; the magnetic field lines are drawn
  solidly; and the flow lines of the plasma are dotted.  The innermost
  field line is in the process of reconnecting into the dashed form.
  The perspective view, (b), shows the closed curve~$\mathscr{E}$
  around which one integrates to derive eq.~(2.9.6) for the EMF in the
  connecting region.}
\label{fig:2.9.1}
\end{figure}

depicted in Figure~2.9.1.  This reconnection gradually converts the two adjacent
cells into a single cell, reducing the overall strength of the magnetic field.

In situations of interest (\S\S4.4 and 5.2), the decrease of field strength due
to reconnection will be counterbalanced by an increase in strength due to
compression or shear of the plasma.

The reconnection process creates very strong electric fields.  In particular,
a simple application of Faraday's law
%
\begin{equation}
\oint_{\mathscr{C}} \mathbf{E}\cdot d\mathbf{l}
  = -\frac{d}{d\tau}\!\int \mathbf{B}\cdot d\mathbf{A},
\end{equation}
%
with the curve~$\mathscr{C}$ as shown in Figure~2.9.1b, shows that the EMF built up
in the reconnecting region is
%
\begin{equation}
\Delta\phi_e
  = (\text{rate of reconnection of magnetic flux})
  \equiv d\Psi/dt.
\tag{2.9.6}
\end{equation}

In situations of interest to us [\S5.12; Lynden-Bell~(1969)\cite{LyndenBell1969}],
these EMF's can be far greater than $10^{10}$~volts.  If the density of the plasma is
sufficiently low to permit long mean-free paths for charged particles, and if the
reconnecting region is sufficiently straight, then these EMF's can accelerate
particles to ultra-relativistic energies; and the particles can then radiate intense
synchrotron radiation.  In this manner regions of reconnecting field lines may be
sources of flares.


% === SECTION 3: Origin of Stellar Black Holes ===
% =============================================================================
% Section 3: The Origin of Stellar Black Holes
% Merged from: agent_11.tex (S3, Table 3.1), agent_12.tex (S3 end)
% =============================================================================

\section{The Origin of Stellar Black Holes}
\label{sec:3}

\textit{Basic references}\quad
\S13.13 of ZN\cite{ZN1971};
Peebles~(1972)\cite{Peebles1972}.

\bigskip\noindent
\textit{Basic ideas and issues}

\medskip\noindent
How much of the mass of the Galaxy is in the form of black holes?  What is the
spectrum of black hole masses?  How many black holes are created per year by
stellar collapse in our Galaxy?  To these questions one has only the vaguest of
answers in 1972.

\textit{If\/} the spectrum of stellar masses at birth is the same elsewhere as in the solar
neighborhood, and always has been the same, then roughly half of the mass of the
Galaxy has been cycled through stars of $M > 2M_\odot$, and such stars are being
born at a rate of about 0.2~per year.  \textit{If\/} the maximum mass of a neutron star is
$2M_\odot$, and \textit{if\/} stars of $M > 2M_\odot$ lose a negligible amount of mass during their
evolutions and deaths, then all of the matter which goes into such stars eventually
winds up in black holes.  Thus, 5~per cent of the mass of the Galaxy might be in
black holes, and new black holes might be forming at a rate of ${\sim}\,0.2$~per
year---\textit{might}.  But very probably not, because the if's which go into this result
are probably not satisfied.

In particular, both observation and theory suggest that mass loss is very
significant during the late stages of stellar evolution and during the death throes
of massive stars.  Mass loss might be so important that black holes are rare
beasts, indeed!

Let us focus attention on mass loss during the death throes, as predicted by
the best numerical calculations to date.  But in doing so, let us keep in mind the
very primitive state of the modern computations---for example, their typical
neglect of the effects of rotation.

Modern computations conclude that stars whose masses exceed the
Chandrasekhar limit, $M\gtrsim 1.4\,M_\odot$, at the endpoint of their evolution should
explode as supernovae.  A supernova explosion can leave behind two remnants: an
expanding gas cloud (the outer part of the original star), and a remnant ``star''.
The remnant ``star'' will be a neutron star if its mass, $M_{\text{remnant}}$, is less
than ${\sim}\,2M_\odot$; but will be a black hole if its mass is greater than
${\sim}\,2M_\odot$.\footnote{In these calculations the initial core and final remnant
have the same mass: $M_\odot$ in one case, and $3M_\odot$ in another.
But since we have assumed that $M_{\text{star}} = M_{\text{core}}$, and since the models
give no mass loss, we must take $2.8\times M_\odot$ and $3\times 32 = M_\odot$ as the
masses of the ``true'' remnants.}

We summarize in Table~\ref{tab:3.1} the results of various calculations of the
supernova process.

%% ============================================================
%%  Table 3.1  (page 387 -- printed upside-down in the original)
%% ============================================================

\begin{table*}[htbp]
\centering
\caption{Summary of Numerical Computations of Supernova Explosions}
\label{tab:3.1}
\renewcommand{\arraystretch}{1.4}
\small
\begin{tabular}{l|c c|c c|p{4.5cm}}
\hline\hline
 & \multicolumn{2}{c|}{Initial Conditions}
 & \multicolumn{2}{c|}{Results}
 & \\[2pt]
\cline{2-5}\\[-8pt]
Authors
 & $\dfrac{M_{\text{core}}}{M_\odot}$
 & $\dfrac{M_{\text{star}}{}^\dagger}{M_\odot}$
 & $\dfrac{M_{\text{remnant}}}{M_\odot}$
 & \parbox{2.2cm}{\centering Envelope\\Kinetic Energy\\($10^{51}$~ergs)}
 & \centering Main Factor Regulating\\the Explosion\arraybackslash \\[4pt]
\hline
Fraley (1968)\cite{Fraley1968}
 & $\gtrsim 40$
 & $\gtrsim 100$
 & 0
 & 3
 & Oxygen detonation in the core at the end of quasistatic evolution \\[4pt]
\hline
Colgate and White (1966)\cite{ColgateWhite1966}
 & 10
 & 30
 & 1.8
 & 1.9
 & Absorption of neutrinos in the envelope \\[4pt]
\hline
Arnett (1966)\cite{Arnett1966}
 & 1.5
 & $\gtrsim 3.5$
 & 0.87
 & 0.1
 & \multirow{2}{4.5cm}{Carbon detonation in the core at the end of quasistatic evolution; degenerate core} \\
 & 1.4
 & 4
 & 0.1
 & 0.1
 & \\[4pt]
\hline
Ivanova, Imshennik, and
Nadezhin (1969)\cite{IvanovaImshennikNadezhin1969}
 & 2
 & $\gtrsim 4$
 & 2.46, 0.95
 & 196
 & \multirow{2}{4.5cm}{Absorption of neutrinos in the envelope} \\
 & 4
 & $\gtrsim 8$
 & 0.56, 0.66
 & $4+$
 & \\[4pt]
\hline
Nadezhin (1969)\cite{Nadezhin1969}
 & 10
 & 30
 & 9.75
 & 0.025
 & Oxygen detonation during the envelope during collapse \\[4pt]
\hline
Hansen and Barkat (1969)\cite{HansenBarkat1969}
 & 1.4
 & 4
 & $\lesssim 0$
 & 0.07
 & Carbon detonation during collapse \\[4pt]
\hline
Arnett (1969)\cite{Arnett1969}
 & 33
 & $\lesssim 100$
 & 0
 & 0
 & \parbox{4.5cm}{No mass loss; $\nu$-neutrinos. The core is opaque to absorption of neutrinos.} \\[4pt]
\hline
Arnett (1967)\cite{Arnett1967}
 & 8
 & 30
 & 9.75
 & 0
 & \parbox{4.5cm}{No mass loss; $\nu$-neutrinos. The core is opaque to absorption of neutrinos.} \\[4pt]
\hline\hline
\end{tabular}

\medskip
\begin{flushleft}
\footnotesize
$\dagger$~In recent calculations by Bruenn~(1972)\cite{Bruenn1972} the thermal
explosion of a star of $M_\odot$ produced a remnant.\\[2pt]
$\ddagger$~This does not change the conclusions given in the text about the
formation of black holes.
\end{flushleft}
\end{table*}

%% ============================================================
%%  Continuation -- discussion of Table 3.1
%% ============================================================

In Table~3.1 $M_{\text{core}}$ is formally the mass of the entire star used in the
numerical calculations.  However, all the calculations assumed a homogeneous
initial model (typically polytropic or isothermal), whereas the theory of stellar
evolution predicts a highly inhomogeneous structure for presupernova stars.  In
particular, stellar evolution predicts a central core which is more or less
homogeneous, surrounded by a huge diffuse envelope of iron, hydrogen, and other
elements.  Thus, we must imagine the initial stars of the supernova calculations to
be the central, homogeneous core; and we must regard those calculations as
neglecting the envelope.  Such neglect is reasonable when one studies supernova
dynamics, because the core is collapsing and reexploding.  The envelope will
generally be thrown off, with little expenditure of work, by the reexploding parts
of the core.  The mass of the envelope might be twice that of the core---or
perhaps only the same as the core, or perhaps even less.  The theory of stellar
evolution is not definite on this point.  Hence the question marks in column~2 of
Table~3.1.

It is clear, from the diverse results shown in Table~3.1, that the theory of
supernova explosions is far from perfect in 1972.  Nevertheless, one can form the
following very tentative conclusions.  (i)~For stars which approach the ends of
their quasistatic evolution with masses less than ${\sim}\,12$ to $30\,M_\odot$, the
supernova explosion may produce a neutron star.  (ii)~For stars with masses
greater than ${\sim}\,12$ to $30\,M_\odot$, the explosion may produce a black hole.  If
this tentative conclusion is correct, then no more than ${\sim}\,1$~per cent of the
mass of the Galaxy should be in the form of black holes today; and new black
holes should be created at a rate no greater than ${\sim}\,0.01$~per year.


% === SECTION 4: Black Holes in the Interstellar Medium ===
% =============================================================================
% Section 4: Black Holes in The Interstellar Medium
% Merged from agents 12--17
% Subsections 4.1 through 4.9
% =============================================================================

\section{Black Holes in the Interstellar Medium}
\label{sec:4}

%%% ====================================================================
%%% 4.1  Accretion of Noninteracting Particles onto a Nonmoving Black Hole
%%% ====================================================================

\subsection{Accretion of Noninteracting Particles onto a Nonmoving Black Hole}
\label{sec:4.1}

Consider a black hole at rest in the interstellar medium; and temporarily treat
the interstellar gas as though it were made up of noninteracting particles
[collisions neglected; mean free paths large compared to the region over which
the hole's gravity makes itself felt; $l_\text{fp} \gg 2GM/c^2$; see below]. In the case of a
Schwarzschild hole, all particles with angular momentum per unit mass
$\ell < 2r_g c$  eventually get captured by the hole (see Box~25.6 of MTW or Figure~12
of ZN). Here $r_g = 2GM/c^2$ is the gravitational radius of hole and $M$ is the hole's
mass. If the particle speeds far from the hole are $v_\infty$, then this condition for
capture corresponds to an impact parameter
%
\begin{equation}
b = \ell / v_\infty = 2r_g c/v_\infty\,.
\tag{4.1.1}
\end{equation}
%
Consequently, the ``capture cross section'' of the hole is
%
\begin{equation}
\sigma_\text{capture} = \pi b^2 = 4\pi r_g^2 (c/v_\infty)^2\,.
\tag{4.1.2}
\end{equation}
%
For a Kerr hole the capture cross section is of this same order of magnitude. The
rest mass per unit time crossing inward through a sphere of $r \gg r_\text{capture}$, with
particles directed into the capture region (solid angle $\Delta\Omega = \sigma/r^2$), is
%
\begin{equation}
\dot{M}_0 = \left(\frac{\rho_\text{cap.}}{4\pi}\right) 4\pi r^2 \Delta\Omega = 4\pi r_g^2 \rho_\infty c^2 / v_\infty\,.
\tag{4.1.3}
\end{equation}
%
This is the rate at which the hole accretes rest mass. Rewritten in typical astro\-nomical units, this accretion rate is
%
\begin{equation}
\dot{M}_0 = 10^{11} \left(\frac{\rho_\infty}{10^{-24}~\text{g~cm}^{-3}}\right) \left(\frac{M}{M_\odot}\right)^2 \left(\frac{v_\infty}{10~\text{km~sec}^{-1}}\right)^{-1} \text{``H~II''}
\tag{4.1.3$'$}
\end{equation}
%
\begin{equation}
\frac{d(M_0/M_\odot)}{d(t/10^{10}~\text{yrs})} = 10^{-13} \left(\frac{\rho_\infty}{10^{-24}~\text{g~cm}^{-3}}\right) \left(\frac{M}{M_\odot}\right)^2 \left(\frac{v_\infty}{10~\text{km~sec}^{-1}}\right)^{-1}\,.
\notag
\end{equation}

[The typical densities and speeds of interstellar gas particles in ionized, ``H~II''
regions are $10^{-24}$~g~cm$^{-3}$ and 10~km~sec$^{-1}$; see e.g.\ Kaplan and Pikel'ner (1970)].
Even if 100 per cent of the inflowing rest mass were somehow converted into
outgoing radiation, the total luminosity would be only
%
\begin{equation}
L_\text{max} = \left(10^{24}~\frac{\text{erg}}{\text{sec}}\right) \left(\frac{\rho_\infty}{10^{-24}~\text{g~cm}^{-3}}\right) \left(\frac{M}{M_\odot}\right)^2 \left(\frac{v_\infty}{10~\text{km~sec}^{-1}}\right)^{-1}
\tag{4.1.4}
\end{equation}
%
---a value much too small to be astronomically interesting. Therefore it is fortunate
that the electrons, ions, and magnetic fields of the interstellar gas interact strongly
enough to make the accretion process obey fluid-dynamic laws rather than the
laws of noninteracting particles (see \S4.6, below).

%%% ====================================================================
%%% 4.2  Adiabatic, Hydrodynamic Accretion onto a Nonmoving Black Hole
%%% ====================================================================

\subsection{Adiabatic, Hydrodynamic Accretion onto a Nonmoving Black Hole}
\label{sec:4.2}

Switch, then, from a noninteracting description of accretion to a hydrodynamic
description. Assume, as above, that the black hole is at rest with respect to the gas
and a Schwarzschild hole, so that the accretion is spherically symmetric. In the
hydrodynamic case the accretion rate $\dot{M}_0$ and all other basic characteristics of the
flow are governed by the gravitational field at distances much greater than the
gravitational radius. (This is because the flow at small radii is supersonic and
therefore cannot influence conditions at large radii.) Thus, one can calculate
the mass flow and other quantities at large radii accurately using the Newtonian
theory of gravitation. Phenomena close to the gravitational radius will be treated
later. The motion of the gas will be assumed adiabatic. Deviations from adiabatic
flow due to radiative losses and radiative transport between gas elements can be
taken into account later by suitably modifying the adiabatic index $\Gamma$.
The form of the flow is governed by two fundamental equations: conservation
of rest mass (``continuity equation''), which we write in the form\footnote{Material for this section is drawn largely from Chapter~13 of ZN, and from Shvartsman and Novikov (1971).}
%
\begin{equation}
4\pi r^2 \rho u = \dot{M}_0 = \text{constant, independent of}~r,
\tag{4.2.1}
\end{equation}
%
and the Euler equation
%
\begin{equation}
u\frac{du}{dr} = -\frac{1}{\rho}\frac{dp}{dr} - \frac{GM}{r^2}\,.
\tag{4.2.2}
\end{equation}
%
[Cf.\ eqs.\ (2.5.23) and (2.5.59).]

Here $\dot{M}_0$ is the total rate of accretion of rest mass, $u$ is the radial
velocity, and $M$ is the mass of the hole.  We assume the gas has constant
adiabatic index~$\Gamma$, so that during the accretion the pressure and
density are related by the adiabatic law $p = K\rho^\Gamma$, and the speed of
sound is given by $a = (\Gamma p/\rho)^{1/2}$.  Let the constants $K$ and $\Gamma$
be given.  Our task is to determine the accretion rate~$\dot{M}_0$ and to
compute as well the distributions of density $\rho(r)$ and velocity $u(r)$
in the flow.\footnote{Notice that our notation differs from that in \S2.
Throughout \S4 we denote (for ease of eyesight)

(speed of sound) $= a$, not $c_s$

(density of rest mass) $= \rho$, not $\rho_0$

No confusion is likely, since nowhere in \S4 shall we deal with
4-accelerations or with total mass-energy; and
$\rho = \rho_0(1+\Pi)$, and because $\Pi \ll 1$ in all regions of the accreting gas.
We shall also retain factors of $G$, $c$, and $k$ (i.e.\ use cgs units)
throughout \S4.}

In place of the Euler equation~(4.2.2) we shall use the Bernoulli
equation (2.5.38) rewritten in the form
%
\begin{equation}
\tfrac{1}{2}\,u^2 + \frac{1}{\Gamma - 1}\,a^2 - \frac{GM}{r}
= \text{constant}
= \frac{1}{\Gamma - 1}\,a_\infty^2\,.
\tag{4.2.3}
\end{equation}
%
[Here the enthalpy has been written in the form
$w = \int \rho^{-1}\,dp = a^2/(\Gamma-1)$;
cf.\ eq.\ (2.5.42).]  The constant has been determined by conditions at
infinity, where the gas is at rest.

Rewrite the law of mass conservation~(4.2.1) with density expressed in terms
of sound speed
%
\begin{equation}
u = \frac{\dot{M}_0}{4\pi\rho_\infty\,r^2}
    \left(\frac{a_\infty}{a}\right)^{2/(\Gamma-1)}\!.
\tag{4.2.4}
\end{equation}

The only unknown parameter in the coupled equations~(4.2.3) and~(4.2.4) is
the accretion rate~$\dot{M}_0$.  Once $\dot{M}_0$ is known, one can readily
solve for the radial distributions of velocity, sound speed, and density.

The mass flux is determined in the following way.  Consider the system of
equations~(4.2.3) and~(4.2.4).  In the $u$--$a$ plane, the Bernoulli
equation~(4.2.3) for fixed radius~$r$ defines an ellipse; each value of $r$
corresponds to a different ellipse.  Similarly, the equation of mass
conservation~(4.2.4) for each value of $r$ defines a hyperbola of fractional
power, $u\,a^{2/(\Gamma-1)} = \text{const}$.  Now, let the accretion
rate~$\dot{M}_0$ be chosen arbitrarily.  For every value of $r$,
equations~(4.2.3) and~(4.2.4) are 2 equations with 2 unknowns, $u$ and $a$.
Solving these equations---i.e.\ finding the intersection point of the ellipse
with the hyperbola---gives $u$ and $a$ for that particular radius.  In other
words, in the $u$, $a$ plane the curve~$u(a)$ is determined parametrically by
the intersections of corresponding ellipses and hyperbolae.  It is clear (see
Figure~\ref{fig:4.2.1}) that for every ellipse-hyperbola pair, there are
either two intersection points, or one point of tangency, or no intersection
at all.  If, for a particular chosen~$\dot{M}_0$, there exists any $r$ at
which the two equations are incompatible for the~$\dot{M}_0$, such flow
cannot occur.  Rejecting this case, we have 2 remaining possibilities, for a
given~$\dot{M}_0$ (see Figure~\ref{fig:4.2.2}):

\begin{enumerate}
\item The ellipse and hyperbola intersect twice for every value of~$r$
  (Figure~\ref{fig:4.2.2}a).\\
  In this case we have 2 separate curves, corresponding to 2 families of
  intersection points.

\item The ellipse and hyperbola meet tangentially
  (Figure~\ref{fig:4.2.2}b) at one value of $r$ (call it~$r_s$), at which
  they cross each other; the two families of intersection points $u(a)$
  cross each other at $r = r_s$.  We shall show below that this
  crossing point necessarily lies on the ``bisectrix'' of the graph,
  $u = a$.
\end{enumerate}

In both the cases, 1 and 2, the curves $u(a)$ describe possible flows of gas
in the gravitational field.  The curves begin on the innermost ellipse,
$r = \infty$.  On the upper curve~$u(a)$, at this ellipse we have $a = 0$,
but $u \neq 0$.  These boundary conditions are not of interest for the
accretion problem,\footnote{This curve describes the ``stellar wind'' by
which some stars eject matter into interstellar space; see Chapter~13 of
ZN.} because accretion requires $u_\infty = 0$,
$a_\infty \neq 0$.  The lower curve~$u(a)$ begins at the point
$u_\infty = 0$, $a_\infty \neq 0$---

\begin{figure}[htbp]
\centering
% \includegraphics[width=\columnwidth]{fig_4_2_1}
\fbox{\parbox{0.9\columnwidth}{\centering [Figure 4.2.1 placeholder]\\[6pt]
The $u$, $a$ ``solution plane'' for solving the coupled Bernoulli
equation~(4.2.3) and law of mass conservation~(4.2.4).
Solid curves: ellipses from the Bernoulli equation at several fixed values
of~$r$ ($r = \infty$ is the innermost ellipse).
Dashed curves: hyperbolae from the continuity equation.
Heavy dots mark intersection points; the locus of these intersections
traces the flow solution $u(a)$.
The diagonal $u = a$ is the ``bisectrix''.
Regions above the bisectrix are supersonic; regions below are subsonic.}}
\caption{The $u$, $a$ ``solution plane'' for solving the coupled Bernoulli
equation~(4.2.3) and law of mass conservation~(4.2.4).}
\label{fig:4.2.1}
\end{figure}

to our desired boundary condition.  Hence we shall restrict attention to the
lower curve.

We must still decide which type of lower curve is reasonable---one that
remains always below the bisectrix (case~1), or one that crosses the
bisectrix (case~2).  In case~1 the flow is subsonic at all radii, $u < a$.
But this requires a large back-pressure at all radii to retard the
inflow---back pressure that can never be provided near the gravitational
radius of a black hole.  The gas must cross the gravitational

\begin{figure}[htbp]
\centering
% \includegraphics[width=\columnwidth]{fig_4_2_2}
\fbox{\parbox{0.9\columnwidth}{\centering [Figure 4.2.2 placeholder]\\[6pt]
Possible forms of adiabatic gas flow in a spherical, Newtonian gravitational
field.  Curves~(a) [case~1 in text] correspond to flow that remains subsonic
or supersonic everywhere.  Curves~(b) [case~2 in text] correspond to flow
for which there is a sonic point, $u = a$.}}
\caption{Possible forms of adiabatic gas flow in a spherical, Newtonian
gravitational field.  Curves~(a) [case~1 in text] correspond to flow that
remains subsonic or supersonic everywhere.  Curves~(b) [case~2 in text]
correspond to flow for which there is a sonic point, $u = a$.}
\label{fig:4.2.2}
\end{figure}

\noindent
radius with the speed of light, as measured in the \emph{proper} reference
frame of an observer there; hence, the flow is surely supersonic near the
gravitational radius.

Thus, the only acceptable solution $u(a)$ is the lower curve of case~2
(Figure~\ref{fig:4.2.2}b).  This curve begins at $u = u_\infty = 0$,
$a = a_\infty \neq 0$, and it crosses from the subsonic region into the
supersonic region at $r = r_s$.  This transition to supersonic flow is
crucial to the accretion process.  It governs the accretion rate~$\dot{M}_0$,
in the same manner as the transition to supersonic flow in the throat of a
rocket nozzle.  In both cases one and only one mass flow rate is compatible
with the required transition to supersonic flow.

Let us calculate the required accretion rate~$\dot{M}_0$.  We begin by
calculating the flow velocity at the transition radius (``sonic
radius'')~$r_s$.  To do this, we rewrite the Euler equation~(4.2.2) in the
form
%
\begin{equation}
\frac{du}{dr} = -\,\frac{a^2\,dp}{p\,dr} - \frac{GM}{r^2}\,.
\tag{4.2.5}
\end{equation}

We differentiate the law of mass conservation~(4.2.1) with respect to $r$
and put $du/dr$ from it into~(4.2.5).  As a result we obtain
%
\begin{equation}
\frac{dp}{dr}\;\frac{(u^2 - a^2)}{p}
  = -\frac{GM}{r^2} + \frac{2u^2}{r}\,,
\tag{4.2.6}
\end{equation}
%
which shows that at the sonic point, where $a = u$,
%
\begin{equation}
2u_s^2 = \frac{GM}{r_s}\,,\qquad \text{or} \quad
u_s^2 = a_s^2 = \tfrac{1}{2}\,\frac{GM}{r_s}\,.
\tag{4.2.7}
\end{equation}

By combining this result with the Bernoulli equation~(4.2.3) we obtain the
speed of sound at the sonic point in terms of the speed of sound at infinity
%
\begin{equation}
a_s = u_s = a_\infty \left(\frac{2}{5-3\Gamma}\right)^{1/2}\!;
\tag{4.2.8}
\end{equation}
%
and using equation~(4.2.7) we obtain the radius at the sonic point
%
\begin{equation}
r_s = \left(\frac{5-3\Gamma}{4}\right)\frac{GM}{a_\infty^2}\,.
\tag{4.2.9}
\end{equation}
%
(Notice that the gravitational potential $GM/r_s$ at the sonic point is of
order~$a_\infty^2$.)

Now it is straightforward to calculate the accretion rate from
equation~(4.2.4):
%
\begin{equation}
\dot{M}_0 = 4\pi\,\alpha\,G^2 M^2\,\rho_\infty\,a_\infty^{-3}
= 4\Gamma^{3/2}\,\alpha\,G^2\,M^2\,\rho_\infty\,
  (m_p / k T_\infty)^{3/2}\,.
\tag{4.2.10}
\end{equation}

Here $\alpha$ is a constant of order unity which depends on~$\Gamma$:
%
\begin{equation}
\alpha \equiv \frac{\pi}{4\Gamma^{3/2}}
  \left(\frac{2}{5-3\Gamma}\right)^{(5-3\Gamma)/[2(\Gamma-1)]}\!.
\tag{4.2.11}
\end{equation}
%
\begin{align}
\alpha &\simeq 1.5  \quad \text{for} \quad \Gamma = 1,    \notag \\
\alpha &\simeq 1.2  \quad \text{for} \quad \Gamma = 1.4,  \notag \\
\alpha &\simeq 0.3  \quad \text{for} \quad \Gamma = 5/3;
\end{align}
%
and we have used the equation of state
$p_\infty = (\rho_\infty / m_p)\,k\,T_\infty$.

This expression for the accretion rate is the main result of our calculation.

Notice that $\dot{M}_0$ depends on~$\Gamma$.  At a typical temperature of
$T_\infty \sim 10^4$\,K the interstellar gas will be partially ionized, so
when it is compressed some energy will go into further ionization.  This means
that the value of $\Gamma$ outside and near the sonic point will be somewhat
below $5/3$.  A value of $\Gamma = 1.4$ might be reasonable for use in
equations~(4.2.10) and~(4.2.11).

It is now straightforward to derive from equations~(4.2.3) and~(4.2.4) the
radial distributions of all quantities.  The general picture is as follows.
Outside the ``radius of influence''
%
\begin{equation}
r_i = 2\,GM/a_\infty^2 = r_g\,(c/a_\infty)^2
\tag{4.2.12}
\end{equation}
%
\[
\simeq 10\,r_g \quad \text{for} \quad \Gamma = 1.4,
\]
%
the pull of the hole hardly makes itself felt, so $\rho$ and $a$ are nearly
equal to their values at ``infinity''; $\rho \simeq \rho_\infty$,
$a \simeq a_\infty$; and the density and sound velocity begin to rise.
After the gas passes through the sonic point~$r_s$, it is in near free-fall
%
\begin{equation}
u = (2GM/r)^{1/2} = a_\infty\,(r_i/r)^{1/2}
\qquad \text{at} \quad r < r_s\,;
\tag{4.2.13a}
\end{equation}
%
so the law of mass conservation requires the density to increase as
%
\begin{equation}
\rho = \frac{\dot{M}_0}{4\pi r^2\,u}
\simeq 0.2\,\rho_\infty\,(r_i/r)^{3/2}
\qquad \text{at} \quad r < r_s\,.
\tag{4.2.13b}
\end{equation}

For adiabatic compression, temperature rises as
$T \propto \rho^{\,\Gamma-1}$; consequently, at $r < r_s$
%
\begin{equation}
T \simeq 0.5\,T_\infty\left(\frac{r_i}{r}\right)^{\frac{3}{2}(\Gamma-1)}
\qquad \text{if} \quad \Gamma = 1.4~\text{near sonic point.}
\tag{4.2.13c}
\end{equation}

The above solution, (4.2.10)--(4.2.13), for adiabatic hydrodynamic accretion
is due originally to \citet{Bondi1952}.

Rewrite the accretion rate~(4.2.10) for hydrodynamic flow in a form similar
to that (eq.~4.1.3) for noninteracting particles
%
\begin{equation}
\dot{M}_0 = \pi\,r_g^2\,\rho_\infty\,c\,(c/a_\infty)^3.
\tag{4.2.14}
\end{equation}

Direct comparison, and use of the approximate equality between speed of sound
$a_\infty$ and proton speeds~$v_\infty$, shows that the hydrodynamic
accretion rate is larger by a factor $(c/a_\infty)^2 \simeq 10^{10}$ than
the accretion rate for independent particles.  The physical reason for this
is clear: gas is distinguished from independent particles by the frequent
collisions of its atoms; these collisions limit the tangential velocities
during infall, but permit radial velocities to grow.

Other, useful forms for the hydrodynamic accretion rate are
%
\begin{equation}
\dot{M}_0 \simeq 10^{5}
\left(\frac{M}{M_\odot}\right)^{\!2}
\left(\frac{\rho_\infty}{10^{-24}\;\text{g\,cm}^{-3}}\right)
\left(\frac{a_\infty}{10\;\text{km\,sec}^{-1}}\right)^{\!-3}
\;\text{,}
\tag{4.2.15a}
\end{equation}
%
\begin{equation}
\dot{M}_0 \simeq \left(1\times 10^{11}\;\frac{\text{g}}{\text{sec}}\right)
\left(\frac{M}{M_\odot}\right)^{\!2}
\left(\frac{\rho_\infty}{10^{-24}\;\text{g\,cm}^{-3}}\right)
\left(\frac{T_\infty}{10^4\;\text{K}}\right)^{\!-3/2}\!.
\tag{4.2.15b}
\end{equation}

%%% ====================================================================
%%% 4.3  Thermal Bremsstrahlung from the Accreting Gas
%%% ====================================================================

\subsection{Thermal Bremsstrahlung from the Accreting Gas}
\label{sec:4.3}

The above idealized case of spherical, adiabatic, hydrodynamic accretion can be
complicated by various physical processes.  Consider, first, the effects of energy
loss due to thermal bremsstrahlung.  To calculate the bremsstrahlung most easily,
we shall assume that the energy loss has a negligible effect on the thermal energy
density of the gas, and then we shall check that this assumption is valid.

When the effects of bremsstrahlung losses are neglected, then the temperature,
density, and velocity distributions retain their adiabatic forms (4.2.13).  One can
readily check that for $T_\infty \sim 10^4$\,K and $\Gamma = 1.4$ the gas remains nonrelativistic,
$kT \ll m_p c^2$ (but just barely so) down to $r \simeq r_s$.  Consequently, the rate at which
one gram of gas at radius $r$ radiates thermal bremsstrahlung is [cf.\ eq.~(2.1.29)]
%
\begin{equation}
\varepsilon_{ff}
= \bigl(5 \times 10^{20}~\text{ergs/(g\,sec)}\bigr)\,
  (\rho / \text{g\,cm}^{-3})
  \left(\frac{T_\infty}{10^4\,\text{K}}\right)^{\!1/2}
  \left(\frac{r_i}{r}\right)^{\!3(\Gamma-1)/2}
\tag{4.3.1}
\end{equation}
%
For comparison, the rate at which adiabatic compression increases the energy of
one gram of gas is [eq.~(2.6.41)]
%
\begin{equation}
\varepsilon_{\text{ad.\ heating}}
= \frac{p}{\rho}\,\frac{d\rho}{\rho\,dt}
= -\frac{3}{2}\,\frac{u}{r}\,\frac{p_\infty}{\rho_\infty}
  \left(\frac{r_i}{r}\right)^{3\Gamma/2}\,.
\tag{4.3.2}
\end{equation}
%
Using expression (4.2.12) for $r_i$, and the thermodynamic relations for a hydro\-gen gas
%
\begin{equation}
a_\infty^2 = (\Gamma p_\infty / \rho_\infty) = \Gamma kT_\infty / m_p\,,
\tag{4.3.3}
\end{equation}
%
we can bring this heating rate into the form
%
\begin{equation}
\varepsilon_{\text{ad.\ heating}}
= \bigl(3 \times 10^4~\text{ergs}/(\text{g\,sec})\bigr)
  \left(\frac{T_\infty}{10^4\,\text{K}}\right)^{\!5/2}
  \left(\frac{M}{M_\odot}\right)^{\!-1}
  \left(\frac{r_i}{r}\right)^{3\Gamma/2}\,.
\tag{4.3.4}
\end{equation}

This heating rate greatly exceeds the free-free loss rate (4.3.1) at all radii.  Hence,
we are justified in our use of the adiabatic forms of $T(r)$, $\rho(r)$, and $u(r)$.  The total
power radiated as thermal bremsstrahlung is
%
\begin{equation}
L_{ff}
= \int_{r_g}^{\infty} \varepsilon_{ff}\,\rho\, 4\pi r^2\, dr
\sim \left(5 \times 10^{31}~\frac{\text{ergs}}{\text{sec}}\right)
  \left(\frac{M}{M_\odot}\right)^{\!3}
  \left(\frac{T_\infty}{10^4\,\text{K}}\right)^{\!-3.3}
  \left(\frac{\rho_\infty}{10^{-24}~\text{g\,cm}^{-3}}\right)^{\!2}\,,
  \quad\text{for }\Gamma = 1.4.
\tag{4.3.5}
\end{equation}

Perhaps a more realistic value for $\Gamma$ would be a little less than $5/3$ down to
$r \sim 10^3 r_g$ at which point $kT \sim m_e c^2$, and then $\Gamma = 4/3$ below that radius.  In
this case $L_{ff}$ would be increased by several orders of magnitude [cf.\ the rela%
tivistic corrections in eq.~(2.1.31)].  Even with such an increase, and even for
$M = 100\,M_\odot$, $L_{ff}$ is so small that it is of little or no observational interest.  For
this reason, we shall not attempt a more rigorous calculation of it.

%%% ====================================================================
%%% 4.4  Influence of Magnetic Fields and Synchrotron Radiation
%%% ====================================================================

\subsection{Influence of Magnetic Fields and Synchrotron Radiation}
\label{sec:4.4}

Up to now we have ignored the fact that the interstellar gas possesses a magnetic
field.  For interstellar temperatures of interest, $T_\infty \sim 10^4$\,K, the magnetic field
will be frozen into the gas (cf.\ \S\ref{sec:2.5}).  The strength of the typical intergalactic
field, $B_\infty \sim 10^{-6}$\,G, is such that its energy density and pressure are not far below
those of the gas
%
\begin{equation}
\frac{B_\infty^2}{8\pi}
\sim 4 \times 10^{-14}~\frac{\text{ergs}}{\text{cm}^3}\,.
\tag{4.4.1}
\end{equation}

At large radii the field may be small enough that one can ignore its influence on
the accretion.  In this case the standard hydrodynamical inflow will stretch each
fluid element out radially (cross-sectional area of fluid element $\propto r^2$, hence
diameter $\propto r$; volume $\propto \rho^{-1} \propto r^{3/2}$, hence radial diameter $\propto r^{-1/2}$).
This radial stretch and tangential compression must quickly convert the initial
field in the fluid element into a nearly radial field with strength
%
\begin{equation}
B \propto \text{(cross sectional area)}^{-1} \propto r^{-2}
\tag{4.4.2}
\end{equation}
%
(conservation of flux).  Hence, the magnetic energy in a given fluid element will
rise as
%
\begin{equation}
E_{\mathrm{mag}}
= \frac{B^2}{8\pi}~\text{(volume of element)}
\propto r^{-4}\,\rho^{-1}
\propto r^{-4}\,r^{3/2}
= r^{-5/2}\,.
\tag{4.4.3}
\end{equation}

Not long after the gas crosses the sonic point, this magnetic energy will become
so large as to exceed the thermal energy in the fluid element
%
\begin{equation}
E_{\mathrm{therm}}
= \frac{3\,kT}{2\,m_p}~\text{(mass of element)}
\propto r^{-(3/2)(\Gamma-1)}\,.
\tag{4.4.2a}
\end{equation}

At this point magnetic pressures must come into play, and the flow will cease
to obey the standard adiabatic hydrodynamic laws of \S\ref{sec:4.2}.

The form of the modified magnetohydrodynamic flow is not understood at
all well today.  The best guesses and calculations to date are those of Shvartsman
(1971).  They suggest a turbulent flow, with reconnection of field lines between
adjacent ``cells'' (cf.\ \S\ref{sec:2.9}), and with a rough equipartition of the gravitational
potential energy between magnetic fields, kinetic infall of gas, turbulent energy,
and thermal kinetic energy of gas particles:
%
\begin{equation}
\bigl\{\text{gravitational energy}\bigr\}
= \frac{GM}{r}
\sim 4u^2
\sim 4\rho v_{\mathrm{turb}}^2
\sim 4\,\frac{3kT}{m_p}
\sim \frac{B^2}{8\pi\rho}
\tag{4.4.2b}
\end{equation}
%
per unit mass.

The mass density corresponding to this equipartition condition,
$4\pi r^2 \rho u = \dot{M}_0$, because $\dot{M}_0$ is determined by conditions
outside and at the sonic point, where the magnetic field is not yet large enough
to strongly influence the flow, we can use equation~(4.2.15b) for $\dot{M}_0$ and write
%
\begin{equation}
\rho
\sim \left(6 \times 10^{-12}~\frac{\text{g}}{\text{cm}^3}\right)
  \left(\frac{\rho_\infty}{10^{-24}~\text{g\,cm}^{-3}}\right)
  \left(\frac{T_\infty}{10^4\,\text{K}}\right)^{\!-3/2}
  \left(\frac{r_i}{r}\right)^{\!3/2}\,.
\tag{4.4.3a}
\end{equation}

Accepting this educated guess as to the nature of the flow, we can ask what
types of radiation might be emitted.  As before~(\S\ref{sec:4.3}) bremsstrahlung is too weak
to be interesting.  However, synchrotron radiation can be rather strong.  Most of
the synchrotron radiation will come from the high-temperature, strong-field
region near the Schwarzschild radius.  There $kT \gg m_e c^2$, so the relativistic
formula~(2.3.13) for synchrotron radiation is applicable,
%
\begin{equation}
\varepsilon_{\mathrm{synch}}
\simeq \bigl(2.2 \times 10^{-10}~\text{ergs}/(\text{g\,sec})\bigr)\,
  \bar{r}_e^2\, B_\perp^2\,;
\tag{4.4.4}
\end{equation}
%
and, because the gas turns out to be optically thin, the total luminosity is
%
\begin{equation}
L_{\mathrm{synch}}
\simeq \int_{r_g}^{\infty}\varepsilon_{\mathrm{synch}}\,\rho\, 4\pi r^2\, dr
\sim \bigl(1 \times 10^{32}~\text{ergs}/\text{sec}\bigr)
  \left(\frac{M}{M_\odot}\right)^{\!2}
  \left(\frac{\rho_\infty}{10^{-24}~\text{g\,cm}^{-3}}\right)^{\!2}
  \left(\frac{T_\infty}{10^4\,\text{K}}\right)^{\!-3/2}\,.
\tag{4.4.5}
\end{equation}

For comparison, the rest mass-energy being accreted is
%
\begin{equation}
\dot{M}_0 c^2
\sim \bigl(1 \times 10^{32}~\text{ergs}/\text{sec}\bigr)
  \left(\frac{M}{M_\odot}\right)^{\!2}
  \left(\frac{\rho_\infty}{10^{-24}~\text{g\,cm}^{-3}}\right)
  \left(\frac{T_\infty}{10^4\,\text{K}}\right)^{\!-3/2}\,.
\tag{4.4.6}
\end{equation}

Since most of the radiation comes from $r \sim 2r_g$ where the thermal energy is
${\sim}\,3$ to 10 per cent of the rest mass-energy, the synchrotron losses account for
${\sim}\,3$ to 10 per cent of the thermal energy available.  However,
for $M \gtrsim 10\,M_\odot$ the effects of the energy losses on the temperature and thence
on the luminosity will not exceed a factor of ${\sim}\,1$.

Equations~(4.4.5) and (4.4.6) predict that a hole of $M \sim 10\,M_\odot$ will convert
${\sim}\,1$ per cent of the rest mass of the infalling gas into synchrotron radiation,
producing a luminosity of ${\sim}\,10^{32}$\,ergs/sec, which is about 3 per cent of the
luminosity of the sun.  The radiation, located at [cf.\ eq.~(2.3.18)] $r \simeq 2r_g$, will have a spectrum
with a broad maximum located at
%
\begin{equation}
\nu_{\mathrm{peak}}
\sim (7 \times 10^{14}~\text{Hz})
  \left(\frac{\rho_\infty}{10^{-24}~\text{g\,cm}^{-3}}\right)^{\!1/2}
  \left(\frac{T_\infty}{10^4\,\text{K}}\right)^{\!-3/4}
\tag{4.4.7}
\end{equation}
%
($\nu \sim 7 \times 10^{14}$\,Hz, $\lambda \sim 4000$\,\AA); and it will die out exponentially for
$\nu \gtrsim 7 \times 10^{14}$\,Hz; cf.\ \S\ref{sec:2.3}).  As Shvartsman (1971) remarks, such a spectrum
and luminosity resemble those of ``DC white-dwarf stars'' (white dwarfs with
continuous, line-free spectra).  This has led Shvartsman to speculate that
``perhaps some of the objects heretofore regarded as type~DC white dwarfs are
actually black holes.''

Instabilities and inhomogeneities in the flow (due, e.g., to reconnection of
field lines) might lead to marked fluctuations in the luminosity of the hole.  The
timescales for such fluctuations might be of the order of the gas travel time from
${\sim}\,10\,r_g$ to ${\sim}\,2r_g$; i.e.
%
\begin{equation}
\Delta t_{\mathrm{fluctuations}}
\sim (10^{-3}~\text{to}~10^{-4}~\text{sec})(M/M_\odot)\,.
\tag{4.4.8}
\end{equation}

Such rapid fluctuations in an object so faint should be very difficult to detect,
even with the 200-inch telescope.  For further details on the fluctuations, see
the end of \S\ref{sec:4.7}.

%%% ====================================================================
%%% 4.5  Interaction of Outflowing Radiation with the Gas
%%% ====================================================================

\subsection{Interaction of Outflowing Radiation with the Gas}
\label{sec:4.5}

As the luminosity $L$ pours out from the region $r \sim 2r_g$, it interacts with the
inflowing gas.  For the luminosities $L \sim 10^{32}$\,ergs/sec and frequencies $\nu \sim 10^{14}$\,Hz
derived in the last section, one can readily verify that the infalling gas is optically
thin to free-free absorption and to synchrotron reabsorption.  (See \S\ref{sec:2.6} for
basic concepts and equations used in such a calculation.)

What of electron scattering?  Electron scattering can exert 3 types of influence:
(i)~If the electron concentration is sufficiently high, it can act as a ``blanket'' to
impede the outflow of the radiation.  To test for such blanketing one can examine
the optical depth of the gas to electron scattering.
%
\begin{equation}
\tau_{es}(r = 2r_g)
= \int_{2r_g}^{\infty}\kappa_{es}\,\rho\, dr
\simeq (1 \times 10^{-9})(M/M_\odot)
  (\rho_\infty / 10^{-24}~\text{g\,cm}^{-3})
  (T_\infty / 10^4\,\text{K})^{-3/2}\,.
\tag{4.5.1}
\end{equation}

(A more careful calculation would replace the nonrelativistic absorption
coefficient $\kappa_{es} = 0.4$\,cm$^2$/g by the coefficient appropriate to relativistic
electrons, since $T \sim 10^{12}$\,K at $r \sim r_g$; such a calculation would also reveal
extreme optical thinness.)  (ii)~Since the outgoing photons have energies
$h\nu \sim 10$\,eV far less than the electron kinetic energies, Compton scattering can
modify the spectrum, producing high-energy photons [see eq.~(2.4.9)].  However,
the extreme optical thinness guarantees that only about one photon in a million
scatters; so this effect can be ignored.  (iii)~The outpouring photons exert a
pressure on the infalling gas when they scatter, and this effect can retard the
inflow.  Notice that the magnitude of the first two effects is governed by the
density of the inflowing gas.  By contrast, the magnitude of the pressure effect
is governed by the luminosity of the outpouring radiation.  Let us calculate this
effect explicitly, at radii sufficiently large that the electrons are nonrelativistic.

Because the electron scattering cross-section $d\sigma/d\Omega = \tfrac{1}{2}r_0^2(1+\cos^2\theta)$ is
symmetric between forward and backward angles, each scattered photon transfers
its full momentum, $h\nu/c$, to the electron.  Hence, the time-averaged
photon force acting on an electron at radius $r$ is
%
\begin{equation}
\langle\text{Force}\rangle = \left\langle\frac{d(\text{momentum})}{dt}\right\rangle
= \int \frac{d(\text{number of photons})}{d(\text{area})\,dt\,d\nu}\, h\nu\,\sigma_e\, d\nu
\tag{4.5.2}
\end{equation}
%
\[
= \sigma_e F = \sigma_e (L/4\pi r^2)\,,
\]
%
where $\sigma_e = (8\pi/3)r_0^2 = 0.657 \times 10^{-24}$\,cm$^2$ is the total scattering cross section,
$F$ is the radiation flux (ergs\,cm$^{-2}$\,sec$^{-1}$) and $L$ is the total luminosity of the hole.
Notice that this time-averaged force is completely independent of the spectrum
of the radiation (so long as most of the photons have $h\nu \ll m_e c^2$ as seen in the
electron rest frame, so that nonrelativistic cross sections are applicable).  It is
also independent of the optical thickness---being equally valid for $\tau \gg 1$, in
which case $F \ll J$ (see \S2.6 for notation).

The pull of gravity must work against this $\sigma_e (L/4\pi r^2)$ force.  The photon force
acts almost entirely on the electrons ($\sigma_e \sim 1$\,(mass)$^{-2}$, so each proton feels a
force $3 \times 10^6$ times weaker than that felt by an electron).  By contrast, the
acceleration of gravity acts equally on electrons and protons.  Hence, the electrons
move outward slightly relative to the protons, creating an electric field that
transmits the photon force $\sigma_e L/4\pi r^2$ from electrons to protons.  The net force
that then acts on each proton (assuming a completely ionized hydrogen gas) is
%
\begin{equation}
\langle\text{Force}\rangle_{\text{total}} = \frac{\sigma_e L}{4\pi r^2} - \frac{GMm_p}{r^2}\,.
\tag{4.5.3}
\end{equation}

Thus, there is a critical luminosity [``Eddington'' (1926) limit]
%
\begin{equation}
L_{\text{crit}} \equiv \frac{4\pi G M m_p}{\sigma_e}
= (1.3 \times 10^{38}~\text{ergs/sec})\,(M/M_\odot)
\tag{4.5.4}
\end{equation}

such that for $L \ll L_{\text{crit}}$ gravity dominates at all radii, but for $L \gg L_{\text{crit}}$ photon
pressure dominates at all radii.  Eddington (1926) derived this limit for the case of
stellar equilibrium.  Zel'dovich and Novikov (1964) derived this limit for the
cases of quasars and accretion.

For the case of current interest---a black hole of $M\sim 1$ to $100\,M_\odot$ swallowing
interstellar gas---the luminosity is far below the critical value, so the photon
pressure can be ignored.

However, in other contexts the accretion rate $\dot{M}_0$ may be so great, and the
efficiency $\xi$ for converting the inflowing mass-energy into outgoing luminosity
$L$ may be sufficiently high that
%
\begin{equation}
\frac{L}{L_{\text{crit}}} = \frac{\xi \dot{M}_0 c^2}{L_{\text{crit}}}
\lesssim \frac{\xi}{0.1} \left(\frac{\dot{M}_0}{10^{18}~\text{g/sec}}\right)
\tag{4.5.5}
\end{equation}

may temporarily exceed unity.  When this happens, the photon pressure will
impede further mass inflow, thereby reducing the luminosity to $L \lesssim L_{\text{crit}}$.  Such
accretion is said to be ``\emph{self-regulated}''.  Self-regulated accretion, with modifications due to lack of spherical symmetry, may be important for black holes in
close binary systems; see \S5.13.  However, it is unimportant for the case at
hand.

The outflowing radiation from a black hole in interstellar space ($L \sim 10^{32}$
ergs/sec, $\nu \sim 10^{14}$--$10^{15}$\,Hz) can also interact with the surrounding interstellar
gas.  One can verify that the radiation is sufficiently intense to keep the gas at a
temperature $T_\infty \sim 10^4$\,K in the neighborhood of $r = r_i$, and to keep it partially
ionized.

%%% ====================================================================
%%% 4.6  Validity of the Hydrodynamical Approximation
%%% ====================================================================

\subsection{Validity of the Hydrodynamical Approximation}
\label{sec:4.6}

The above model for accretion onto an interstellar black hole relies crucially on
the validity of the hydrodynamical approximation at and outside the sonic point,
and on equipartition assumptions inside the sonic point.  If the gas were to behave
like independent particles rather than like a fluid, then the accretion rate would
be greatly reduced (see \S4.1).  Thus, it is crucial to check the validity of the
hydrodynamical approximation.

In an ionized hydrogen plasma \emph{without} magnetic fields, the deflection of
protons away from straight-line paths is caused primarily by the cumulative
effects of many small-angle Coulomb scatterings.  The total distance a proton
must travel before the ``random-walk'' deflections add up to an angle of ${\sim}90^\circ$
is [see, e.g., Shkarofsky et al.\ (1966) or Pikel'ner (1961)]
%
\begin{equation}
\lambda_p = \frac{9}{4\pi}\,\frac{(kT)^2}{n_p e^4 L_c}
\simeq (7 \times 10^{12}~\text{cm})\left(\frac{\rho_\infty}{10^{-24}~\text{g\,cm}^{-3}}\right)^{-1}
\left(\frac{T}{10^4\,\text{K}}\right)^{2}\,.
\tag{4.6.1}
\end{equation}

Here $L_c$ is a ``Coulomb-logarithm'' (analogue of Gaunt factor in theory of
bremsstrahlung) given by
%
\begin{equation}
L_c \simeq 23 + \tfrac{3}{2}\ln (n_e/\text{cm}^{-3});
\tag{4.6.2}
\end{equation}

$n_e$ and $n_p$ are the number densities of electrons and (ionized) protons; and we
have evaluated $\lambda_p$ for the case of a fully ionized hydrogen plasma.  The plasma
will behave like a fluid on scales $l \gg \lambda_p$, but like independent particles on scales
$l \ll \lambda_p$.

The scale of interest for determining the accretion rate is the sonic radius
%
\begin{equation}
r_s = \frac{5 - 3\Gamma}{4}\,\frac{GM}{a_\infty^2}
\simeq (10^{13}~\text{cm})\left(\frac{M}{M_\odot}\right)
\left(\frac{T_\infty}{10^4\,\text{K}}\right)^{-1}\,,
\tag{4.6.3}
\end{equation}

at which point $T \sim T_\infty$, $\rho \sim \rho_\infty$, so
%
\begin{equation}
\lambda_p(r_s)/r_s \sim 0.7 \left(\frac{M}{M_\odot}\right)^{-1}
\left(\frac{T_\infty}{10^4\,\text{K}}\right)^{-1}\,.
\tag{4.6.4}
\end{equation}

Thus, if one ignores the magnetic field, then the gas will not quite behave like a
fluid near the sonic point---and it may fail even more to be fluid-like at small
radii; cf.\ \S\,13.4 of ZN.

The interstellar magnetic field, of strength $B_\infty \sim 10^{-6}$\,G, can ``save'' the
hydrodynamic approximation.  It deflects protons away from straight-line
motion in a distance of the order of the Larmour radius
%
\begin{equation}
\lambda_L = \frac{m_p c\,v}{eB}
\simeq (1 \times 10^{8}~\text{cm})
\left(\frac{B}{10^{-6}~\text{G}}\right)^{-1}
\left(\frac{T}{10^4\,\text{K}}\right)^{1/2}
\tag{4.6.5}
\end{equation}

and this is true even when the thermal pressure of the plasma exceeds the
magnetic pressure so that the field gets dragged about by the gas.  The Larmour
radius is small compared to all macroscopic scales of interest (e.g.\ small compared
to $r_s$ when one uses the values $B \sim B_\infty$ and $T \sim T_\infty$ appropriate to $r_s$).  Hence,
the hydrodynamical approximation is fairly well justified.

%%% ====================================================================
%%% 4.7  Gas Flow Near the Horizon
%%% ====================================================================

\subsection{Gas Flow Near the Horizon}
\label{sec:4.7}

Near the horizon, $r \sim r_g$, relativistic effects will modify the Newtonian picture
of accretion.  It is in fact difficult to solve for the relativistic flow using the
relativistic Bernoulli equation (2.5.33) and the law of mass conservation (2.5.23$'$).
Such a solution is analogous to the Newtonian hydrodynamical solution studied
in \S4.2.  However, such a solution is partially irrelevant because it ignores the
role of the magnetic field, and such a solution is not needed because the intensity
of the gravitational pull, as viewed in stationary frames near the horizon, is so
great as to guarantee supersonic flow with qualitatively the same form as free-particle
fall.  Thus, from a knowledge of geodesic orbits one can infer the main
features of the flow.

Consider, first, spherical accretion onto a nonrotating hole (Schwarzschild
gravitational field).  A fluid element in the accreting gas, like a freely falling
particle, reaches the horizon in finite proper time.  The moment of crossing the
horizon is in no way peculiar, as seen by proper time.  The density does not
reach infinity there; and the temperature does not reach infinity.  In fact,
because the equation for radial geodesics in the Schwarzschild geometry
%
\begin{equation}
ds^2 = (1 - r_g/r)\,c^2\,dt^2 + (1 - r_g/r)^{-1}\,dr^2 + r^2\,d\Omega^2
\tag{4.7.1}
\end{equation}

has the first integral
%
\begin{equation}
(dr/d\tau)^2 = 2GM/r \quad\text{(for fall from rest at } r \gg r_g\text{)},
\tag{4.7.2}
\end{equation}

and because the law of rest-mass conservation $\boldsymbol{\nabla}\cdot(\rho \mathbf{u}) = 0$ has the first integral
%
\begin{equation}
4\pi r^2 \sqrt{-g}\, \rho\,(dr/d\tau) = \dot{M}_0\,,
\tag{4.7.3}
\end{equation}

the density of rest mass must vary as
%
\begin{equation}
\rho = \frac{\dot{M}_0}{4\pi r^2} \left(\frac{r}{2GM}\right)^{1/2}
\simeq \left(6 \times 10^{-12}~\text{g\,cm}^{-3}\right)
\left(\frac{M}{M_\odot}\right)^{-1}
\left(\frac{T_\infty}{10^4\,\text{K}}\right)^{-3/2}
\,(r/r_g)^{-3/2}\,.
\tag{4.7.4}
\end{equation}

This must be as true near and inside the horizon, $r \lesssim r_g$, as it is in the Newtonian
realm $r \gg r_g$.  Similarly, the temperature will retain its Newtonian form (4.4.2b):
%
\begin{equation}
T \simeq (10^{12}~\text{K})(r_g/r)\,;
\tag{4.7.5}
\end{equation}

and the synchrotron emissivity will retain its Newtonian value (4.4.4).  However,
the radiation which reaches infinity will be sharply cut off as the gas element
passes through the horizon.
%
\begin{equation}
\left(\text{fraction of emitted radiation that}\atop
\text{reaches } r = \infty \text{ from infalling gas at radius } r_f\right)
\simeq \frac{27}{64}\left(1 - \frac{r_g}{r_f}\right)\,.
\tag{4.7.6}
\end{equation}

In previous sections we have taken rough account of this cutoff and of the
accompanying redshift by retaining the Newtonian luminosity down to $r = 2r_g$.

In the case of a rotating (Kerr) hole, near and inside the ergosphere the
dragging of inertial frames will swing the infalling gas into orbital rotation
about the hole.  As the gas approaches the horizon, its angular velocity as seen
from infinity must approach the angular velocity of the horizon,
%
\begin{equation}
\Omega \to \Omega_{\text{horizon}} = \frac{a}{r_+^2 + a^2}
= \frac{c^3}{2GM}
\quad\text{for ``maximally rotating hole'', } a = M
\tag{4.7.7}
\end{equation}
%
(geometrized units, $c = G = 1$).

Here $a$ is the hole's angular momentum per unit mass.  Any extra-luminous hot
spot in the gas (e.g.\ a region where magnetic field lines are reconnecting; cf.\
\S2.9) must orbit the hole with this angular velocity as it falls inward.  The
brightness of the hole as seen at Earth will be modulated with a period
$P \simeq 4\pi r_g/\Omega \gtrsim (10^{-4}~\text{sec})(M/M_\odot)$ as a result of the Doppler shift of the spot's
light and the bending of its rays.  (A factor of $2\pi/\Omega$ is the orbital period; the
other factor of $2\pi/\Omega$ is the light travel time between the radius at which the
spot is located at the beginning of one period and the radius to which it has
fallen at the end of the period.)  Thus, one might expect quasiperiodic
fluctuations in the light from a black hole living alone in the interstellar medium.

%%% ====================================================================
%%% 4.8  Accretion onto a Moving Hole
%%% ====================================================================

\subsection{Accretion onto a Moving Hole}
\label{sec:4.8}

Turn attention now from a hole at rest in the interstellar medium, to a hole that
moves with a speed much larger than the sound speed, $u_\infty \gg a_\infty$. Examine the
resulting accretion in the rest frame of the hole. At radii $r \gg 2GM/u_\infty^2$, the kinetic
energy per unit mass, $\frac{1}{2}u_\infty^2$, is far larger than the potential energy, $GM/r$; so the
pull of the hole has no effect on the gas flow. At $r \sim 2GM/u_\infty^2$, the pull of the
hole becomes significant. Because the flow is supersonic the gas will respond to
that pull in the same manner as would noninteracting particles; its pressure
cannot have an influence. (One can see this, for example, in the Euler equation
%
\begin{equation}
\text{const.} = \tfrac{1}{2}v^2 + \frac{a^2}{\Gamma - 1} - \frac{GM}{r}\,,
\tag{4.8.1}
\end{equation}

where the speed-of-sound (enthalpy) term---which accounts for pressure effects---is
negligible compared to the kinetic-energy term.) Thus, the flow lines will be
identical to the trajectories of test particles: they will be hyperbolae
(Fig.~\ref{fig:4.8.1}a).

After passing around the hole, the trajectories of test particles intersect
(particle collisions; dotted line of Figure~\ref{fig:4.8.1}a). In the gas-dynamic case such
intersection of flow lines is prevented by rapidly mounting pressure. Consequently,
a shock front must develop around the black hole (Fig.~\ref{fig:4.8.1}b).
Outside the shock front the flow lines will be hyperbolic test-particle trajectories.
Behind the shock front, they will not. The shock will be located roughly where
potential energy equals kinetic energy, so its characteristic size $l_s$ as shown in
Figure~\ref{fig:4.8.1}b will be
%
\begin{equation}
l_s \sim GM/u_\infty^2\,.
\tag{4.8.2}
\end{equation}

In the shock, the gas will lose most of its velocity perpendicular to the shock
front [``strong shock'' in terminology of \S\,2.7; ``strong'' because $u_\infty/a_\infty ={}$
(speed of gas)/(speed of sound) $\gg 1$ on front side]; but it will retain all of its
velocity parallel to the front (denote this by $u_\|$). For a given gas element, if the
remaining kinetic energy behind the front greatly exceeds the potential energy,
$\frac{1}{2}u_\|^2 \gg GM/r$, then the gas element will escape the pull of the hole. If
$\frac{1}{2}u_\|^2 \lesssim GM/r$, then it cannot escape. The dividing line between escape and
capture is a dotted line in Figure~\ref{fig:4.8.1}b; and the corresponding impact parameter
is labelled $b_\text{capture}$.

We can calculate $b_\text{capture}$ in order of magnitude by idealizing the shock as
confined to the thin dotted region of Figure~\ref{fig:4.8.1}a (Hoyle and Lyttleton~1939).
The value of $b_\text{capture}$ will correspond to an orbit for which
%
\begin{equation}
\tfrac{1}{2}v_\perp^2 = GM/x \quad \text{at}\quad y = 0.
\tag{4.8.3}
\end{equation}

Straightforward examination of Kepler orbits reveals that
%
\begin{equation}
b_\text{capture} = 2GM/u_\infty^2\,.
\tag{4.8.3$'$}
\end{equation}

Since the mass flux in the gas at large radii is $\rho_\infty u_\infty$, the accretion rate is
%
\begin{equation}
\dot{M}_0 = (\pi b_\text{capture}^2)\,\rho_\infty\, u_\infty
= 4\pi G^2 M^2 \rho_\infty / u_\infty^3\,.
\tag{4.8.4}
\end{equation}

\begin{figure}[htbp]
\centering
% \includegraphics[width=\columnwidth]{fig_4_8_1}
\fbox{\parbox{0.9\columnwidth}{\centering [Figure 4.8.1 placeholder]\\[6pt]
(a)~Trajectories of test particles (and of supersonic gas)
relative to a black hole, as viewed in the rest frame of the hole.\\[4pt]
(b)~The flow lines of supersonic gas
relative to a black hole, as viewed in the rest frame of the hole.
The trajectories and flow lines are virtually the same, until
the gas encounters the shock front.}}
\caption{\textit{Figure~4.8.1.} The trajectories of test particles~(a), and the flow lines
of supersonic gas~(b) relative to a black hole, as viewed in the rest frame of the hole.
The trajectories and flow lines are virtually the same, until the gas encounters the
shock front.}
\label{fig:4.8.1}
\end{figure}

More careful calculations give accretion rates which differ from this by
multiplicative factors of ${\sim}\,0.5$ to $1.0$ (Bondi and Hoyle~1944).

Notice that the accretion rate~(4.8.4) for the supersonic case, and the rate
(4.2.10) for a hole at rest are identical except for the replacement of sound
speed $a_\infty$ by speed of flow $u_\infty$, and except for a numerical coefficient of order
unity. In the intermediate case of $u_\infty \sim a_\infty$, one can use the hybrid formula
%
\begin{equation}
\dot{M}_0 \simeq \frac{4\pi G^2 M^2 \rho_\infty}{(a_\infty^2 + u_\infty^2)^{3/2}}
\tag{4.8.5}
\end{equation}
%
\begin{equation}
\dot{M}_0 \simeq \left(1 \times 10^{11}~\frac{\text{g}}{\text{sec}}\right)
\frac{(M/M_\odot)^2\;(\rho_\infty/10^{-24}~\text{g~cm}^{-3})}
{\bigl[(u_\infty/10~\text{km~sec}^{-1})^2 + (a_\infty/10~\text{km~sec}^{-1})^2\bigr]^{3/2}}\,.
\notag
\end{equation}
%
(Bondi~1952, Hunt~1971).

[Salpeter~(1964) gave the first application of these formulas to black holes; all
previous work dealt with normal stars.]

It is worth noting that a significant fraction of the kinetic energy released in
the shock
%
\begin{equation}
\frac{dE}{dt} \simeq \tfrac{1}{2}\dot{M}_0\, u_\infty^2
\simeq \left(10^{24}~\text{ergs~sec}^{-1}\right)
\left(\frac{M}{M_\odot}\right)^{\!2}
\left(\frac{\rho_\infty}{10^{-24}~\text{g~cm}^{-3}}\right)
\left(\frac{u_\infty}{10~\text{km~sec}^{-1}}\right)^{-1/2}
\tag{4.8.6}
\end{equation}
%
may go into excitation of unionized atoms and thence into light as the atoms
decay, thus producing a shock front that glows (albeit weakly). See, e.g.,
Pikel'ner~(1961), and \S\S\,9 and~10 of Kaplan~(1966).

The motion of the gas behind the shock is practically unknown, particularly
when one takes account of magnetic fields. Many different types of instabilities
may occur, both here and in the case of a stationary black hole. For example at
radii $r < l_s$, where the inflow has become supersonic again, turbulence may
create shock waves which convert kinetic energy into heat (Bisnovatyi-Kogan and
Sunyaev~1971). A variety of plasma instabilities may also arise. And the
reconnection of magnetic field lines will surely be important. One can only guess that
the overall, time-averaged flow will resemble the equipartition model of
Shvartsman~(\S\,4.4).

One of the factors that may be important near the gravitational radius, where
most of the luminosity is produced, is angular momentum. Suppose that the accreting
gas has nonzero angular momentum per unit mass, $\mathscr{L}$, in the direction of motion
of the hole. If the angular momentum per unit mass, $\mathscr{L}$, exceeds
$r_g c$, then centrifugal forces will become important before the infalling gas reaches
the horizon; the gas will be thrown into circulating orbits; and only after viscous
stresses have transported away the excess angular momentum will the gas fall
down the hole. (See \S\,5 for a situation similar to this.) The resulting viscous
heating and lengthened time spent outside the hole, as well as the changed flow
pattern, may affect the outpouring luminosity significantly. Moreover, when the
hole moves from a region with angular momentum of gas in one direction to a
region with angular momentum in another, the gas flow may become particularly
violent for a while near the horizon; and a strong flare of luminosity may be
produced (Salpeter~1964; Shvartsman~1971). Almost nothing is known today
about these issues.

Is the specific angular momentum of the accreting gas likely to exceed $r_g c$?
Interstellar gas is accreted from a region of diameter
%
\begin{equation}
2b_\text{capture} = \frac{2r_g}{u_\infty^2/c^2}
= (3 \times 10^{14}~\text{cm})
\left(\frac{M}{M_\odot}\right)
\left(\frac{u_\infty}{10~\text{km~sec}^{-1}}\right)^{-2}\,.
\tag{4.8.7}
\end{equation}

(For $u_\infty < a_\infty$, we must replace $u_\infty$ by $a_\infty$.) The interstellar gas is turbulent with
the turbulent velocity $v_\text{turb}$ on scales $l$ given roughly by
%
\begin{equation}
v_\text{turb} = \left(10^5~\frac{\text{cm}}{\text{sec}}\right)
\left(\frac{l}{3 \times 10^{20}~\text{cm}}\right)^{q}\,,
\tag{4.8.8}
\end{equation}
%
where $q$ is a coefficient believed to lie between $1/2$ and~$1$ (Kaplan and Pikel'ner
1970), and the coefficient $10^6$\,cm/sec is taken from astronomical observations.
Let us take $q = 0.75$ as a best guess. Then the specific angular momentum
associated with this turbulence, for a gas element of size $l$, is
%
\begin{equation}
\mathscr{L}_\text{turb} \simeq v_\text{turb}\, l
\simeq \left(3 \times 10^{26}~\frac{\text{cm}^2}{\text{sec}}\right)
\left(\frac{l}{3 \times 10^{20}~\text{cm}}\right)^{1.75}\,.
\tag{4.8.9}
\end{equation}

Consequently, the specific angular momentum of the gas that accretes from a
region of size $l = 2b_\text{capture}$ should be
%
\begin{equation}
\frac{\mathscr{L}_\text{turb}}{r_g c}
\simeq 1 \times
\left(\frac{M}{M_\odot}\right)^{3/4}
\left(\frac{u_\infty}{10~\text{km~sec}^{-1}}\right)^{-7/2}\,.
\tag{4.8.10}
\end{equation}

Evidently, the angular momentum may be important in some cases and may not
be in others.

%%% ====================================================================
%%% 4.9  Optical Appearance of Hole: Summary
%%% ====================================================================

\subsection{Optical Appearance of Hole: Summary}
\label{sec:4.9}

In summary, a black hole alone in the interstellar medium may be expected to
emit ${\sim}\,10^{29}$ to $10^{35}$\,ergs/sec of synchrotron radiation, concentrated in the
optical part of the spectrum (${\sim}\,10^{14}$\,Hz to $10^{15}$\,Hz). The output may
fluctuate on timescales of ${\sim}\,10^{-2}$ to $10^{-4}$ seconds due to instabilities, and to
orbital motion of hot spots about the hole (\S\S\,4.4, 4.7, 4.8). There may also
be strong flares when the hole moves out of one turbulent cell of interstellar
gas into another. The time interval between flares would be about
%
\begin{equation}
\Delta t_\text{flares}
\simeq 2b_\text{capture}/u_\infty
\simeq (10~\text{years})
\left(\frac{M}{M_\odot}\right)
\left(\frac{u_\infty}{10~\text{km~sec}^{-1}}\right)^{-3}\,.
\tag{4.9.1}
\end{equation}

It will probably be difficult observationally to distinguish accreting, isolated
black holes from accreting, old neutron stars without magnetic fields. The only
distinguishing feature will be the emission from the surface of the neutron star.


% === SECTION 5: Binary Systems and Galactic Nuclei ===
% Section 5: Black Holes in Binary Star Systems and in the Nuclei of Galaxies
% Placeholder - to be populated by translation agents 17-26
\section{Black Holes in Binary Star Systems and in the Nuclei of Galaxies}
\label{sec:5}

% Content will be merged from agent_17.tex through agent_26.tex


% === SECTION 6: Cosmological Black Holes ===
% =============================================================================
% Section 6: White Holes and Black Holes of Cosmological Origin
% Merged from: agent_26.tex (S6.1), agent_27.tex (S6.2--6.4)
% =============================================================================

\section{White Holes and Black Holes of Cosmological Origin}
\label{sec:6}

%%%%%%%%%%%%%%%%%%%%%%%%%%%%%%%%%%%%%%%%%%%%%%%%%%%%%%%%%%%%%%%%%%%
%% §6.1 White Holes, Grey Holes, and Black Holes
%%%%%%%%%%%%%%%%%%%%%%%%%%%%%%%%%%%%%%%%%%%%%%%%%%%%%%%%%%%%%%%%%%%

\subsection{White Holes, Grey Holes, and Black Holes}
\label{sec:6.1}

All black holes studied in previous sections were regarded as remnants of the
relativistic collapse of massive stars. However, the formation of stars from rarefied
gas is possible only in relatively late stages of the evolution of the Universe
($t \gtrsim 10^7$~years). In earlier stages, according to current cosmological models,
radiation pressure impeded the growth of condensations; so individual bodies
could not form from \emph{small} perturbations of the primeval gas. In the absence of
star formation, there correspondingly should have been no black-hole formation.
However, this conclusion might be wrong. It is quite possible that near the
beginning of the cosmological expansion the matter distribution and metric were
\emph{highly} inhomogeneous and anisotropic. In this case the condensation of individual
bodies would have been possible. Moreover, because in the early stages of the hot
universe the mass-energy in radiation and in particle-antiparticle pairs greatly
exceeded that in baryons, bodies which condensed then must have been made
primarily of photons and pairs. However, because such an ultrarelativistic gas has
an adiabatic index $\Gamma \leq 4/3$, any body made from it is unstable against gravita\-
tional collapse. Thus, one is led to imagine primordial condensations that, like
Wheeler's~(1955) \cite{Wheeler1955} geons, were made primarily from photons and pairs, but unlike
geons, quickly collapsed to form black holes.

Another possible type of strong inhomogeneity in the early universe is a
``white hole'' (Novikov 1964; Ne'eman 1965) \cite{Novikov1964, Neeman1965}. Suppose that some isolated
portions of the universe failed to begin their expansion at the same moment as
the rest of the universe. As measured by external observers, the delay in the
initial expansion of a given region (or ``core'') might be arbitrarily great; and the
delays might vary from core to core. When a ``lagging core'' eventually begins to
expand, and emerges through its gravitational radius, external observers should
see an explosion with a release of tremendous energy. Such an object is, in some
sense, a concrete realization of Ambartsumyan's (1961, 1964) concept of a super\-
dense ``D-body''. (Thus, one can regard ``D-body'', ``white hole'', and ``lagging
core'' as different names for the same concept.)

Finally, if a lagging core, when it begins to expand, has insufficient energy to
emerge through its gravitational radius into the external universe, then one can
regard it as a ``grey hole''. Thus, in the idealized spherical case, the matter of a
grey hole is ``born'' in the Kruskal diagram, staying always to the left of the center line; it moves
upward in the Kruskal diagram, staying always by a past horizon or a future
horizon; and it finally dives to its death in the terminal $r = 0$ singularity. What
would all three types of cosmological holes actually exist in Nature? What
existence be their properties? Theorists have very little observational way to rule out their
about them.

In these notes we shall dwell on only two points: (i) the growth of cosmological
holes by accretion, and (ii) a limit on how many baryons can be inside cosmological
holes.

%% ----------------------------------------------------------------
%% 6.2  The Growth of Cosmological Holes by Accretion
%% ----------------------------------------------------------------

\subsection{The Growth of Cosmological Holes by Accretion}
\label{sec:6.2}

If cosmological holes have existed since near the beginning of the universe,
then during the era when the energy density in radiation exceeded that in matter,
accretion of radiation onto the holes must have been important (Zel'dovich and
Novikov 1966).  For a rough estimate of the effects of such accretion we can use
equation~(4.2.14) for the accretion of gas onto a stationary hole, taking the
velocity of sound to be $a_\infty = c/\sqrt{3}\,c$:
%
\begin{equation}
\tag{6.2.1}
dM/dt \simeq r_g^2 \rho^{\text{rad}}\!.
\end{equation}
%
In standard cosmological models the radiation energy density varies as
$\rho^{\text{rad}} = 1/Gt^2$.

Putting this into the growth equation~(6.2.1), and integrating from the moment
$t_0$ when a hole of initial mass $M_0$ first forms, we obtain for the final mass, at
$t > t_0$:
%
\begin{equation}
\tag{6.2.2}
M = M_0 (1 - G M_0 / c^3 t_0)^{-1}.
\end{equation}

Notice that the final mass diverges for $t_0 \to GM_0/c^3$.  However, the method of
calculation breaks down for this same limit because, for $t_0 \sim GM_0/c^3$ the
characteristic time scale for the accretion is the same as that for the change of
$\rho^{\text{rad}}$, so the ``steady-state'' hypothesis on which equation
(6.2.1) is based.  To determine whether the accretion is catastrophically great at
early times ($t_0 \lesssim GM_0/c^3$), one must solve the accretion problem in a nonsteady,
cosmological context.  One would not be surprised if the final answer depends
sensitively on the particular choice of initial conditions ($t_0 = GM_0/c^3$ or
$t_0 = 0.01\, GM_0/c^3$ or $t_0 = 10^{-10} GM_0/c^3$).  But the problem has not yet been
solved.

%% ----------------------------------------------------------------
%% 6.3  Limit on the Number of Baryons in Cosmological Holes
%% ----------------------------------------------------------------

\subsection{Limit on the Number of Baryons in Cosmological Holes}
\label{sec:6.3}

Independently of the issue of accretion, one can place a tight limit on what
fraction~$\alpha$ of all baryons in the Universe were swallowed into cosmological holes
in the early stages.

Suppose that prior to some particular early moment~$t_*$, a fraction~$\alpha$ of all
baryons had been swallowed into holes.  At the moment~$t_*$, the ratio of rest mass-energy in baryons to mass-energy in radiation and pairs was tiny
%
\begin{equation}
\tag{6.3.1}
\beta \equiv \left[\frac{\rho_0}{\rho^{\text{rad}} + \rho^{\text{(pairs)}}}\right]_{t_*}
\sim 10^{-7} (t_*/1\;\text{sec})^{1/2} \ll 1.
\end{equation}
%
However, as the universe subsequently expanded, this ratio increased until today
it is far greater than unity.  Moreover, the total mass-energy in holes divided by
the volume of the universe ($\alpha \rho_0$) changed in the same manner as
$\rho_0$ ($\propto$~vol$^{-1}$), if no new holes were formed after~$t_*$, and if accretion was negligible.
Otherwise it increased relative to~$\rho_0$.  This means that
%
\begin{equation}
\tag{6.3.2}
\frac{\alpha}{\beta(1-\alpha)}
\leq \left[\frac{\rho^{\text{(holes)}}}{\rho_0}\right]_{t_*}
\leq \left[\frac{\rho^{\text{(holes)}}}{\rho_0}\right]_{\text{today}}
< 80
\end{equation}
%
Here $\alpha/[\beta(1-\alpha)]$ is the value of $[\rho^{\text{(holes)}}/\rho_0]_{t_*}$ if the holes were all created from
primeval plasma at the moment~$t_*$; if they were created even earlier, when rest
mass was even less important, $[\rho^{\text{(holes)}}/\rho_0]_*$ would be even bigger.  The number
80 is a generous observational upper limit on the ratio of nonluminous matter
to luminous matter in the universe today.

Equation~(6.3.2) for $\beta \ll 1/80$ can be rewritten as
%
\begin{equation}
\tag{6.3.3}
\alpha < 1/80\beta \sim 10^{-5} (t_*/1\;\text{sec})^{-1/2}\,.
\end{equation}

This is a very tight limit on the amount of rest mass that could have been down
black holes at early times~$t_*$.

One can also place a tight limit on the amount of mass-energy that white holes
have spewed forth as radiation in recent times (at cosmological redshifts $z \lesssim 10$).
That mass-energy cannot exceed the total mass-energy in radiation today,
%
\begin{equation}
\tag{6.3.4}
\rho_{\text{from white holes}} < \rho_0^{\text{(rad)}}
= 5 \times 10^{-13}\;\text{erg/cm}^3
= 6 \times 10^{-34}\;\text{g/cm}^3\,.
\end{equation}

%% ----------------------------------------------------------------
%% 6.4  Caution
%% ----------------------------------------------------------------

\subsection{Caution}
\label{sec:6.4}

We conclude with a word of caution.  This cosmological section has dealt with
objects about which both theory and observation are equivocal.  There is no firm
theoretical reason for believing that cosmological holes exist, though theory
surely permits them.

By contrast, theory virtually demands the existence of black holes formed
by stellar collapse.  It would be astonishing indeed if every star in the Galaxy
somehow managed to avoid the black-hole fate!  [See Hoyle, Fowler, Burbidge,
and Burbidge (1964).]

The clear discovery of black holes, we expect, is only a few months or years
away.  (The X-ray source Cyg~X-1 is an excellent candidate.)  But cosmological
holes might never be found and might, in fact, not exist at all.  On the other
hand, we might eventually discover that they are the energizers of quasars and
of explosions in galactic nuclei.


% === References ===
\bibliography{NT-Disk-translated}

\end{document}
